% Transformations du plan : définitions analytiques et composées — Mathématiques 1ère C — Cours signé OnBuch+
% Programme officiel MINESEC. Compiler avec : tectonic N13-transformations-plan.tex
\newcommand{\DOCMATIERE}{Mathématiques}
\newcommand{\DOCNIVEAU}{1ère C}
\newcommand{\DOCMODULE}{Configurations et transformations élémentaires du plan}
\newcommand{\DOCLECON}{N13}
\newcommand{\DOCTITRE}{Transformations du plan : définitions analytiques et composées}
% OnBuch+ — préambule « cours signés » ENRICHI (compilé avec tectonic / XeLaTeX)
% Système visuel « Pop » d'OnBuch+ : fond crème (jamais blanc), bordures encre
% épaisses, ombres « dures » décalées, titres Archivo Black en capitales.
% Version MATHÉMATIQUES : graphes (pgfplots/tikz), boîtes pédagogiques
% (définition, propriété, méthode, attention, le savais-tu, exercice résolu,
% exercice, TP, bilan), page de garde et signature OnBuch+ ancrée en bas.
%
% Macros attendues dans main.tex AVANT \input{preamble} :
%   \DOCMATIERE  \DOCNIVEAU  \DOCMODULE  \DOCLECON (numéro)  \DOCTITRE
\documentclass[11pt]{article}

\usepackage{fontspec}
\usepackage[french]{babel}
\usepackage[a4paper,margin=2cm,top=3.4cm,bottom=2.8cm,headheight=1.4cm,footskip=1.3cm]{geometry}
\usepackage{amsmath,amssymb,amsfonts}
\usepackage{mathtools} % \xmapsto, \coloneqq... souvent utilisés par les modèles
\usepackage{cancel} % simplifications barrées
% \abs{z} (module, valeur absolue) et \norm{u} (norme), utilisés par les modèles
\providecommand{\abs}{}\let\abs\relax\DeclarePairedDelimiter{\abs}{\lvert}{\rvert}
\providecommand{\norm}{}\let\norm\relax\DeclarePairedDelimiter{\norm}{\lVert}{\rVert}
\usepackage[table]{xcolor} % [table] : fournit aussi \rowcolors (alternance de lignes)
\usepackage{pagecolor}
\usepackage{tikz}
\usetikzlibrary{arrows.meta,positioning,calc,shapes.geometric,shapes.misc,shapes.arrows,decorations.pathmorphing,decorations.pathreplacing,decorations.markings,patterns,shadows,mindmap,trees,fit,backgrounds,matrix,intersections,angles,quotes}
\usepackage{pgfplots}
\usepackage{pgf-pie} % diagrammes circulaires \pie (statistiques)
\pgfplotsset{compat=1.18}
\usepgfplotslibrary{fillbetween,groupplots} % aires entre courbes (calcul intégral)
\usepackage{enumitem}
\usepackage{fancyhdr}
% [most] charge toutes les bibliothèques tcolorbox (listings, minted, raster,
% documentation...) et gonfle inutilement la mémoire TeX sur un document long
% (~300 boîtes) : on ne charge que ce qui est réellement utilisé.
\usepackage[skins,breakable]{tcolorbox}
\usepackage{graphicx}
\usepackage{colortbl}
\usepackage{booktabs}
\usepackage{tabularx}
\usepackage{diagbox} % case d'en-tête diagonale des tableaux à double entrée
% Colonnes extensibles centrée / gauche / droite (C, L, R), souvent utilisées par les modèles
\newcolumntype{C}{>{\centering\arraybackslash}X}
\newcolumntype{L}{>{\raggedright\arraybackslash}X}
\newcolumntype{R}{>{\raggedleft\arraybackslash}X}
\usepackage{array}
\usepackage{multicol}
\usepackage{siunitx}
% parse-numbers=false : accepte aussi \SI{2,5 \times 10^{-3}}{...}
\sisetup{locale=FR,output-decimal-marker={,},group-minimum-digits=4,parse-numbers=false}
\DeclareSIUnit\ohm{\ensuremath{\symup{\Omega}}} % U+2126 absent de la police
\usepackage{qrcode}
% \degree : commande fréquemment utilisée par les modèles pour un angle ou une
% température en dehors de \SI{}{} (non fournie par les paquets chargés).
\providecommand{\degree}{^{\circ}}
% \qed (fin de démonstration), utilisé par les modèles sans amsthm
\providecommand{\qed}{\hfill$\square$}
% \barre : double barre des valeurs interdites dans un tableau de variations (tkz-tab)
\providecommand{\barre}{\ensuremath{\|}}
% environnement proof (amsthm non chargé) utilisé par les modèles
\ifdefined\proof\else\newenvironment{proof}[1][Démonstration]{\par\noindent\textit{#1.}\ }{\qed\par}\fi
\usepackage{lastpage}
\usepackage{hyperref}
\hypersetup{hidelinks}

% ============================================================
% Palette « Pop » OnBuch+ — mêmes hex que la classe P (plus_ui.dart)
% ============================================================
\definecolor{poporange}{HTML}{F59321}
\definecolor{popdark}{HTML}{E07A0C}
\definecolor{popcream}{HTML}{FFF6E8}
\definecolor{popink}{HTML}{1C1714}
\definecolor{poppurple}{HTML}{7A5AE0}
\definecolor{popgreen}{HTML}{2E9E6A}
\definecolor{poppink}{HTML}{E0457B}
\definecolor{popblue}{HTML}{2F7DE1}
\definecolor{popgold}{HTML}{C98A12}
\definecolor{popmuted}{HTML}{8A7F76}
\definecolor{popline}{HTML}{EADFD2}
\definecolor{popgray}{HTML}{8A7F76}
\colorlet{popgrey}{popgray}
% Alias tolérants (le modèle nomme parfois les couleurs différemment)
\colorlet{popwhite}{white}
\colorlet{popblack}{popink}
\colorlet{popred}{poppink}
\colorlet{popyellow}{popgold}
% Teintes claires pour fonds de tableaux / zones
\colorlet{popblueL}{popblue!10!white}
\colorlet{poporangeL}{poporange!14!white}
\colorlet{popgreenL}{popgreen!12!white}
\colorlet{poppurpleL}{poppurple!12!white}
\colorlet{poppinkL}{poppink!12!white}
\colorlet{popgoldL}{popgold!14!white}

\pagecolor{popcream}

% ============================================================
% Typographie
% ============================================================
\defaultfontfeatures{Ligatures=TeX}
\setmainfont{PlusJakartaSans-Medium.ttf}[Path=../../../fonts/,BoldFont=PlusJakartaSans-Bold.ttf]
% Maths en sans-serif (Fira Math) avec chiffres et lettres droites de Plus
% Jakarta, pour que les formules se fondent dans le texte.
\usepackage[math-style=ISO,bold-style=ISO]{unicode-math}
\setmathfont{FiraMath-Regular.otf}[Path=../../../fonts/]
\setmathfont{PlusJakartaSans-Medium.ttf}[Path=../../../fonts/,range={up/num,up/{Latin,latin},\mathup}]
\usepackage{icomma} % virgule décimale sans espace en mode math
% Caractères mathématiques Unicode absents des polices de texte (Plus Jakarta,
% Archivo Black) : rendus par leur équivalent mathématique, en texte comme en
% formule — sinon ils s'affichent en carré vide (ex. « Arithmétique dans ℤ »).
\usepackage{newunicodechar}
\newunicodechar{ℝ}{\ensuremath{\mathbb{R}}}
\newunicodechar{ℤ}{\ensuremath{\mathbb{Z}}}
\newunicodechar{ℂ}{\ensuremath{\mathbb{C}}}
\newunicodechar{ℕ}{\ensuremath{\mathbb{N}}}
\newunicodechar{ℚ}{\ensuremath{\mathbb{Q}}}
\newunicodechar{𝔻}{\ensuremath{\mathbb{D}}}
\newunicodechar{α}{\ensuremath{\alpha}}
\newunicodechar{β}{\ensuremath{\beta}}
\newunicodechar{γ}{\ensuremath{\gamma}}
\newunicodechar{δ}{\ensuremath{\delta}}
\newunicodechar{ε}{\ensuremath{\varepsilon}}
\newunicodechar{θ}{\ensuremath{\theta}}
\newunicodechar{λ}{\ensuremath{\lambda}}
\newunicodechar{μ}{\ensuremath{\mu}}
\newunicodechar{σ}{\ensuremath{\sigma}}
\newunicodechar{τ}{\ensuremath{\tau}}
\newunicodechar{φ}{\ensuremath{\varphi}}
\newunicodechar{ψ}{\ensuremath{\psi}}
\newunicodechar{ω}{\ensuremath{\omega}}
\newunicodechar{Σ}{\ensuremath{\Sigma}}
% majuscules produites par \MakeUppercase dans les titres (α-aminés -> Α)
\newunicodechar{Α}{\ensuremath{\alpha}}
\newunicodechar{Β}{\ensuremath{\beta}}
\newunicodechar{Γ}{\ensuremath{\Gamma}}
\newunicodechar{Δ}{\ensuremath{\Delta}}
\newunicodechar{Ε}{\ensuremath{\varepsilon}}
\newunicodechar{Θ}{\ensuremath{\Theta}}
\newunicodechar{Λ}{\ensuremath{\Lambda}}
\newunicodechar{Μ}{\ensuremath{\mu}}
\newunicodechar{Τ}{\ensuremath{\tau}}
\newunicodechar{Φ}{\ensuremath{\Phi}}
\newunicodechar{Ψ}{\ensuremath{\Psi}}
\newunicodechar{Ω}{\ensuremath{\Omega}}
\newunicodechar{↦}{\ensuremath{\mapsto}}
\newunicodechar{∈}{\ensuremath{\in}}
\newunicodechar{∉}{\ensuremath{\notin}}
\newunicodechar{∧}{\ensuremath{\wedge}}
\newunicodechar{⇔}{\ensuremath{\Leftrightarrow}}
\newunicodechar{⟺}{\ensuremath{\Longleftrightarrow}}
\newunicodechar{⇒}{\ensuremath{\Rightarrow}}
\newunicodechar{⟹}{\ensuremath{\Longrightarrow}}
\newunicodechar{∘}{\ensuremath{\circ}}
\newunicodechar{∪}{\ensuremath{\cup}}
\newunicodechar{∩}{\ensuremath{\cap}}
\newunicodechar{≡}{\ensuremath{\equiv}}
\newunicodechar{⊂}{\ensuremath{\subset}}
\newunicodechar{⊥}{\ensuremath{\perp}}
\newunicodechar{∃}{\ensuremath{\exists}}
\newunicodechar{∀}{\ensuremath{\forall}}
\newunicodechar{∛}{\ensuremath{\sqrt[3]{\,}}}
\newunicodechar{∜}{\ensuremath{\sqrt[4]{\,}}}
\newunicodechar{⃗}{\ensuremath{{}^{\to}}}
\newunicodechar{ⁿ}{\ifmmode^{n}\else\textsuperscript{\ensuremath{n}}\fi}
\newunicodechar{ˣ}{\ifmmode^{x}\else\textsuperscript{\ensuremath{x}}\fi}
\newunicodechar{ᵇ}{\ifmmode^{b}\else\textsuperscript{\ensuremath{b}}\fi}
\newunicodechar{ʳ}{\ifmmode^{r}\else\textsuperscript{\ensuremath{r}}\fi}
\newunicodechar{ᵘ}{\ifmmode^{u}\else\textsuperscript{\ensuremath{u}}\fi}
\newunicodechar{ʸ}{\ifmmode^{y}\else\textsuperscript{\ensuremath{y}}\fi}
\newunicodechar{ᵃ}{\ifmmode^{a}\else\textsuperscript{\ensuremath{a}}\fi}
\newunicodechar{ᵉ}{\ifmmode^{e}\else\textsuperscript{\ensuremath{e}}\fi}
\newunicodechar{ᵏ}{\ifmmode^{k}\else\textsuperscript{\ensuremath{k}}\fi}
\newunicodechar{ᵖ}{\ifmmode^{p}\else\textsuperscript{\ensuremath{p}}\fi}
\newunicodechar{ᴺ}{\ifmmode^{N}\else\textsuperscript{\ensuremath{N}}\fi}
\newunicodechar{ᶜ}{\ifmmode^{c}\else\textsuperscript{\ensuremath{c}}\fi}
\newunicodechar{ʰ}{\ifmmode^{h}\else\textsuperscript{\ensuremath{h}}\fi}
\newunicodechar{ᵠ}{\ifmmode^{q}\else\textsuperscript{\ensuremath{q}}\fi}
\newunicodechar{ₙ}{\ifmmode_{n}\else\textsubscript{\ensuremath{n}}\fi}
\newunicodechar{ₐ}{\ifmmode_{a}\else\textsubscript{\ensuremath{a}}\fi}
\newunicodechar{ᵢ}{\ifmmode_{i}\else\textsubscript{\ensuremath{i}}\fi}
\newunicodechar{ᵣ}{\ifmmode_{r}\else\textsubscript{\ensuremath{r}}\fi}
\newunicodechar{ₖ}{\ifmmode_{k}\else\textsubscript{\ensuremath{k}}\fi}
\newunicodechar{ᵦ}{\ifmmode_{\beta}\else\textsubscript{\ensuremath{\beta}}\fi}
\newunicodechar{ₓ}{\ifmmode_{x}\else\textsubscript{\ensuremath{x}}\fi}
\newfontfamily\popdisplay{ArchivoBlack-Regular.ttf}[Path=../../../fonts/]
\newfontfamily\poplogo{SpaceGrotesk-Bold.ttf}[Path=../../../fonts/]
\newfontfamily\popsemi{PlusJakartaSans-SemiBold.ttf}[Path=../../../fonts/]
\newfontfamily\popxbold{PlusJakartaSans-ExtraBold.ttf}[Path=../../../fonts/]
\newfontfamily\popxboldls{PlusJakartaSans-ExtraBold.ttf}[Path=../../../fonts/,LetterSpace=3.0]

\setlist[enumerate]{leftmargin=1.7em,itemsep=5pt,topsep=4pt}
\setlist[itemize]{leftmargin=1.7em,itemsep=4pt,topsep=4pt,label={\color{poporange}\textbullet}}
\setlength{\parindent}{0pt}
\setlength{\parskip}{5pt}
\renewcommand{\arraystretch}{1.25}

% pgfplots : style Pop par défaut
\pgfplotsset{
  every axis/.append style={axis line style={popink,line width=1pt},tick style={popink},
    grid style={popline,line width=0.5pt},label style={font=\small},tick label style={font=\footnotesize}},
  popcurve/.style={poporange,line width=1.8pt,no markers,smooth},
}
% Même style disponible dans un \draw TikZ simple (hors pgfplots)
\tikzset{popcurve/.style={poporange,line width=1.8pt}}

% ============================================================
% Pill / squiggle
% ============================================================
\newcommand{\poppill}[3]{%
  \tikz[baseline=-0.55ex]\node[draw=popink,line width=1.2pt,rounded corners=2.6pt,%
    fill=#2,inner xsep=6.5pt,inner ysep=3pt]{%
    \popxboldls\fontsize{7}{7}\selectfont\color{#3}\MakeUppercase{#1}};%
}
\newcommand{\popsquiggle}[1]{%
  \tikz[baseline=-0.4ex]{
    \draw[#1,line width=2.6pt,line cap=round]
      plot [smooth] coordinates {(0,0.09) (0.34,0.01) (0.68,0.21) (1.02,0.01) (1.36,0.10)};
  }%
}
% Mot-clé surligné (marqueur orange sous le texte)
\newcommand{\cle}[1]{{\popxbold\color{popdark}#1}}

% ============================================================
% En-tête / pied
% ============================================================
\pagestyle{fancy}
\fancyhf{}
\renewcommand{\headrulewidth}{0pt}
\renewcommand{\footrulewidth}{0pt}
\lhead{\raisebox{-0.15ex}{\poplogo\fontsize{13}{13}\selectfont\color{popink}OnBuch\color{poporange}+}}
\rhead{\poppill{\DOCMATIERE\ \textbullet\ \DOCNIVEAU\ \textbullet\ Leçon \DOCLECON}{white}{popink}}
\lfoot{\popxbold\fontsize{7.6}{7.6}\selectfont\color{popmuted}\MakeUppercase{Cours signé OnBuch+}\ \textbullet\ {\popsemi\DOCTITRE}}
\rfoot{\tikz[baseline=-0.6ex]\node[rounded corners=8.5pt,fill=popink,minimum height=17pt,minimum width=17pt,inner xsep=5pt,inner sep=0]{\popxbold\fontsize{8}{8}\selectfont\color{popcream}\thepage\,/\,\pageref*{LastPage}};}

% ============================================================
% Titres
% ============================================================
\newcommand{\chaptitre}[2]{%
  \begin{tcolorbox}[enhanced,colback=poporange,colframe=popink,boxrule=2.6pt,arc=11pt,
      drop shadow=popink,left=15pt,right=15pt,top=12pt,bottom=11pt,before skip=3pt,after skip=18pt]
    \poppill{#2}{popink}{popcream}\par\vspace{9pt}
    {\raggedright\hyphenpenalty=10000\exhyphenpenalty=10000\popdisplay\fontsize{22}{24}\selectfont\color{popcream}\MakeUppercase{#1}\par}
  \end{tcolorbox}
}
% Section numérotée : pastille orange avec le numéro + titre Archivo
\newcounter{popsec}
\newcommand{\coursec}[1]{%
  \stepcounter{popsec}\setcounter{popsub}{0}%
  \par\vspace{16pt}\noindent
  \begin{minipage}{\linewidth}
  \tikz[baseline=-0.7ex]\node[draw=popink,line width=1.6pt,fill=poporange,rounded corners=4pt,minimum size=22pt,inner sep=0]{\popdisplay\fontsize{12}{12}\selectfont\color{popcream}\thepopsec};%
  \hspace{8pt}{\popdisplay\fontsize{15}{17}\selectfont\color{popink}\MakeUppercase{#1}}\par
  \vspace{4pt}\noindent{\color{poporange}\rule{36pt}{3.6pt}}
  \end{minipage}\par\nopagebreak\vspace{8pt}
}
\newcounter{popsub}
\newcommand{\courssub}[1]{%
  \stepcounter{popsub}%
  \par\vspace{9pt}\noindent{\popxbold\fontsize{11.5}{13}\selectfont\color{popink}{\color{poporange}\thepopsec.\thepopsub}\hspace{6pt}#1}\par\nopagebreak\vspace{4pt}
}

% ============================================================
% Boîtes pédagogiques — même recette : fond blanc, bordure épaisse colorée,
% ombre dure encre, badge plein. Titre optionnel [#1].
% ============================================================
\tcbset{popbase/.style={breakable,enhanced,colback=white,boxrule=2.3pt,arc=9pt,
  shadow={3pt}{-3pt}{0pt}{popink},left=14pt,right=14pt,top=11pt,bottom=11pt,before skip=11pt,after skip=15pt,
  fontupper=\popsemi\fontsize{10.6}{14.5}\selectfont\color{popink},
  fontlower=\popsemi\fontsize{10.6}{14.5}\selectfont\color{popink}}}
% Coche dessinée (le glyphe ✓ n'existe pas dans les polices du document)
\newcommand{\popcheck}[1]{\tikz[baseline=-0.5ex]\draw[#1,line width=1.6pt,line cap=round,line join=round](-0.13,0.0)--(-0.04,-0.09)--(0.13,0.1);}
\newcommand{\popboxhead}[3]{\poppill{#1}{#2}{white}\ifx\relax#3\relax\else\hspace{6pt}{\popxbold\fontsize{10.6}{12}\selectfont\color{popink}#3}\fi\par\vspace{7pt}}

\newtcolorbox{aretenir}[1][]{popbase,colframe=popgreen,before upper={\popboxhead{À retenir}{popgreen}{#1}}}
\newtcolorbox{exemplebox}[1][]{popbase,colframe=poporange,before upper={\popboxhead{Exemple}{poporange}{#1}}}
\newtcolorbox{definition}[1][]{popbase,colframe=popblue,before upper={\popboxhead{Définition}{popblue}{#1}}}
\newtcolorbox{propriete}[1][]{popbase,colframe=poppurple,before upper={\popboxhead{Propriété}{poppurple}{#1}}}
\newtcolorbox{methode}[1][]{popbase,colframe=popgold,before upper={\popboxhead{Méthode}{popgold}{#1}}}
\newtcolorbox{attention}[1][]{popbase,colframe=poppink,before upper={\popboxhead{Attention}{poppink}{#1}}}
\newtcolorbox{experience}[1][]{popbase,colframe=popblue,colback=popblueL,before upper={\popboxhead{Expérience}{popblue}{#1}}}
% « Le savais-tu ? » : carte violette pleine (contexte, culture, Cameroun)
\newtcolorbox{savaistu}[1][]{popbase,colback=poppurple,colframe=popink,
  fontupper=\popsemi\fontsize{10.4}{14.2}\selectfont\color{white},
  before upper={\poppill{Le savais-tu\,?}{white}{poppurple}\ifx\relax#1\relax\else\hspace{6pt}{\popxbold\color{white}#1}\fi\par\vspace{7pt}}}
% Exercice résolu : énoncé en haut, \tcblower puis la solution sur fond crème
\newcounter{popexo}\newcounter{popexr}
\newtcolorbox{exoresolu}[1][]{popbase,colframe=poporange,colbacklower=popcream,
  segmentation style={popink,line width=1.2pt,dashed},
  before upper={\stepcounter{popexr}\popboxhead{Exercice résolu \thepopexr}{poporange}{#1}},
  before lower={\poppill{Solution}{popink}{popcream}\par\vspace{6pt}}}
% Exercice (non résolu en place) — la correction est dans la partie Corrigés
\newtcolorbox{exercice}[1][]{popbase,colframe=popink,
  before upper={\stepcounter{popexo}\popboxhead{Exercice \thepopexo}{popink}{#1}}}
% Corrigé
\newtcolorbox{corrige}[1][]{popbase,colframe=popgreen,colback=popgreenL,
  before upper={\popboxhead{Corrigé}{popgreen}{#1}}}
% Objectifs / prérequis / bilan
\newtcolorbox{objectifs}{popbase,colframe=popink,colback=poporangeL,
  before upper={\popboxhead{Objectifs de la leçon}{popink}{}}}
\newtcolorbox{prerequis}{popbase,colframe=popink,colback=popblueL,
  before upper={\popboxhead{Prérequis}{popblue}{}}}
\newtcolorbox{bilan}{popbase,colframe=popink,colback=white,boxrule=2.8pt,
  before upper={\popboxhead{Fiche bilan}{poporange}{L'essentiel à connaître pour l'examen}}}
% Figure encadrée avec légende
\newtcolorbox{popfigure}[1][]{enhanced,breakable=false,colback=white,colframe=popline,boxrule=1.4pt,arc=8pt,
  left=10pt,right=10pt,top=10pt,bottom=8pt,before skip=10pt,after skip=12pt,halign=center,
  lower separated=false,halign lower=center,
  fontlower=\popsemi\fontsize{9}{11.5}\selectfont\color{popmuted},
  #1}
\newcommand{\legende}[1]{\par\vspace{6pt}{\popsemi\fontsize{9}{11.5}\selectfont\color{popmuted}#1}}

% ============================================================
% Signature OnBuch+
% ============================================================
\newcommand{\poplogoinline}[1]{{\poplogo\fontsize{#1}{#1}\selectfont\color{popink}OnBuch\color{poporange}+}}
\newcommand{\signatureonbuch}{%
  \begin{tcolorbox}[enhanced,colback=popink,colframe=popink,boxrule=2.2pt,arc=10pt,
      drop shadow=poporange,left=14pt,right=14pt,top=10pt,bottom=10pt,before skip=0pt,after skip=0pt]
    \tikz[baseline=-0.7ex]\node[circle,fill=popgreen,draw=popcream,line width=1.3pt,minimum size=22pt,inner sep=0]{\popcheck{white}};%
    \hspace{9pt}\begin{minipage}[c]{0.62\linewidth}
      {\popxbold\fontsize{12}{13}\selectfont\color{popcream}Signé\ \textbullet\ {\poplogo OnBuch\color{poporange}+}}\par\vspace{2pt}
      {\raggedright\popsemi\fontsize{8}{10.5}\selectfont\color{popline}\MakeUppercase{Équipe pédagogique OnBuch+}\\\MakeUppercase{Conforme au programme officiel MINESEC}\par}
    \end{minipage}\hfill
    {\poplogo\fontsize{22}{22}\selectfont\color{popcream}OnBuch\color{poporange}+}
  \end{tcolorbox}%
}

% ============================================================
% Page de garde
% #1 titre · #2 matière · #3 niveau · #4 module · #5 n° leçon · #6 durée ·
% #7 description / objectifs (texte ou itemize)
% ============================================================
\newcommand{\pagegarde}[7]{%
  \thispagestyle{empty}
  \newgeometry{a4paper,margin=1.7cm,top=1.6cm,bottom=1.4cm}%
  \begin{tikzpicture}[remember picture,overlay]
    \fill[poporange!18] ([xshift=-3.2cm,yshift=-3.6cm]current page.north east) circle (4.6cm);
    \fill[poppurple!14] ([xshift=2.4cm,yshift=7.6cm]current page.south west) circle (3.8cm);
  \end{tikzpicture}%
  {\poplogo\fontsize{30}{30}\selectfont\color{popink}OnBuch\color{poporange}+}\hfill
  \raisebox{0.6ex}{\poppill{Cours signé \textbullet\ Premium}{popink}{popcream}}\par
  \vspace{1pt}\popsquiggle{poppurple}\par
  \vspace{8pt}
  \begin{tcolorbox}[enhanced,colback=poporange,colframe=popink,boxrule=2.8pt,arc=13pt,
      drop shadow=popink,left=18pt,right=18pt,top=12pt,bottom=13pt,after skip=10pt]
    \poppill{#2 \textbullet\ #3}{popink}{popcream}\hspace{5pt}\poppill{#4}{white}{popink}\par\vspace{8pt}
    {\popdisplay\fontsize{14}{14}\selectfont\color{popink}LEÇON #5}\par\vspace{4pt}
    {\raggedright\hyphenpenalty=10000\exhyphenpenalty=10000\popdisplay\fontsize{25}{27}\selectfont\color{popcream}\MakeUppercase{#1}\par}
    \vspace{8pt}
    {\popxbold\fontsize{9.5}{9.5}\selectfont\color{popink}\MakeUppercase{Durée conseillée : #6}}
  \end{tcolorbox}
  \begin{tcolorbox}[enhanced,colback=white,colframe=popink,boxrule=2.2pt,arc=10pt,
      drop shadow=popink,left=16pt,right=16pt,top=10pt,bottom=10pt,after skip=8pt]
    \poppill{Dans cette leçon}{poppurple}{white}\par\vspace{5pt}
    \popsemi\fontsize{9.3}{12.6}\selectfont\color{popink}#7
  \end{tcolorbox}
  \noindent\begin{minipage}[t]{0.487\linewidth}
    \begin{tcolorbox}[enhanced,colback=popgreen,colframe=popink,boxrule=2.2pt,arc=10pt,
        drop shadow=popink,left=11pt,right=11pt,top=10pt,bottom=10pt,height=3.1cm,valign=center]
      \tcbox[colback=white,colframe=popink,boxrule=1.3pt,arc=5pt,boxsep=0pt,nobeforeafter,
        left=3pt,right=3pt,top=3pt,bottom=3pt,valign=center]{\qrcode[height=1.9cm]{https://play.google.com/store/apps/details?id=com.luvvix.onbuch&hl=fr}}%
      \hspace{8pt}\begin{minipage}[c]{0.5\linewidth}
        {\popdisplay\fontsize{11}{12.5}\selectfont\color{white}\MakeUppercase{Télécharge OnBuch}}\par\vspace{4pt}
        {\raggedright\popsemi\fontsize{8.4}{11}\selectfont\color{white}Gratuit \textbullet\ Google Play\\Tuteur IA, annales, concours, résultats\par}
      \end{minipage}
    \end{tcolorbox}
  \end{minipage}\hfill
  \begin{minipage}[t]{0.487\linewidth}
    \begin{tcolorbox}[enhanced,colback=poppurple,colframe=popink,boxrule=2.2pt,arc=10pt,
        drop shadow=popink,left=11pt,right=11pt,top=10pt,bottom=10pt,height=3.1cm,valign=center]
      \tikz[baseline=-0.7ex]\node[circle,fill=white,draw=popink,line width=1.3pt,minimum size=1.6cm,inner sep=0]{\popdisplay\fontsize{24}{24}\selectfont\color{poppurple}+};%
      \hspace{8pt}\begin{minipage}[c]{0.6\linewidth}
        {\popdisplay\fontsize{11}{12.5}\selectfont\color{white}\MakeUppercase{Passe à OnBuch+}}\par\vspace{4pt}
        {\raggedright\popsemi\fontsize{8.4}{11}\selectfont\color{white}Tous les cours signés, Tuteur IA illimité, fiches PDF — onglet OnBuch+ de l'app\par}
      \end{minipage}
    \end{tcolorbox}
  \end{minipage}
  \vspace{8pt}
  \popcontenu
  \vfill
  \signatureonbuch
  \restoregeometry
  \newpage
}

% Rangée « Ce cours contient » (page de garde)
\newcommand{\poptile}[3]{% #1 couleur #2 gros texte #3 légende
  \begin{tcolorbox}[enhanced,colback=#1,colframe=popink,boxrule=2pt,arc=9pt,shadow={2.5pt}{-2.5pt}{0pt}{popink},
      left=6pt,right=6pt,top=9pt,bottom=9pt,halign=center,valign=center,height=1.75cm,width=\linewidth,nobeforeafter]
    {\popdisplay\fontsize{12.5}{13}\selectfont\color{white}\MakeUppercase{#2}}\par\vspace{3pt}
    {\centering\popsemi\fontsize{7.8}{9.5}\selectfont\color{white}#3\par}
  \end{tcolorbox}}
\newcommand{\popcontenu}{%
  \par\noindent{\popxboldls\fontsize{7.5}{7.5}\selectfont\color{popmuted}CE COURS SIGNÉ CONTIENT}\par\vspace{7pt}
  \noindent\begin{minipage}[t]{0.235\linewidth}\poptile{popblue}{Cours}{complet, illustré, pas à pas}\end{minipage}\hfill
  \begin{minipage}[t]{0.235\linewidth}\poptile{poporange}{Méthodes}{et exercices résolus}\end{minipage}\hfill
  \begin{minipage}[t]{0.235\linewidth}\poptile{poppink}{Exercices}{type examen + corrigés}\end{minipage}\hfill
  \begin{minipage}[t]{0.235\linewidth}\poptile{popgreen}{Fiche bilan}{l'essentiel à retenir}\end{minipage}\par}

% Sommaire
\newtcolorbox{sommaire}{enhanced,colback=white,colframe=popink,boxrule=2.2pt,arc=10pt,
  drop shadow=popink,left=18pt,right=18pt,top=13pt,bottom=13pt,before skip=6pt,after skip=20pt,
  before upper={\poppill{Sommaire}{poppurple}{white}\par\vspace{9pt}\popsemi\fontsize{10.6}{14.5}\selectfont\color{popink}}}

% Signature de fin de cours — poussée tout en bas de la dernière page
\newcommand{\findecours}{%
  \par\vfill\nopagebreak
  \begin{center}\popsquiggle{poporange}\end{center}\vspace{4pt}
  \signatureonbuch
}
\AtBeginDocument{\def\llbracket{\mathopen{[\mkern-3.5mu[}}\def\rrbracket{\mathclose{]\mkern-3.5mu]}}} % intervalles d'entiers (glyphe ⟦ absent de la police)
% Glyphes absents de Fira Math : redéfinis à partir de glyphes disponibles
\AtBeginDocument{%
  \def\setminus{\mathbin{\backslash}}\let\smallsetminus\setminus
  \def\vdots{\mathord{\vbox{\baselineskip3.2pt\lineskiplimit0pt\kern4pt\hbox{.}\hbox{.}\hbox{.}}}}%
  \def\ddots{\mathinner{\mkern1mu\raise6.5pt\hbox{.}\mkern2mu\raise3.6pt\hbox{.}\mkern2mu\raise0.7pt\hbox{.}\mkern1mu}}%
  \def\triangle{\mathord{\scalebox{0.9}{$\symup{\Delta}$}}}%
  \def\nabla{\mathord{\raisebox{\depth}{\scalebox{1}[-1]{$\symup{\Delta}$}}}}%
  \def\checkmark{\popcheck{popdark}}%
  \def\star{\mathbin{\ast}}\def\bigstar{\mathord{\ast}}%
  \def\diamond{\mathbin{\scalebox{0.6}{$\Diamond$}}}%
  \def\bigcup{\mathop{\mathchoice{\raisebox{-0.3em}{\scalebox{1.7}{$\cup$}}}{\scalebox{1.25}{$\cup$}}{\cup}{\cup}}\displaylimits}%
  \def\bigcap{\mathop{\mathchoice{\raisebox{-0.3em}{\scalebox{1.7}{$\cap$}}}{\scalebox{1.25}{$\cap$}}{\cap}{\cap}}\displaylimits}%
  \def\lesseqqgtr{\mathrel{\lessgtr}}\def\bigoplus{\mathop{\oplus}\displaylimits}%
}
\providecommand{\ssi}{si et seulement si} % abréviation fréquente
\AtBeginDocument{\providecommand{\wideparen}{\overparen}} % arc de cercle

%% FIGFIX-BEGIN v5 — correctifs globaux des figures (docs/onbuch-generation/tools/figfix)
\usepackage{adjustbox}\usepackage{etoolbox}\tcbuselibrary{raster}
% toute figure TikZ/pgfplots plus large que la ligne est réduite à la largeur disponible
\BeforeBeginEnvironment{tikzpicture}{\begin{adjustbox}{max width=\linewidth}}
\AfterEndEnvironment{tikzpicture}{\end{adjustbox}}
% légendes pgfplots lisibles par-dessus les courbes
\pgfplotsset{every axis/.append style={legend style={fill=white,fill opacity=0.92,text opacity=1,draw=black!15,inner sep=3pt}}}
% étiquettes d'arêtes (syntaxe edge["..."]) avec fond blanc
\tikzset{every edge quotes/.append style={fill=white,inner sep=1.2pt}}
%% FIGFIX-END
\begin{document}

\pagegarde{Transformations du plan : définitions analytiques et composées}{Mathématiques}{1ère C}{Configurations et transformations élémentaires du plan}{N13}{14 h}{Cette leçon établit les définitions analytiques rigoureuses des translations et homothéties dans un repère, puis explore la structure algébrique des composées de transformations élémentaires (rotations, homothéties, translations) pour résoudre des problèmes géométriques complexes.
\vspace{4pt}
\begin{multicols}{2}\small
\begin{itemize}
\item À la fin de cette leçon, je sais définir analytiquement une translation et une homothétie dans un repère orthonormé.
\item Je sais reconnaître une translation ou une homothétie à partir de son écriture analytique (équations ou matrice).
\item Je sais démontrer qu'une application du plan dans lui-même est une transformation (bijection) et déterminer analytiquement sa réciproque.
\item Je sais caractériser la composée de deux rotations de même centre (nature, angle, centre).
\item Je sais déterminer la nature et les éléments caractéristiques de la composée de deux homothéties de même centre.
\item Je sais déterminer la nature et les éléments caractéristiques de la composée d'une homothétie et d'une translation (homothétie ou translation).
\end{itemize}\end{multicols}}

\begin{sommaire}
\begin{multicols}{2}
\begin{enumerate}
\item 1. Rappels et cadre général : Transformations du plan
\item 2. Définition analytique de la translation
\item 3. Définition analytique de l'homothétie
\item 4. Composée de deux rotations de même centre
\item 5. Composée de deux homothéties de même centre
\item 6. Composée d'une homothétie et d'une translation
\item 7. Composée d'une homothétie et d'une rotation de même centre (Similitude directe)
\item Composées de symétries orthogonales, de rotations et d'homothéties de centres distincts
\item Activité d'intégration
\item Méthodes et exercices résolus
\item Exercices
\item Corrigés des exercices
\item Fiche bilan
\end{enumerate}
\end{multicols}
\end{sommaire}

\begin{objectifs}
\begin{itemize}
\item À la fin de cette leçon, je sais définir analytiquement une translation et une homothétie dans un repère orthonormé.
\item Je sais reconnaître une translation ou une homothétie à partir de son écriture analytique (équations ou matrice).
\item Je sais démontrer qu'une application du plan dans lui-même est une transformation (bijection) et déterminer analytiquement sa réciproque.
\item Je sais caractériser la composée de deux rotations de même centre (nature, angle, centre).
\item Je sais déterminer la nature et les éléments caractéristiques de la composée de deux homothéties de même centre.
\item Je sais déterminer la nature et les éléments caractéristiques de la composée d'une homothétie et d'une translation (homothétie ou translation).
\item Je sais déterminer la nature et les éléments caractéristiques de la composée d'une homothétie et d'une rotation de même centre (similitude directe).
\item J'utilise les propriétés des transformations composées pour démontrer des propriétés géométriques sur des configurations planes (alignements, parallélismes, rapports de longueurs).
\end{itemize}
\end{objectifs}

\begin{prerequis}
\begin{itemize}
\item Notion de repère orthonormé et coordonnées d'un vecteur (Seconde).
\item Définition géométrique et propriétés de la translation, rotation, homothétie, symétrie (Seconde).
\item Calcul vectoriel : somme, produit par un réel, colinéarité, coordonnées du milieu (Seconde/Première).
\item Résolution d'équations linéaires et systèmes 2x2 (Seconde).
\item Notion d'application, d'image, d'antécédent, d'injection, de surjection, de bijection (Début Première).
\end{itemize}
\end{prerequis}

\newpage

\coursec{Rappels et cadre général : Transformations du plan}

Au carrefour du Wouri à Douala, un agent de circulation fait pivoter sa main de $90^\circ$ puis de $45^\circ$ pour guider les véhicules : quelle rotation unique équivaut à ces deux mouvements successifs ? Sur le chantier du barrage de Nachtigal, un plan de masse est photographié par un drone qui zoome (homothétie) puis se déplace latéralement (translation) : comment retrouver la position réelle d'un pylône sur la carte ? Au marché Mokolo, un artisan reproduit un motif de pagne en le glissant (translation) puis en l'agrandissant (homothétie) : l'ordre des opérations change-t-il le résultat final ?

Toutes ces questions ont un point commun : elles font agir, l'une après l'autre, des \cle{transformations du plan}. Avant de composer les transformations (sections 4 à 7 de cette leçon) et de les décrire par des formules (sections 2 et 3), il faut d'abord répondre à une question plus fondamentale : \emph{qu'est-ce qu'une transformation du plan, exactement ?} C'est l'objet de cette première section, qui pose le cadre général de toute la leçon.

\courssub{Qu'est-ce qu'une transformation du plan ?}

\textbf{Intuition.} Imagine le plan comme une immense feuille de papier calque posée sur la carte de la ville de Yaoundé. \og Transformer \fg{} le plan, c'est déplacer \emph{chaque} point de la feuille vers un nouveau point, selon une règle précise et uniforme : on peut glisser la feuille (translation), la faire pivoter autour d'une punaise (rotation), l'agrandir ou la rétrécir depuis un point fixe (homothétie). Mais il y a une exigence essentielle : après le déplacement, on doit pouvoir \emph{revenir en arrière} de façon unique. Aucun point ne doit être \og écrasé \fg{} sur un autre, et aucune zone ne doit rester sans provenance. Cette exigence s'appelle la \cle{bijectivité}.

\begin{definition}[Transformation du plan]
On appelle \cle{transformation du plan} $\mathcal{P}$ toute application \emph{bijective} de $\mathcal{P}$ dans lui-même, c'est-à-dire toute application $f : \mathcal{P} \to \mathcal{P}$ vérifiant les deux conditions :
\begin{itemize}
\item $f$ est \cle{injective} : deux points distincts ont des images distinctes ;
\item $f$ est \cle{surjective} : tout point du plan est l'image d'au moins un point.
\end{itemize}
Autrement dit : pour tout point $M'$ du plan, il existe un \emph{unique} point $M$ tel que $f(M) = M'$.
\end{definition}

\begin{attention}[Application ou transformation ?]
Erreur fréquente : confondre \og application du plan dans lui-même \fg{} et \og transformation \fg{}. Toute transformation est une application, mais \textbf{l'inverse est faux}. Par exemple, la projection orthogonale sur une droite $(D)$ envoie le plan \emph{dans} le plan, mais elle écrase toute une droite perpendiculaire sur un seul point : elle n'est pas injective. Ce n'est \emph{pas} une transformation. De même, l'application constante qui envoie tout point sur un point fixe $O$ n'est ni injective ni surjective.
\end{attention}

\begin{exemplebox}[Contre-exemples fondamentaux]
\begin{itemize}
\item \textbf{Projection orthogonale sur l'axe des abscisses} : $p(x ; y) = (x ; 0)$. Les points $A(1 ; 2)$ et $B(1 ; 5)$ ont la même image $(1 ; 0)$ : $p$ n'est pas injective. De plus, le point $(0 ; 3)$ n'a aucun antécédent : $p$ n'est pas surjective.
\item \textbf{Application constante} : $c(M) = O$ pour tout $M$. Tout le plan s'écrase sur un seul point.
\end{itemize}
Ces applications sont pourtant bien définies sur tout le plan : ce sont des applications, pas des transformations.
\end{exemplebox}

\begin{popfigure}
\begin{tikzpicture}[scale=1.0]
% --- Gauche : non injective ---
\begin{scope}[xshift=-4.2cm]
\draw[popline, dashed] (-1.4,-1.4) rectangle (1.4,1.4);
\draw[popline, dashed] (3.0,-1.4) rectangle (5.8,1.4);
\node[popdark, font=\small\bfseries] at (0,1.8) {Non injective};
\fill[popblue] (-0.7,0.6) circle (2pt) node[left, popdark]{$A$};
\fill[popblue] (0.7,0.6) circle (2pt) node[right, popdark]{$B$};
\fill[poporange] (4.4,0.4) circle (2.5pt) node[right, popdark]{$M'$};
\draw[-{Stealth}, thick, popblue] (-0.55,0.6) .. controls (1.6,1.3) .. (4.2,0.55);
\draw[-{Stealth}, thick, popblue] (0.85,0.6) .. controls (2.6,1.3) .. (4.5,0.6);
\node[popmuted, font=\small] at (2.2,-1.9) {Deux points, une même image};
\end{scope}
% --- Droite : non surjective ---
\begin{scope}[xshift=4.6cm]
\draw[popline, dashed] (-1.4,-1.4) rectangle (1.4,1.4);
\draw[popline, dashed] (3.0,-1.4) rectangle (5.8,1.4);
\node[popdark, font=\small\bfseries] at (2.2,1.8) {Non surjective};
\fill[popblue] (0,0.4) circle (2pt) node[left, popdark]{$C$};
\fill[poporange] (4.0,0.3) circle (2.5pt) node[right, popdark]{$N'$};
\fill[popgreen] (4.6,-0.8) circle (2.5pt) node[right, popdark]{$P'$};
\draw[-{Stealth}, thick, popblue] (0.2,0.4) .. controls (2.0,1.0) .. (3.85,0.4);
\draw[-{Stealth}, thick, popgreen, dashed] (4.6,-0.8) -- (4.6,-1.35);
\node[popmuted, font=\small] at (2.2,-1.9) {$P'$ n'a aucun antécédent};
\end{scope}
\end{tikzpicture}
\legende{Deux défauts qui empêchent une application d'être une transformation : à gauche, deux points ont la même image (non injective) ; à droite, un point n'a pas d'antécédent (non surjective).}
\end{popfigure}

\courssub{Réciproque d'une transformation}

Puisqu'une transformation $f$ est bijective, chaque point image $M'$ provient d'un unique point $M$. On peut donc \og défaire \fg{} la transformation : c'est la \cle{réciproque}.

\begin{definition}[Transformation réciproque]
Soit $f$ une transformation du plan. Il existe une unique transformation, notée $f^{-1}$ et appelée \cle{transformation réciproque} de $f$, telle que :
\[ f \circ f^{-1} = f^{-1} \circ f = \mathrm{Id}_{\mathcal{P}} \]
où $\mathrm{Id}_{\mathcal{P}}$ est l'identité du plan ($\mathrm{Id}_{\mathcal{P}}(M) = M$ pour tout point $M$). Concrètement : $f(M) = M' \iff f^{-1}(M') = M$.
\end{definition}

\begin{propriete}[Réciproques des transformations usuelles]
\begin{itemize}
\item La réciproque de la translation de vecteur $\vec{u}$ est la translation de vecteur $-\vec{u}$.
\item La réciproque de la rotation de centre $\Omega$ et d'angle $\theta$ est la rotation de centre $\Omega$ et d'angle $-\theta$.
\item La réciproque de l'homothétie de centre $\Omega$ et de rapport $k \neq 0$ est l'homothétie de centre $\Omega$ et de rapport $\dfrac{1}{k}$.
\item La réciproque d'une symétrie (axiale ou centrale) est elle-même : on dit qu'elle est \cle{involutive}.
\end{itemize}
Ces transformations conservent les propriétés géométriques fondamentales : alignement, milieux, contacts ; les isométries conservent en outre les distances et les angles.
\end{propriete}

\courssub{Représentation analytique d'une application du plan}

Le plan est muni d'un repère $(O ; \vec{i}, \vec{j})$. Décrire une application $f$ du plan dans lui-même revient à donner les coordonnées $(x' ; y')$ de l'image $M' = f(M)$ en fonction des coordonnées $(x ; y)$ de $M$ :
\[ \begin{cases} x' = f_1(x, y) \\ y' = f_2(x, y) \end{cases} \]
C'est la \cle{définition analytique} (ou représentation analytique) de $f$. Dans cette leçon, nous nous intéressons surtout aux \cle{applications affines}, celles dont les formules sont du premier degré en $x$ et $y$ :
\[ \begin{cases} x' = ax + by + p \\ y' = cx + dy + q \end{cases} \qquad\text{que l'on écrit sous forme matricielle :}\qquad \overrightarrow{OM'} = A\,\overrightarrow{OM} + B, \]
avec $A = \begin{pmatrix} a & b \\ c & d \end{pmatrix}$ la \cle{matrice linéaire} (ou partie linéaire) et $B = \begin{pmatrix} p \\ q \end{pmatrix}$ la partie de translation.

\begin{exemplebox}[Lecture d'une définition analytique]
L'application $f : (x ; y) \mapsto (2x + y + 1 \,;\, x - y - 3)$ est affine, de matrice linéaire $A = \begin{pmatrix} 2 & 1 \\ 1 & -1 \end{pmatrix}$ et de partie de translation $B = \begin{pmatrix} 1 \\ -3 \end{pmatrix}$. L'image du point $M(1 ; 2)$ est $M'(2\times 1 + 2 + 1 \,;\, 1 - 2 - 3) = (5 ; -4)$.
\end{exemplebox}

\courssub{Critère analytique de bijectivité}

Voici le résultat central de cette section : il permet de décider, par un simple calcul, si une application affine est une transformation.

\begin{propriete}[Critère du déterminant]
Soit $f$ l'application affine du plan définie par $\overrightarrow{OM'} = A\,\overrightarrow{OM} + B$, avec $A = \begin{pmatrix} a & b \\ c & d \end{pmatrix}$. Alors :
\[ f \text{ est une transformation du plan} \iff \det(A) = ad - bc \neq 0. \]
Dans ce cas, la réciproque $f^{-1}$ est aussi affine, donnée par $\overrightarrow{OM} = A^{-1}\left(\overrightarrow{OM'} - B\right)$, où $A^{-1} = \dfrac{1}{\det(A)}\begin{pmatrix} d & -b \\ -c & a \end{pmatrix}$.
\end{propriete}

\begin{experience}[Pourquoi le déterminant décide-t-il de tout ?]
Cherchons si, $M'(x' ; y')$ étant donné, le système $\begin{cases} ax + by = x' - p \\ cx + dy = y' - q \end{cases}$ admet une solution unique $(x ; y)$.
\begin{enumerate}
\item Multiplions la première ligne par $d$, la seconde par $b$, et soustrayons : $(ad - bc)x = d(x' - p) - b(y' - q)$.
\item De même, en éliminant $x$ : $(ad - bc)y = a(y' - q) - c(x' - p)$.
\item Si $\det(A) = ad - bc \neq 0$, on peut diviser : on obtient \emph{une unique} solution $(x ; y)$ pour chaque $(x' ; y')$. L'application est donc bijective : c'est une transformation.
\item Si $\det(A) = 0$, alors soit le système n'a aucune solution ($M'$ sans antécédent : non surjectif), soit il en a une infinité (plusieurs points ont la même image : non injectif). Dans les deux cas, $f$ n'est pas une transformation.
\end{enumerate}
\end{experience}

\begin{methode}[Vérifier qu'une application est une transformation]
\begin{enumerate}
\item Écrire l'application sous forme matricielle $\overrightarrow{OM'} = A\,\overrightarrow{OM} + B$ : identifier $A$ et $B$.
\item Calculer le déterminant de la matrice linéaire : $\det(A) = ad - bc$.
\item Si $\det(A) \neq 0$ : c'est une bijection, donc une transformation. Exprimer la réciproque en calculant $A^{-1}$ puis en résolvant $\overrightarrow{OM} = A^{-1}\left(\overrightarrow{OM'} - B\right)$.
\item Si $\det(A) = 0$ : ce n'est pas une transformation (exhiber un contre-exemple le confirme).
\end{enumerate}
\end{methode}

\begin{exoresolu}[Une application affine est-elle une transformation ?]
\textbf{Énoncé.} Le plan est muni d'un repère $(O ; \vec{i}, \vec{j})$. On considère l'application $f$ qui à tout point $M(x ; y)$ associe le point $M'(x' ; y')$ tel que
\[ \begin{cases} x' = 2x + y \\ y' = x - y \end{cases} \]
Montrer que $f$ est une transformation du plan et déterminer l'expression analytique de sa réciproque $f^{-1}$.
\tcblower
\textbf{Solution.} Suivons la méthode.

\textbf{Étape 1 : forme matricielle.} On a $\overrightarrow{OM'} = A\,\overrightarrow{OM}$ avec
\[ A = \begin{pmatrix} 2 & 1 \\ 1 & -1 \end{pmatrix}, \qquad B = \begin{pmatrix} 0 \\ 0 \end{pmatrix}. \]

\textbf{Étape 2 : déterminant.} $\det(A) = 2 \times (-1) - 1 \times 1 = -2 - 1 = -3 \neq 0$.

\textbf{Étape 3 : conclusion et réciproque.} Comme $\det(A) = -3 \neq 0$, $f$ est une bijection du plan : c'est une \cle{transformation}. Calculons la matrice inverse :
\[ A^{-1} = \frac{1}{-3}\begin{pmatrix} -1 & -1 \\ -1 & 2 \end{pmatrix} = \begin{pmatrix} \dfrac{1}{3} & \dfrac{1}{3} \\[4pt] \dfrac{1}{3} & -\dfrac{2}{3} \end{pmatrix}. \]
Comme $B$ est nul, $\overrightarrow{OM} = A^{-1}\overrightarrow{OM'}$, c'est-à-dire :
\[ \begin{cases} x = \dfrac{x' + y'}{3} \\[4pt] y = \dfrac{x' - 2y'}{3} \end{cases} \]
Donc $f^{-1} : (x ; y) \mapsto \left(\dfrac{x + y}{3} \,;\, \dfrac{x - 2y}{3}\right)$.

\textbf{Vérification.} Prenons $M(1 ; 1)$ : $f(M) = (3 ; 0)$. Appliquons $f^{-1}$ à $(3 ; 0)$ : on obtient $\left(\frac{3 + 0}{3} ; \frac{3 - 0}{3}\right) = (1 ; 1)$. On retombe bien sur $M$.
\end{exoresolu}

\begin{attention}[Pièges de rédaction]
\begin{itemize}
\item Ne conclus jamais \og c'est une transformation \fg{} sans avoir calculé le déterminant (ou prouvé l'existence et l'unicité de l'antécédent).
\item Le déterminant se calcule sur la \textbf{matrice linéaire $A$ uniquement} : les constantes $p$ et $q$ (partie de translation) n'interviennent pas dans le critère.
\item Dans $A^{-1}$, n'oublie pas le facteur $\dfrac{1}{\det(A)}$ devant la matrice $\begin{pmatrix} d & -b \\ -c & a \end{pmatrix}$, et attention aux changements de signes.
\end{itemize}
\end{attention}

\courssub{Composition des transformations : une structure de groupe}

Effectuer deux transformations l'une après l'autre, c'est les \cle{composer} : si $f$ puis $g$ agissent, la transformation résultante est $g \circ f$ (attention à l'ordre de lecture : $f$ agit d'abord). C'est exactement la situation de l'agent de circulation du Wouri : pivoter de $90^\circ$ puis de $45^\circ$, c'est composer deux rotations.

\begin{propriete}[Le groupe des transformations du plan]
L'ensemble des transformations du plan, muni de la composition $\circ$, vérifie :
\begin{itemize}
\item \textbf{Stabilité} : la composée de deux transformations est une transformation (la composée de deux bijections est une bijection) ;
\item \textbf{Associativité} : pour toutes transformations $f, g, h$ : $(f \circ g) \circ h = f \circ (g \circ h)$ ;
\item \textbf{Élément neutre} : l'identité $\mathrm{Id}_{\mathcal{P}}$ vérifie $f \circ \mathrm{Id}_{\mathcal{P}} = \mathrm{Id}_{\mathcal{P}} \circ f = f$ ;
\item \textbf{Symétrisation} : toute transformation $f$ admet une réciproque $f^{-1}$ telle que $f \circ f^{-1} = f^{-1} \circ f = \mathrm{Id}_{\mathcal{P}}$.
\end{itemize}
On dit que $(\mathcal{T}, \circ)$ est un \cle{groupe}, appelé groupe des transformations du plan. De plus, $(g \circ f)^{-1} = f^{-1} \circ g^{-1}$ : \emph{l'ordre s'inverse dans la réciproque d'une composée}.
\end{propriete}

\begin{attention}[La composition n'est pas commutative en général]
En général, $g \circ f \neq f \circ g$ : l'ordre des opérations compte ! C'est précisément la question de l'artisan du marché Mokolo : glisser puis agrandir ne donne pas le même résultat qu'agrandir puis glisser. Nous étudierons cela en détail dans la section 6. (Cas particulier rassurant : deux rotations de \emph{même centre} commutent entre elles.)
\end{attention}

\courssub{Classification des transformations affines}

Les transformations que nous étudierons dans cette leçon sont toutes affines. Elles se rangent en grandes familles selon les propriétés géométriques qu'elles conservent.

\begin{popfigure}
\begin{tikzpicture}[
  grow=down,
  level 1/.style={sibling distance=7.2cm, level distance=1.7cm},
  level 2/.style={sibling distance=2.6cm, level distance=1.7cm},
  every node/.style={draw, rounded corners=3pt, align=center, font=\small, thick},
  edge from parent/.style={draw=popmuted, thick, -{Stealth}}
]
\node[fill=popblueL, draw=popblue, align=center] {Applications du plan\\ dans lui-même}
  child { node[fill=poporangeL, draw=poporange, align=center] {Applications affines\\ $\overrightarrow{OM'} = A\overrightarrow{OM} + B$}
    child { node[fill=popgreenL, draw=popgreen, align=center] {Transformations affines\\ $\det(A) \neq 0$}
      child { node[fill=poppurpleL, draw=poppurple, font=\footnotesize, align=center] {Isométries\\ (translations, rotations,\\ symétries)} }
      child { node[fill=poppinkL, draw=poppink, font=\footnotesize, align=center] {Similitudes\\ (homothéties,\\ similitudes directes)} }
      child { node[fill=popgoldL, draw=popgold, font=\footnotesize, align=center] {Affinités\\ (autres cas)} }
    }
  };
\end{tikzpicture}
\legende{Arbre de classification : parmi toutes les applications du plan, les applications affines ; parmi elles, celles dont le déterminant est non nul sont les transformations affines, qui se répartissent en isométries, similitudes et affinités.}
\end{popfigure}

\begin{aretenir}[L'essentiel de la section]
\begin{itemize}
\item Une \cle{transformation du plan} est une application \emph{bijective} du plan dans lui-même : injective et surjective.
\item Toute transformation $f$ admet une unique réciproque $f^{-1}$ avec $f \circ f^{-1} = \mathrm{Id}_{\mathcal{P}}$.
\item Une application affine $\overrightarrow{OM'} = A\overrightarrow{OM} + B$ est une transformation \textbf{si et seulement si} $\det(A) \neq 0$ ; sa réciproque s'obtient avec $A^{-1}$.
\item Les transformations forment un groupe pour la composition : stabilité, associativité, neutre $\mathrm{Id}_{\mathcal{P}}$, réciproque ; mais la composition n'est pas commutative en général.
\end{itemize}
\end{aretenir}

\begin{savaistu}[Des transformations partout autour de nous]
Les logiciels de cartographie utilisés par les services cadastraux du Cameroun représentent chaque parcelle de terrain par des coordonnées et appliquent des transformations affines pour passer d'une photographie aérienne à la carte officielle : le critère $\det(A) \neq 0$ garantit qu'aucune parcelle ne \og disparaît \fg{} ni ne se superpose à une autre. En infographie, les jeux vidéo et les logiciels de dessin n'ont cessé d'utiliser, depuis les travaux du mathématicien allemand Félix Klein (programme d'Erlangen, 1872), l'idée que la géométrie est l'étude des propriétés conservées par des groupes de transformations : c'est exactement la démarche de cette leçon.
\end{savaistu}

\begin{exercice}[Pour s'entraîner]
Dans un repère $(O ; \vec{i}, \vec{j})$, on considère les applications :
\[ f : (x ; y) \mapsto (3x - 2y + 1 \,;\, x + y), \qquad g : (x ; y) \mapsto (2x - y \,;\, -4x + 2y + 5). \]
\begin{enumerate}
\item Montrer que $f$ est une transformation du plan et déterminer l'expression analytique de $f^{-1}$.
\item Montrer que $g$ n'est pas une transformation. Exhiber deux points distincts ayant la même image par $g$.
\end{enumerate}
\end{exercice}

\begin{corrige}[Exercice 1]
\textbf{1.} La matrice linéaire de $f$ est $A = \begin{pmatrix} 3 & -2 \\ 1 & 1 \end{pmatrix}$, et $\det(A) = 3 \times 1 - (-2) \times 1 = 5 \neq 0$ : $f$ est une transformation. On a
\[ A^{-1} = \frac{1}{5}\begin{pmatrix} 1 & 2 \\ -1 & 3 \end{pmatrix}. \]
Le système $\begin{cases} x' = 3x - 2y + 1 \\ y' = x + y \end{cases}$ équivaut à $\begin{cases} x' - 1 = 3x - 2y \\ y' = x + y \end{cases}$, d'où
\[ \begin{cases} x = \dfrac{(x' - 1) + 2y'}{5} = \dfrac{x' + 2y' - 1}{5} \\[4pt] y = \dfrac{-(x' - 1) + 3y'}{5} = \dfrac{-x' + 3y' + 1}{5} \end{cases} \]
Donc $f^{-1} : (x ; y) \mapsto \left(\dfrac{x + 2y - 1}{5} \,;\, \dfrac{-x + 3y + 1}{5}\right)$.

\textbf{2.} La matrice linéaire de $g$ est $A = \begin{pmatrix} 2 & -1 \\ -4 & 2 \end{pmatrix}$, et $\det(A) = 2 \times 2 - (-1)(-4) = 4 - 4 = 0$ : $g$ n'est \emph{pas} une transformation. Contre-exemple : $M_1(0 ; 0)$ a pour image $(0 ; 5)$, et $M_2(1 ; 2)$ a pour image $(2 - 2 \,;\, -4 + 4 + 5) = (0 ; 5)$. Deux points distincts ont la même image : $g$ n'est pas injective.
\end{corrige}

\coursec{Définition analytique de la translation}

Dans la section précédente, nous avons posé le cadre général : une \cle{transformation du plan} est une bijection du plan dans lui-même, et toute transformation admet une réciproque qui est elle-même une transformation. Nous avons aussi rappelé que les translations, homothéties, rotations et symétries sont les transformations élémentaires au programme. Il est temps maintenant de passer du langage géométrique au langage des \cle{coordonnées} : c'est toute la puissance de la géométrie analytique. Placer un repère, c'est donner un « code » à chaque point ; décrire une transformation analytiquement, c'est donner la \emph{recette de calcul} qui transforme le code de $M$ en celui de son image $M'$. Commençons par la plus simple de toutes : la translation.

\courssub{Rappel géométrique : la translation de vecteur $\vec{u}$}

Imagine un car de l'agence « Garanti Express » qui quitte Douala pour Yaoundé. Chaque passager, chaque siège, chaque valise effectue \emph{exactement le même déplacement} : même direction, même sens, même distance. C'est cela, une translation : un glissement \emph{global} du plan, sans rotation, sans déformation, sans retournement.

\begin{definition}[Translation de vecteur $\vec{u}$ (rappel)]
Soit $\vec{u}$ un vecteur du plan. On appelle \cle{translation de vecteur $\vec{u}$}, notée $t_{\vec{u}}$, l'application du plan dans lui-même qui, à tout point $M$, associe le point $M'$ tel que :
\[ \overrightarrow{MM'} = \vec{u}. \]
On note $M' = t_{\vec{u}}(M)$. Le vecteur $\vec{u}$ est l'\cle{élément caractéristique} de la translation.
\end{definition}

\begin{popfigure}
\begin{tikzpicture}[scale=1.05,>=Stealth]
% Repère
\draw[->,popmuted] (-0.8,0) -- (7.2,0) node[below left] {$x$};
\draw[->,popmuted] (0,-0.8) -- (0,3.6) node[below left] {$y$};
\draw[popmuted] (0,0) node[below left] {$O$};
% Points A et B
\coordinate (A) at (1,0.8);
\coordinate (B) at (3,1.6);
% Vecteur u = (3,1)
\coordinate (Ap) at ($(A)+(3,1)$);
\coordinate (Bp) at ($(B)+(3,1)$);
% Figures
\fill[popblue!15,draw=popblue,thick] (A) circle (0.05) node[below left,popdark] {$A(1\,;\,0{,}8)$};
\fill[popblue!15,draw=popblue,thick] (B) circle (0.05) node[below,popdark] {$B(3\,;\,1{,}6)$};
\fill[poporange!20,draw=poporange,thick] (Ap) circle (0.05) node[above right,popdark] {$A'(4\,;\,1{,}8)$};
\fill[poporange!20,draw=poporange,thick] (Bp) circle (0.05) node[above,popdark] {$B'(6\,;\,2{,}6)$};
% Vecteurs
\draw[->,very thick,popgreen] (A) -- (Ap) node[midway,sloped,above] {$\vec{u}$};
\draw[->,very thick,popgreen] (B) -- (Bp) node[midway,sloped,above] {$\vec{u}$};
\draw[->,very thick,popgreen] (0.2,3.1) -- (3.2,3.6) node[midway,sloped,above] {$\vec{u}(3\,;\,1)$};
% Segments AA' et BB' pointillés vers A'B'
\draw[dashed,popmuted] (A) -- (B);
\draw[dashed,popmuted] (Ap) -- (Bp);
\end{tikzpicture}
\legende{La translation $t_{\vec{u}}$ avec $\vec{u}(3\,;\,1)$ : $\overrightarrow{AA'} = \overrightarrow{BB'} = \vec{u}$. Tous les points « glissent » du même vecteur.}
\end{popfigure}

\begin{attention}[Le vecteur nul]
Si $\vec{u} = \vec{0}$, alors $\overrightarrow{MM'} = \vec{0}$, c'est-à-dire $M' = M$ pour tout point $M$ : on obtient l'\cle{identité} du plan. L'identité est donc une translation \emph{particulière} (la translation de vecteur nul). C'est un point que les élèves oublient souvent : lorsqu'on démontre qu'une composée de transformations est une translation, obtenir l'identité est une conclusion tout à fait acceptable — c'est la translation de vecteur nul.
\end{attention}

\courssub{Écriture analytique dans un repère $(O\,;\vec{\imath},\vec{\jmath})$}

Fixons maintenant un repère $(O\,;\vec{\imath},\vec{\jmath})$ du plan. Tout point $M$ a des coordonnées $(x\,;\,y)$, ce qui signifie $\overrightarrow{OM} = x\vec{\imath} + y\vec{\jmath}$. Soit $\vec{u}$ un vecteur de coordonnées $(x_u\,;\,y_u)$, et soit $M'(x'\,;\,y')$ l'image de $M(x\,;\,y)$ par $t_{\vec{u}}$.

Par définition, $\overrightarrow{MM'} = \vec{u}$. Or les coordonnées de $\overrightarrow{MM'}$ sont $(x' - x\,;\,y' - y)$. L'égalité vectorielle se traduit donc, coordonnée par coordonnée, par :
\[ x' - x = x_u \qquad \text{et} \qquad y' - y = y_u. \]

\begin{definition}[Écriture analytique de la translation]
Dans un repère $(O\,;\vec{\imath},\vec{\jmath})$, la translation de vecteur $\vec{u}(x_u\,;\,y_u)$ est l'application :
\[ t_{\vec{u}} : \begin{cases} x' = x + x_u \\ y' = y + y_u \end{cases} \]
Autrement dit : $t_{\vec{u}} : (x\,;\,y) \longmapsto (x + x_u\,;\,y + y_u)$.
\end{definition}

Cette écriture est d'une simplicité remarquable : \textbf{traduire un point, c'est ajouter les coordonnées du vecteur à ses coordonnées}. On dit que la translation est une application \cle{affine} dont la partie linéaire est l'identité.

\begin{exemplebox}[Image d'un point par une translation]
Dans un repère orthonormé, soit $\vec{u}(-2\,;\,5)$ et $A(4\,;\,1)$. Déterminons les coordonnées de $A' = t_{\vec{u}}(A)$.
\[ x_{A'} = x_A + x_u = 4 + (-2) = 2 \qquad ; \qquad y_{A'} = y_A + y_u = 1 + 5 = 6. \]
Donc $A'(2\,;\,6)$. Vérification géométrique : $\overrightarrow{AA'}(2-4\,;\,6-1) = (-2\,;\,5) = \vec{u}$. \checkmark
\end{exemplebox}

\courssub{Forme matricielle : $X' = X + U$}

L'écriture analytique se condense élégamment en langage matriciel. Associons à chaque point $M(x\,;\,y)$ sa \cle{matrice colonne} $X = \begin{pmatrix} x \\ y \end{pmatrix}$, et au vecteur $\vec{u}$ la colonne $U = \begin{pmatrix} x_u \\ y_u \end{pmatrix}$. Le système
\[ \begin{cases} x' = x + x_u \\ y' = y + y_u \end{cases} \quad \text{s'écrit} \quad \begin{pmatrix} x' \\ y' \end{pmatrix} = \begin{pmatrix} x \\ y \end{pmatrix} + \begin{pmatrix} x_u \\ y_u \end{pmatrix}, \quad \text{c'est-à-dire} \quad \boxed{X' = X + U}. \]

Regardons de plus près la structure de cette application. Une application de la forme $X' = AX + B$, où $A$ est une matrice $2\times 2$ (la \cle{partie linéaire}) et $B$ une colonne (le \cle{terme constant}), est appelée application affine. Pour la translation :
\[ \begin{pmatrix} x' \\ y' \end{pmatrix} = \underbrace{\begin{pmatrix} 1 & 0 \\ 0 & 1 \end{pmatrix}}_{I_2} \begin{pmatrix} x \\ y \end{pmatrix} + \begin{pmatrix} x_u \\ y_u \end{pmatrix}. \]

\begin{propriete}[Caractérisation matricielle d'une translation]
Une application $f$ du plan dans lui-même, d'écriture $X' = AX + B$ dans un repère, est une \cle{translation} si et seulement si sa partie linéaire est la matrice identité : $A = I_2 = \begin{pmatrix} 1 & 0 \\ 0 & 1 \end{pmatrix}$.
\begin{itemize}
\item Si $B \neq \begin{pmatrix} 0 \\ 0 \end{pmatrix}$, c'est la translation de vecteur $B$.
\item Si $B = \begin{pmatrix} 0 \\ 0 \end{pmatrix}$, c'est l'identité (translation de vecteur nul).
\end{itemize}
On remarque que $\det(I_2) = 1$ : la translation \emph{conserve les aires et l'orientation}.
\end{propriete}

\begin{experience}[Démonstration guidée : pourquoi la partie linéaire est-elle $I_2$ ?]
Soit $f$ une application affine d'expression $f : (x\,;\,y) \mapsto (ax + by + c\,;\,dx + ey + g)$. Nous allons montrer que si $f$ est une translation, alors $a = e = 1$ et $b = d = 0$.
\begin{enumerate}
\item \textbf{Traduisons l'hypothèse.} Dire que $f$ est une translation de vecteur $\vec{u}(x_u\,;\,y_u)$ signifie : pour tout $M(x\,;\,y)$, $\overrightarrow{Mf(M)} = \vec{u}$, c'est-à-dire que le vecteur $\overrightarrow{Mf(M)}$ est \emph{constant} (ne dépend pas de $M$).
\item \textbf{Calculons $\overrightarrow{Mf(M)}$.} Ses coordonnées sont :
\[ \big( ax + by + c - x \,;\, dx + ey + g - y \big) = \big( (a-1)x + by + c \,;\, dx + (e-1)y + g \big). \]
\item \textbf{Imposons la constance.} Pour que cette expression ne dépende ni de $x$ ni de $y$, il faut et il suffit que les coefficients de $x$ et $y$ soient nuls : $a - 1 = 0$, $b = 0$, $d = 0$, $e - 1 = 0$.
\item \textbf{Conclusion.} $a = 1$, $b = 0$, $d = 0$, $e = 1$ : la partie linéaire est bien $I_2$, et le vecteur de la translation est $(c\,;\,g)$. $\blacksquare$
\end{enumerate}
Cette démonstration, très formatrice, est le modèle de toutes les questions du type « reconnaître une transformation à partir de son écriture analytique ».
\end{experience}

\courssub{Réciproque d'une translation}

La translation $t_{\vec{u}}$ est une bijection du plan : intuitivement, si l'on a glissé de $\vec{u}$, on peut « revenir en arrière » en glissant de $-\vec{u}$.

\begin{propriete}[Réciproque d'une translation]
La translation $t_{\vec{u}}$ est une transformation du plan, et sa réciproque est la translation de vecteur opposé :
\[ \left( t_{\vec{u}} \right)^{-1} = t_{-\vec{u}}. \]
Analytiquement, si $t_{\vec{u}} : (x\,;\,y) \mapsto (x + x_u\,;\,y + y_u)$, alors $t_{-\vec{u}} : (x\,;\,y) \mapsto (x - x_u\,;\,y - y_u)$.
\end{propriete}

\begin{experience}[Vérification : la composée donne l'identité]
Vérifions que $t_{-\vec{u}} \circ t_{\vec{u}} = \mathrm{Id}_{\mathcal{P}}$. Partons de $M(x\,;\,y)$.
\begin{enumerate}
\item $t_{\vec{u}}(M) = M_1(x + x_u\,;\,y + y_u)$.
\item $t_{-\vec{u}}(M_1) = \big( (x + x_u) - x_u \,;\, (y + y_u) - y_u \big) = (x\,;\,y) = M$.
\end{enumerate}
Donc $t_{-\vec{u}} \circ t_{\vec{u}}(M) = M$ pour tout point $M$ : la composée est l'identité. On vérifierait de même $t_{\vec{u}} \circ t_{-\vec{u}} = \mathrm{Id}_{\mathcal{P}}$. Ceci prouve que $t_{\vec{u}}$ est bijective, de réciproque $t_{-\vec{u}}$. $\blacksquare$
\end{experience}

\begin{exoresolu}[Translation, image et réciproque]
\textbf{Énoncé.} Dans un repère $(O\,;\vec{\imath},\vec{\jmath})$, on considère l'application $f : (x\,;\,y) \mapsto (x + 3\,;\,y - 2)$.
\begin{enumerate}
\item Montrer que $f$ est une translation dont on précisera le vecteur.
\item Déterminer l'image du point $C(-1\,;\,4)$.
\item Déterminer l'antécédent du point $D(5\,;\,1)$.
\item Donner l'écriture analytique de la réciproque $f^{-1}$.
\end{enumerate}
\tcblower
\textbf{Solution détaillée.}
\begin{enumerate}
\item L'écriture est de la forme $(x\,;\,y) \mapsto (x + a\,;\,y + b)$ avec $a = 3$ et $b = -2$ : la partie linéaire est l'identité. Donc $f$ est la \textbf{translation de vecteur $\vec{u}(3\,;\,-2)$}.
\item $f(C) = (-1 + 3\,;\,4 - 2) = (2\,;\,2)$. Donc $C'(2\,;\,2)$.
\item Chercher l'antécédent de $D$, c'est résoudre $\begin{cases} x + 3 = 5 \\ y - 2 = 1 \end{cases}$, soit $\begin{cases} x = 2 \\ y = 3 \end{cases}$. L'antécédent de $D$ est le point $E(2\,;\,3)$. (On a appliqué $f^{-1}$ à $D$ sans le savoir !)
\item $f^{-1} = t_{-\vec{u}}$ avec $-\vec{u}(-3\,;\,2)$, donc :
\[ f^{-1} : (x\,;\,y) \longmapsto (x - 3\,;\,y + 2). \]
Vérification : $f^{-1}(D) = (5-3\,;\,1+2) = (2\,;\,3) = E$. \checkmark
\end{enumerate}
\end{exoresolu}

\courssub{Reconnaître une translation à partir d'une écriture analytique}

Au Probatoire, on te donnera souvent une application définie par ses formules et on te demandera sa \emph{nature} et ses \emph{éléments caractéristiques}. Voici la démarche systématique.

\begin{methode}[Reconnaître une translation]
Soit $f$ l'application définie analytiquement par $f : (x\,;\,y) \mapsto (ax + by + c\,;\,dx + ey + g)$.
\begin{enumerate}
\item \textbf{Identifier la partie linéaire} : les coefficients de $x$ et $y$ dans chaque formule, c'est-à-dire la matrice $A = \begin{pmatrix} a & b \\ d & e \end{pmatrix}$.
\item \textbf{Tester l'identité} : vérifier que $a = 1$, $e = 1$, $b = 0$, $d = 0$, c'est-à-dire $A = I_2$.
\item \textbf{Conclure} : si oui, $f$ est la translation de vecteur $\vec{u}(c\,;\,g)$ (les termes constants). Si non, ce n'est \emph{pas} une translation (ce pourrait être une homothétie, une rotation, une symétrie...).
\end{enumerate}
\end{methode}

\begin{attention}[Pièges fréquents]
\begin{itemize}
\item \textbf{Confondre les rôles des coefficients.} Dans $(x + 3\,;\,y - 2)$, le vecteur est $(3\,;\,-2)$, pas $(-3\,;\,2)$ ni $(2\,;\,-3)$. Le signe devant $y$ appartient au vecteur !
\item \textbf{Oublier de vérifier la partie linéaire.} L'application $(x\,;\,y) \mapsto (2x + 3\,;\,2y - 2)$ n'est \emph{pas} une translation (le coefficient de $x$ est $2$, pas $1$) : c'est une homothétie de rapport $2$ (nous y reviendrons à la section 3).
\item \textbf{Croire qu'une translation a un point fixe.} Une translation de vecteur \emph{non nul} n'a \textbf{aucun} point fixe : $M' = M$ exigerait $\vec{u} = \vec{0}$. Si tu trouves un point fixe, ta transformation n'est pas une translation non triviale.
\end{itemize}
\end{attention}

\begin{exoresolu}[Reconnaissance d'une application]
\textbf{Énoncé.} On considère les applications définies dans un repère par :
\[ f : (x\,;\,y) \mapsto (x - 5\,;\,y + 7) \qquad \text{et} \qquad g : (x\,;\,y) \mapsto (x + 2y - 5\,;\,y + 7). \]
Pour chacune, dire s'il s'agit d'une translation et préciser, le cas échéant, son vecteur.
\tcblower
\textbf{Solution détaillée.}
\begin{itemize}
\item \textbf{Pour $f$ :} la partie linéaire est $\begin{pmatrix} 1 & 0 \\ 0 & 1 \end{pmatrix} = I_2$ (coefficient de $x$ : $1$ ; de $y$ : $0$ dans la première ligne ; de $x$ : $0$ ; de $y$ : $1$ dans la seconde). Donc $f$ est la translation de vecteur $\vec{u}(-5\,;\,7)$.
\item \textbf{Pour $g$ :} la première formule contient le terme $+2y$ : la partie linéaire est $\begin{pmatrix} 1 & 2 \\ 0 & 1 \end{pmatrix} \neq I_2$. Donc $g$ \textbf{n'est pas une translation}, même si les termes constants $(-5\,;\,7)$ sont les mêmes que pour $f$. (Géométriquement, $g$ est un « cisaillement » composé avec une translation : le déplacement dépend de l'ordonnée du point.)
\end{itemize}
\end{exoresolu}

\courssub{Propriété fondamentale : conservation des vecteurs}

Voici la propriété qui fait toute la richesse géométrique de la translation.

\begin{propriete}[Une translation conserve les vecteurs]
Soit $t_{\vec{u}}$ une translation, $A$ et $B$ deux points, $A' = t_{\vec{u}}(A)$ et $B' = t_{\vec{u}}(B)$. Alors :
\[ \overrightarrow{A'B'} = \overrightarrow{AB}. \]
Autrement dit, l'image d'un vecteur par une translation est \emph{le vecteur lui-même}.
\end{propriete}

\begin{experience}[Démonstration (exigible)]
Par la relation de Chasles :
\[ \overrightarrow{A'B'} = \overrightarrow{A'A} + \overrightarrow{AB} + \overrightarrow{BB'} = -\vec{u} + \overrightarrow{AB} + \vec{u} = \overrightarrow{AB}, \]
car $\overrightarrow{AA'} = \overrightarrow{BB'} = \vec{u}$ par définition de la translation. $\blacksquare$
\end{experience}

De cette propriété découlent immédiatement toutes les conservations que tu connais : une translation conserve les \cle{distances} (car $\|\overrightarrow{A'B'}\| = \|\overrightarrow{AB}\|$), donc les \cle{angles}, les \cle{aires}, le \cle{parallélisme}, l'\cle{orthogonalité}, les \cle{milieux}, l'\cle{alignement}. L'image d'une droite est une droite \emph{parallèle}, l'image d'un cercle de rayon $R$ est un cercle de même rayon $R$ dont le centre est l'image du centre. C'est pourquoi on dit que la translation est une \cle{isométrie} : elle transporte les figures sans les déformer.

\begin{exemplebox}[Justifier une propriété d'une configuration]
Soit $ABCD$ un parallélogramme. Montrons, à l'aide d'une translation, que ses diagonales se coupent en leur milieu.
Puisque $ABCD$ est un parallélogramme, $\overrightarrow{AD} = \overrightarrow{BC}$. Considérons la translation $t$ de vecteur $\overrightarrow{AB}$ : elle envoie $A$ sur $B$, et comme $\overrightarrow{DC} = \overrightarrow{AB}$, elle envoie $D$ sur $C$. La translation conserve les milieux, donc elle envoie le milieu de $[AC]$ sur le milieu de $[BD]$... plus directement : $t$ envoie le segment $[AD]$ sur $[BC]$, donc le milieu de $[AD]$ sur le milieu de $[BC]$. Ce type d'argument — « telle transformation envoie telle figure sur telle figure, donc elle en conserve les propriétés » — est exactement celui attendu au Probatoire pour « utiliser les transformations pour justifier des propriétés des configurations planes ».
\end{exemplebox}

\courssub{Complément pour la Terminale : coordonnées homogènes}

\begin{savaistu}[La translation en une seule matrice : les coordonnées homogènes]
Tu as remarqué que la translation s'écrit $X' = X + U$ : une \emph{somme}, pas un produit matriciel. C'est un cas particulier gênant en informatique, où l'on voudrait écrire \emph{toutes} les transformations sous la forme $X' = MX$ pour pouvoir les composer par simple produit de matrices. L'astuce, que tu rencontreras en Terminale et qui est au cœur de l'infographie, consiste à ajouter une troisième coordonnée fixe égale à $1$ : on représente $M(x\,;\,y)$ par $\begin{pmatrix} x \\ y \\ 1 \end{pmatrix}$. Alors :
\[ \begin{pmatrix} x' \\ y' \\ 1 \end{pmatrix} = \begin{pmatrix} 1 & 0 & x_u \\ 0 & 1 & y_u \\ 0 & 0 & 1 \end{pmatrix} \begin{pmatrix} x \\ y \\ 1 \end{pmatrix} = \begin{pmatrix} x + x_u \\ y + y_u \\ 1 \end{pmatrix}. \]
La translation devient une matrice $3\times 3$ « par blocs » $\begin{pmatrix} I_2 & U \\ 0 & 1 \end{pmatrix}$. Composer deux translations revient alors à multiplier deux matrices : c'est ainsi que les moteurs de jeux vidéo enchaînent des milliers de déplacements par seconde.
\end{savaistu}

\begin{savaistu}[Les translations dans ta poche]
Quand tu fais glisser une photo sur l'écran de ton téléphone, quand un personnage se déplace dans un jeu vidéo, quand un logiciel de CAO (conception assistée par ordinateur) déplace une pièce mécanique sans la tourner ni l'étirer : ce sont des translations, calculées exactement avec les formules $x' = x + x_u$, $y' = y + y_u$ que tu viens d'apprendre. Les cartels publicitaires lumineux de Douala qui font défiler un texte appliquent, image par image, une translation de vecteur constant. Les mathématiques de cette leçon tournent des millions de fois par seconde autour de toi !
\end{savaistu}

\begin{exercice}[Pour t'entraîner]
Dans un repère $(O\,;\vec{\imath},\vec{\jmath})$, on donne les applications :
\[ f_1 : (x\,;\,y) \mapsto (x - 1\,;\,y + 4), \qquad f_2 : (x\,;\,y) \mapsto (-x + 1\,;\,y + 4), \qquad f_3 : (x\,;\,y) \mapsto (x\,;\,y). \]
\begin{enumerate}
\item Lesquelles sont des translations ? Préciser le vecteur.
\item Soit $A(2\,;\,-3)$. Calculer $f_1(A)$, puis l'antécédent de $B(0\,;\,5)$ par $f_1$.
\item Donner l'écriture analytique de $f_1^{-1}$ et vérifier que $f_1^{-1} \circ f_1(A) = A$.
\end{enumerate}
\end{exercice}

\begin{corrige}[Exercice : reconnaissance et réciproque]
\begin{enumerate}
\item $f_1$ : partie linéaire $I_2$ $\Rightarrow$ translation de vecteur $\vec{u}(-1\,;\,4)$. $f_2$ : le coefficient de $x$ est $-1 \neq 1$ $\Rightarrow$ ce n'est pas une translation (sa partie linéaire est $\begin{pmatrix} -1 & 0 \\ 0 & 1 \end{pmatrix}$ : c'est une symétrie orthogonale composée avec une translation). $f_3$ : c'est l'identité, c'est-à-dire la translation de vecteur nul $\vec{0}$.
\item $f_1(A) = (2 - 1\,;\,-3 + 4) = (1\,;\,1)$. Antécédent de $B(0\,;\,5)$ : on résout $x - 1 = 0$ et $y + 4 = 5$, d'où $(x\,;\,y) = (1\,;\,1)$. (Remarque : c'est le point $f_1(A)$ !)
\item $f_1^{-1} : (x\,;\,y) \mapsto (x + 1\,;\,y - 4)$. Vérification : $f_1^{-1}(f_1(A)) = f_1^{-1}(1\,;\,1) = (1 + 1\,;\,1 - 4) = (2\,;\,-3) = A$. \checkmark
\end{enumerate}
\end{corrige}

\begin{aretenir}[L'essentiel de la section]
\begin{itemize}
\item Dans un repère, la translation de vecteur $\vec{u}(x_u\,;\,y_u)$ s'écrit : $\boxed{(x'\,;\,y') = (x + x_u\,;\,y + y_u)}$, soit matriciellement $X' = X + U$.
\item Sa partie linéaire est l'identité $I_2$ (de déterminant $1$) : c'est le \emph{critère de reconnaissance}.
\item Sa réciproque est la translation de vecteur $-\vec{u}$ : $(x\,;\,y) \mapsto (x - x_u\,;\,y - y_u)$.
\item Une translation conserve les vecteurs : $\overrightarrow{A'B'} = \overrightarrow{AB}$ ; c'est donc une isométrie qui conserve distances, angles, aires, milieux et alignements.
\item L'identité est la translation de vecteur nul : ne l'oublie pas dans tes conclusions.
\end{itemize}
\end{aretenir}

\coursec{Définition analytique de l'homothétie}

Dans la section précédente, nous avons vu que la translation se traduit analytiquement par une simple \emph{addition} de vecteur : $X' = X + U$. Nous allons maintenant découvrir une transformation d'un autre type : celle qui \emph{agrandit} ou \emph{réduit} les figures. Tu la connais déjà sans le savoir : quand tu zoomes sur une photo avec ton téléphone, quand un cartographe du MINRESI réalise un plan de Yaoundé à l'échelle $1/25\,000$, ou quand un photographe de quartier agrandit un portrait, c'est une \cle{homothétie} qui est à l'œuvre. Notre objectif : traduire cette idée géométrique en formules, comme nous l'avons fait pour la translation.

\courssub{Rappel géométrique et intuition}

\begin{definition}[Homothétie (rappel géométrique)]
Soit $C$ un point du plan et $k$ un réel \textbf{non nul}. L'\cle{homothétie} de centre $C$ et de rapport $k$, notée $h_{C,k}$, est l'application du plan dans lui-même qui, à tout point $M$, associe le point $M'$ défini par :
\[ \overrightarrow{CM'} = k\,\overrightarrow{CM}. \]
\end{definition}

Lisons attentivement cette égalité vectorielle. Elle dit trois choses à la fois :
\begin{itemize}
\item les points $C$, $M$ et $M'$ sont \textbf{alignés} (les vecteurs sont colinéaires) ;
\item la distance est multipliée par $|k|$ : $CM' = |k|\,CM$ ;
\item si $k > 0$, $M'$ est du \emph{même côté} de $C$ que $M$ ; si $k < 0$, $M'$ est de \emph{l'autre côté} (la figure est « retournée » par rapport au centre).
\end{itemize}

\begin{attention}[Le rapport $k$ ne peut pas être nul]
Si l'on autorisait $k = 0$, alors $\overrightarrow{CM'} = \vec{0}$ pour tout point $M$ : tous les points du plan seraient envoyés sur le centre $C$ ! Cette application constante n'est \textbf{pas bijective} (on ne peut pas « remonter » de $C$ vers $M$), donc ce n'est \textbf{pas une transformation du plan}. Retiens bien : une homothétie exige $k \neq 0$.
\end{attention}

\courssub{Écriture analytique lorsque le centre est l'origine du repère}

Plaçons-nous dans le cas le plus simple : le plan est muni d'un repère $(O\,;\vec{\imath},\vec{\jmath})$ et le centre de l'homothétie est précisément l'origine $O$. Soit $M(x\,;y)$ et $M'(x'\,;y')$ son image par $h_{O,k}$. La relation $\overrightarrow{OM'} = k\,\overrightarrow{OM}$ se traduit coordonnée par coordonnée :

\begin{propriete}[Écriture analytique, centre à l'origine]
L'homothétie de centre $O$ (origine du repère) et de rapport $k \neq 0$ s'écrit :
\[ \begin{cases} x' = kx \\ y' = ky \end{cases} \]
Sous forme matricielle, en posant $X = \begin{pmatrix} x \\ y \end{pmatrix}$ et $X' = \begin{pmatrix} x' \\ y' \end{pmatrix}$ :
\[ X' = kX = \begin{pmatrix} k & 0 \\ 0 & k \end{pmatrix} \begin{pmatrix} x \\ y \end{pmatrix}. \]
La \cle{matrice linéaire} de l'homothétie est donc $kI_2$, où $I_2 = \begin{pmatrix} 1 & 0 \\ 0 & 1 \end{pmatrix}$. Son déterminant vaut $\det(kI_2) = k^2 > 0$ : c'est le facteur par lequel les \textbf{aires} sont multipliées.
\end{propriete}

Deux cas particuliers à connaître par cœur :
\begin{itemize}
\item $k = 1$ : $X' = X$, c'est l'\cle{identité} du plan (rien ne bouge) ;
\item $k = -1$ : $X' = -X$, c'est la \cle{symétrie centrale} de centre $O$, que l'on peut aussi voir comme la rotation de centre $O$ et d'angle $\pi$ ($180°$).
\end{itemize}

\begin{popfigure}
\begin{center}
\begin{tikzpicture}[scale=1.6,>=Stealth]
\draw[->,popmuted] (-1.6,0) -- (2.8,0) node[below left]{$x$};
\draw[->,popmuted] (0,-1.3) -- (0,2.3) node[below left]{$y$};
\coordinate (O) at (0,0);
\coordinate (M) at (1,0.8);
\coordinate (Mp) at (2,1.6);
\coordinate (Mq) at (-0.5,-0.4);
\fill[popblue] (O) circle (1.3pt) node[below left]{$O$};
\fill[popink] (M) circle (1.3pt) node[above right]{$M(1\,;0{,}8)$};
\fill[poporange] (Mp) circle (1.3pt) node[above right]{$M'(2\,;1{,}6)$};
\fill[popgreen] (Mq) circle (1.3pt) node[below left]{$M''(-0{,}5\,;-0{,}4)$};
\draw[->,very thick,popink] (O) -- (M) node[midway,sloped,above,font=\small]{$\overrightarrow{OM}$};
\draw[->,very thick,poporange] (M) -- (Mp);
\draw[->,very thick,popgreen] (O) -- (Mq) node[midway,sloped,below,font=\small]{$-\tfrac12\overrightarrow{OM}$};
\draw[dashed,popmuted] (Mq) -- (Mp);
\end{tikzpicture}
\end{center}
\legende{Homothéties de centre $O$ : $M' = h_{O,2}(M)$ (agrandissement, même côté) et $M'' = h_{O,-1/2}(M)$ (réduction et retournement).}
\end{popfigure}

\courssub{Homothétie de centre quelconque : la formule fondamentale}

En général, le centre $C(c_x\,;c_y)$ n'est pas l'origine du repère. L'astuce consiste à \emph{décomposer le trajet} : pour aller de $O$ à $M'$, on passe par $C$. Par la relation de Chasles :
\[ \overrightarrow{OM'} = \overrightarrow{OC} + \overrightarrow{CM'} = \overrightarrow{OC} + k\,\overrightarrow{CM} = \overrightarrow{OC} + k\left(\overrightarrow{OM} - \overrightarrow{OC}\right). \]
En factorisant, on obtient la forme analytique générale.

\begin{definition}[Homothétie de centre $C$, rapport $k$ (analytique)]
Soit $C$ de vecteur position $C = \begin{pmatrix} c_x \\ c_y \end{pmatrix}$ et $k \neq 0$. L'homothétie $h_{C,k}$ admet l'écriture analytique :
\[ X' = kX + (1-k)C, \qquad \text{c'est-à-dire} \qquad \begin{cases} x' = kx + (1-k)c_x \\ y' = ky + (1-k)c_y \end{cases} \]
\end{definition}

\begin{popfigure}
\begin{center}
\begin{tikzpicture}[scale=1.5,>=Stealth]
\draw[->,popmuted] (-0.6,0) -- (4.4,0) node[below left]{$x$};
\draw[->,popmuted] (0,-0.6) -- (0,3.2) node[below left]{$y$};
\coordinate (O) at (0,0);
\coordinate (C) at (1,0.6);
\coordinate (M) at (2.2,1.4);
\coordinate (Mp) at (3.4,2.2);
\fill[popblue] (O) circle (1.3pt) node[below left]{$O$};
\fill[poppurple] (C) circle (1.3pt) node[below left]{$C$};
\fill[popink] (M) circle (1.3pt) node[above]{$M$};
\fill[poporange] (Mp) circle (1.3pt) node[above right]{$M'$};
\draw[->,thick,poppurple] (O) -- (C) node[midway,below,font=\small]{$\overrightarrow{OC}$};
\draw[->,thick,popink] (C) -- (M) node[midway,sloped,above,font=\small]{$\overrightarrow{CM}$};
\draw[->,very thick,poporange] (M) -- (Mp) node[midway,sloped,above,font=\small]{};
\draw[->,thick,popgreen,dashed] (O) -- (Mp) node[pos=0.6,sloped,below,font=\small]{$\overrightarrow{OM'}$};
\node[poporange,font=\small] at (3.05,1.62) {$k\,\overrightarrow{CM}$};
\end{tikzpicture}
\end{center}
\legende{Centre $C$ distinct de l'origine : $\overrightarrow{OM'} = \overrightarrow{OC} + k\,\overrightarrow{CM}$, ici avec $k = 2$.}
\end{popfigure}

Vérifions la cohérence : si $X = C$, alors $X' = kC + (1-k)C = C$. Le centre est bien un \textbf{point fixe} — et c'est le seul quand $k \neq 1$.

\begin{propriete}[Réciproque d'une homothétie]
L'homothétie $h_{C,k}$ ($k \neq 0$) est bijective. Sa réciproque est l'homothétie de \textbf{même centre} et de rapport $\dfrac{1}{k}$ :
\[ h_{C,k}^{-1} = h_{C,\,1/k}. \]
En effet, $\overrightarrow{CM'} = k\,\overrightarrow{CM}$ équivaut à $\overrightarrow{CM} = \dfrac{1}{k}\,\overrightarrow{CM'}$.
\end{propriete}

\courssub{Reconnaître une homothétie à partir d'une écriture analytique}

Au Probatoire, on te donnera souvent une application définie par des formules et on te demandera sa \emph{nature} et ses \emph{éléments caractéristiques}. Voici le critère décisif.

\begin{propriete}[Caractérisation analytique d'une homothétie]
Une application affine $f : X \mapsto AX + B$ est une homothétie si et seulement si sa matrice linéaire est un multiple non nul de l'identité : $A = kI_2$ avec $k \neq 0$.
\begin{itemize}
\item Si $k \neq 1$, $f$ est l'homothétie de rapport $k$ et de centre l'unique point fixe $C$, solution de $C = kC + B$, c'est-à-dire $C = \dfrac{B}{1-k}$.
\item Si $k = 1$ et $B = \vec{0}$, $f$ est l'identité ; si $k = 1$ et $B \neq \vec{0}$, $f$ est une \textbf{translation}, pas une homothétie.
\end{itemize}
\end{propriete}

\begin{experience}[Démonstration guidée : le centre est l'unique point fixe]
Montrons que, pour $k \neq 1$, l'application $f(X) = kX + B$ admet un unique point fixe et qu'il s'agit d'une homothétie.
\begin{enumerate}
\item Un point fixe vérifie $X = f(X)$, c'est-à-dire $X = kX + B$, soit $(1-k)X = B$.
\item Comme $k \neq 1$, on peut diviser par $1 - k$ : il existe une \textbf{unique} solution $C = \dfrac{B}{1-k}$.
\item Réécrivons alors $f$ autour de $C $ : pour tout $X$,
\[ f(X) - C = kX + B - C = kX + (1-k)C - C = k(X - C). \]
\item Ainsi $\overrightarrow{Cf(M)} = k\,\overrightarrow{CM}$ : $f$ est bien l'homothétie de centre $C$ et de rapport $k$. $\square$
\end{enumerate}
\end{experience}

\begin{methode}[Déterminer le centre et le rapport d'une homothétie donnée analytiquement]
\begin{enumerate}
\item \textbf{Identifier $k$} : c'est le coefficient commun sur la diagonale de la matrice (le facteur devant $x$ dans $x'$, qui doit être le même que devant $y$ dans $y'$).
\item \textbf{Vérifier la forme} : les termes croisés (en $y$ dans $x'$, en $x$ dans $y'$) doivent être nuls, sinon ce n'est pas une homothétie.
\item \textbf{Trouver le centre} (si $k \neq 1$) : résoudre le système des points fixes $\begin{cases} x = kx + b_1 \\ y = ky + b_2 \end{cases}$, soit $c_x = \dfrac{b_1}{1-k}$ et $c_y = \dfrac{b_2}{1-k}$.
\item \textbf{Conclure} en rédigeant : « $f$ est l'homothétie de centre $C(c_x\,;c_y)$ et de rapport $k$ ».
\end{enumerate}
\end{methode}

\begin{exemplebox}[Exemple chiffré complet]
Soit $f : (x\,;y) \mapsto (2x+1\,;\;2y-3)$.
\begin{itemize}
\item La matrice linéaire est $A = \begin{pmatrix} 2 & 0 \\ 0 & 2 \end{pmatrix} = 2I_2$ : c'est bien un multiple non nul de l'identité, donc $f$ est une homothétie de rapport $k = 2$.
\item Centre : on résout $x = 2x + 1$ d'où $x = -1$ ; et $y = 2y - 3$ d'où $y = 3$. Le centre est $C(-1\,;3)$.
\item Vérification par la formule : $(1-k)C = -(C) = (1\,;-3)$, et l'on retrouve bien $x' = 2x + 1$, $y' = 2y - 3$.
\item Réciproque : $f^{-1} = h_{C,\,1/2}$, d'écriture $X' = \dfrac{1}{2}X + \dfrac{1}{2}C$, soit $\begin{cases} x' = \dfrac{x - 1}{2} \\[2pt] y' = \dfrac{y + 3}{2} \end{cases}$.
\end{itemize}
\end{exemplebox}

\begin{attention}[Pièges fréquents au Probatoire]
\begin{itemize}
\item \textbf{Signe de $k$} : si $k < 0$, l'image est de l'autre côté du centre. Ne dis jamais qu'une homothétie de rapport négatif « agrandit » sans préciser qu'elle retourne aussi la figure.
\item \textbf{Confusion translation/homothétie} : $(x\,;y) \mapsto (x+1\,;y-3)$ est une translation (matrice $I_2$, donc $k = 1$) ; $(x\,;y) \mapsto (2x+1\,;2y-3)$ est une homothétie ($k = 2$). Regarde d'abord la matrice !
\item \textbf{Centre mal calculé} : le centre n'est généralement \emph{pas} le point $(b_1\,;b_2)$. Il faut résoudre le système des points fixes.
\item $k = 0$ n'est \textbf{jamais} une homothétie : l'application serait constante, non bijective.
\end{itemize}
\end{attention}

\begin{exoresolu}[Reconnaître et caractériser une homothétie]
Le plan est muni d'un repère $(O\,;\vec{\imath},\vec{\jmath})$. On considère l'application $g$ qui à tout point $M(x\,;y)$ associe $M'(x'\,;y')$ tel que $x' = -3x + 8$ et $y' = -3y - 4$.
\begin{enumerate}
\item Montrer que $g$ est une homothétie dont on précisera le rapport.
\item Déterminer les coordonnées de son centre $\Omega$.
\item Donner l'écriture analytique de la réciproque $g^{-1}$.
\end{enumerate}
\tcblower
\textbf{Solution.}
\begin{enumerate}
\item La matrice linéaire de $g$ est $A = \begin{pmatrix} -3 & 0 \\ 0 & -3 \end{pmatrix} = -3I_2$. C'est un multiple non nul de l'identité, donc $g$ est une homothétie de rapport $k = -3$.
\item Le centre est l'unique point fixe. On résout :
\[ x = -3x + 8 \iff 4x = 8 \iff x = 2 \quad ; \quad y = -3y - 4 \iff 4y = -4 \iff y = -1. \]
Donc $\Omega(2\,;-1)$. Vérification : $(1-k)\Omega = 4\Omega = (8\,;-4)$, conforme à l'écriture de $g$.
\item La réciproque est $h_{\Omega,\,-1/3}$ : $X' = -\dfrac{1}{3}X + \left(1 + \dfrac{1}{3}\right)\Omega = -\dfrac{1}{3}X + \dfrac{4}{3}\Omega$, soit
\[ \begin{cases} x' = -\dfrac{x}{3} + \dfrac{8}{3} \\[4pt] y' = -\dfrac{y}{3} - \dfrac{4}{3} \end{cases} \]
\end{enumerate}
\end{exoresolu}

\begin{savaistu}[Des homothéties aux fractales]
Les \cle{fractales} — ces objets géométriques infiniment détaillés comme le \emph{flocon de Koch} ou l'\emph{ensemble de Mandelbrot} — se construisent en itérant des homothéties (souvent combinées à des rotations et translations) : on parle de \emph{systèmes de fonctions itérées}. Le flocon de Koch, par exemple, s'obtient en reproduisant indéfiniment un motif réduit par des homothéties de rapport $\tfrac{1}{3}$. Ces objets, loin d'être de simples curiosités, servent à modéliser les côtes découpées (comme celles du littoral atlantique camerounais vers Kribi), les réseaux de téléphonie mobile ou encore les images compressées. Les homothéties que tu étudies aujourd'hui sont littéralement les briques de ces constructions modernes.
\end{savaistu}

\begin{aretenir}[L'essentiel de la section]
\begin{itemize}
\item Homothétie de centre $C$, rapport $k \neq 0$ : $\overrightarrow{CM'} = k\,\overrightarrow{CM}$, soit analytiquement $X' = kX + (1-k)C$.
\item Centre à l'origine : $x' = kx$, $y' = ky$ ; matrice $kI_2$, déterminant $k^2$.
\item Cas particuliers : $k = 1$ (identité), $k = -1$ (symétrie centrale).
\item Réciproque : homothétie de même centre, rapport $\dfrac{1}{k}$.
\item Reconnaissance : $f(X) = AX + B$ est une homothétie ssi $A = kI_2$, $k \neq 0$ ; le centre est l'unique point fixe $C = \dfrac{B}{1-k}$ (pour $k \neq 1$).
\end{itemize}
\end{aretenir}

\begin{exercice}[Pour t'entraîner]
Dans un repère $(O\,;\vec{\imath},\vec{\jmath})$, on donne les applications :
\[ f_1 : \begin{cases} x' = 3x - 4 \\ y' = 3y + 2 \end{cases} \qquad f_2 : \begin{cases} x' = 2x + y \\ y' = 2y - 1 \end{cases} \qquad f_3 : \begin{cases} x' = x + 5 \\ y' = y - 2 \end{cases} \]
Pour chacune, dire s'il s'agit d'une homothétie, d'une translation ou d'aucune des deux ; dans le cas d'une homothétie, préciser le rapport, le centre et l'écriture analytique de la réciproque.
\end{exercice}

\begin{corrige}[Exercice : Pour t'entraîner]
\begin{itemize}
\item $f_1$ : matrice $3I_2$, donc homothétie de rapport $k = 3$. Centre : $x = 3x - 4 \Rightarrow x = 2$ ; $y = 3y + 2 \Rightarrow y = -1$. Centre $C(2\,;-1)$. Réciproque : $h_{C,\,1/3}$, soit $x' = \dfrac{x}{3} + \dfrac{4}{3}$, $y' = \dfrac{y}{3} - \dfrac{2}{3}$.
\item $f_2$ : la matrice $\begin{pmatrix} 2 & 1 \\ 0 & 2 \end{pmatrix}$ n'est pas un multiple de l'identité (terme $y$ dans $x'$) : ce n'est ni une homothétie ni une translation.
\item $f_3$ : matrice $I_2$ avec $B = (5\,;-2) \neq \vec{0}$ : c'est la translation de vecteur $\vec{u}(5\,;-2)$, pas une homothétie.
\end{itemize}
\end{corrige}

\coursec{Composée de deux rotations de même centre}

Dans la section précédente, nous avons donné la définition analytique d'une homothétie : dans un repère de centre $O$, le point $M(x\,;y)$ a pour image $M'(kx\,;ky)$. Nous savons donc maintenant décrire une transformation par des formules sur les coordonnées. Nous allons exploiter ce point de vue analytique pour répondre à une question très naturelle : \emph{que se passe-t-il lorsqu'on enchaîne deux rotations de même centre ?} Intuitivement, tourner de $30$ degrés puis encore de $45$ degrés autour du même point revient à tourner de $75$ degrés. Nous allons démontrer cette intuition, lui donner un cadre algébrique précis (matrices, nombres complexes), et en tirer toutes les conséquences utiles au Probatoire.

\courssub{Rappel analytique : l'expression d'une rotation de centre $O$}

Plaçons-nous dans le plan muni d'un repère orthonormé direct $(O\,;\vec{i},\vec{j})$. Rappelons ce qu'est une rotation.

\begin{definition}[Rotation de centre $O$ et d'angle $\alpha$]
Soit $O$ un point du plan et $\alpha$ un réel. La \cle{rotation} de centre $O$ et d'angle $\alpha$, notée $r_{(O,\alpha)}$ ou simplement $R_\alpha$ lorsque le centre est fixé, est la transformation du plan qui :
\begin{itemize}
\item laisse $O$ invariant : $r(O)=O$ ;
\item à tout point $M$ distinct de $O$ associe le point $M'$ tel que
\[ OM' = OM \qquad \text{et} \qquad \left(\overrightarrow{OM}\,,\overrightarrow{OM'}\right) = \alpha \ [2\pi]. \]
\end{itemize}
\end{definition}

\begin{propriete}[Expression analytique et matricielle d'une rotation de centre $O$]
Dans un repère orthonormé direct $(O\,;\vec{i},\vec{j})$ d'origine $O$, si $M(x\,;y)$ a pour image $M'(x'\,;y')$ par la rotation de centre $O$ et d'angle $\alpha$, alors
\[
\begin{cases} x' = x\cos\alpha - y\sin\alpha \\ y' = x\sin\alpha + y\cos\alpha \end{cases}
\qquad\text{c'est-à-dire}\qquad
\begin{pmatrix} x' \\ y' \end{pmatrix} = R_\alpha \begin{pmatrix} x \\ y \end{pmatrix},
\quad \text{avec} \quad
R_\alpha = \begin{pmatrix} \cos\alpha & -\sin\alpha \\ \sin\alpha & \cos\alpha \end{pmatrix}.
\]
La matrice $R_\alpha$ est appelée \cle{matrice de rotation} d'angle $\alpha$.
\end{propriete}

Cette matrice a une structure remarquable qu'il faut savoir reconnaître à l'œil : la première colonne est $(\cos\alpha\,;\sin\alpha)$, image du vecteur $\vec{i}$, et la seconde $(-\sin\alpha\,;\cos\alpha)$, image du vecteur $\vec{j}$. En particulier, $R_0 = \begin{pmatrix} 1 & 0 \\ 0 & 1 \end{pmatrix}$ est la matrice identité, ce qui correspond bien au fait qu'une rotation d'angle nul ne bouge rien.

\begin{attention}[Bien vérifier le centre]
La formule matricielle ci-dessus n'est valable que si le centre de la rotation est l'\textbf{origine du repère}. Si le centre est un point $\Omega(a\,;b)$ quelconque, il faut d'abord se ramener à $\Omega$ par translation des coordonnées : poser $X = x - a$, $Y = y - b$, appliquer la matrice, puis revenir. Dans toute cette section, le centre commun des rotations est $O$, origine du repère.
\end{attention}

\courssub{Composée de deux rotations de même centre : le théorème}

\begin{experience}[Découverte : enchaîner deux rotations]
Sur une feuille, place un point $O$ et un point $M$ tel que $OM = 4$ cm. À l'aide d'un rapporteur :
\begin{enumerate}
\item construis $M_1$, image de $M$ par la rotation de centre $O$ et d'angle $\alpha = 40$ degrés ;
\item construis $M_2$, image de $M_1$ par la rotation de centre $O$ et d'angle $\beta = 50$ degrés ;
\item mesure l'angle $\left(\overrightarrow{OM}\,,\overrightarrow{OM_2}\right)$.
\end{enumerate}
Tu mesures $90$ degrés, c'est-à-dire $\alpha + \beta$. De plus $OM_2 = OM_1 = OM$ : les distances à $O$ sont conservées à chaque étape. Tout se passe comme si $M_2$ était l'image de $M$ par une unique rotation de centre $O$ et d'angle $\alpha+\beta$. Reste à le démontrer pour tous les points et tous les angles.
\end{experience}

\begin{propriete}[Théorème — Composée de deux rotations de même centre (démonstration exigible)]
Soit $O$ un point du plan, $\alpha$ et $\beta$ deux réels. La composée de la rotation de centre $O$ et d'angle $\alpha$ suivie de la rotation de centre $O$ et d'angle $\beta$ est la rotation de centre $O$ et d'angle $\alpha + \beta$ :
\[ r_{(O,\beta)} \circ r_{(O,\alpha)} = r_{(O,\alpha+\beta)}. \]
En particulier, cette composée est commutative : $r_{(O,\beta)} \circ r_{(O,\alpha)} = r_{(O,\alpha)} \circ r_{(O,\beta)}$.
\end{propriete}

\begin{experience}[Démonstration guidée par le calcul matriciel]
Notons $R_\alpha$ et $R_\beta$ les matrices des deux rotations dans un repère orthonormé direct d'origine $O$. Pour un point $M$ de coordonnées $\begin{pmatrix} x \\ y \end{pmatrix}$, l'image par la première rotation a pour coordonnées $R_\alpha \begin{pmatrix} x \\ y \end{pmatrix}$, puis l'image finale a pour coordonnées $R_\beta R_\alpha \begin{pmatrix} x \\ y \end{pmatrix}$. Calculons le produit $R_\beta R_\alpha$.

\textbf{Étape 1 : poser le produit.}
\[
R_\beta R_\alpha =
\begin{pmatrix} \cos\beta & -\sin\beta \\ \sin\beta & \cos\beta \end{pmatrix}
\begin{pmatrix} \cos\alpha & -\sin\alpha \\ \sin\alpha & \cos\alpha \end{pmatrix}.
\]

\textbf{Étape 2 : calculer chaque coefficient.}
\begin{align*}
R_\beta R_\alpha &=
\begin{pmatrix}
\cos\beta\cos\alpha - \sin\beta\sin\alpha & -\cos\beta\sin\alpha - \sin\beta\cos\alpha \\
\sin\beta\cos\alpha + \cos\beta\sin\alpha & -\sin\beta\sin\alpha + \cos\beta\cos\alpha
\end{pmatrix}.
\end{align*}

\textbf{Étape 3 : appliquer les formules d'addition trigonométriques.}
Rappelons : $\cos(a+b) = \cos a\cos b - \sin a\sin b$ et $\sin(a+b) = \sin a\cos b + \cos a\sin b$. On obtient :
\[
R_\beta R_\alpha =
\begin{pmatrix}
\cos(\alpha+\beta) & -\sin(\alpha+\beta) \\
\sin(\alpha+\beta) & \cos(\alpha+\beta)
\end{pmatrix}
= R_{\alpha+\beta}.
\]

\textbf{Étape 4 : conclure.} La composée a pour matrice $R_{\alpha+\beta}$, qui est exactement la matrice de la rotation de centre $O$ et d'angle $\alpha+\beta$. Comme deux applications ayant même expression analytique sont égales, on a bien
\[ r_{(O,\beta)} \circ r_{(O,\alpha)} = r_{(O,\alpha+\beta)}. \qquad \blacksquare \]
\end{experience}

\begin{experience}[Autre démonstration, par les nombres complexes]
Dans le plan complexe d'origine $O$, la rotation de centre $O$ et d'angle $\alpha$ s'écrit $z' = e^{i\alpha}z$. Enchaînons : $M(z) \xrightarrow{r_\alpha} M_1(z_1 = e^{i\alpha}z) \xrightarrow{r_\beta} M_2(z_2 = e^{i\beta}z_1)$. Alors
\[ z_2 = e^{i\beta}\left(e^{i\alpha}z\right) = e^{i\beta}e^{i\alpha}z = e^{i(\alpha+\beta)}z, \]
car l'exponentielle complexe transforme la somme en produit. On reconnaît l'écriture complexe de la rotation de centre $O$ et d'angle $\alpha+\beta$. $\blacksquare$
\end{experience}

\begin{aretenir}[L'idée centrale]
Composer deux rotations de \textbf{même centre}, c'est \textbf{additionner leurs angles}. Le centre ne change pas, les angles s'ajoutent. C'est l'analogue géométrique de l'addition des nombres réels — et c'est exactement ce que traduisent les formules d'addition trigonométriques.
\end{aretenir}

\begin{popfigure}
\begin{center}
\begin{tikzpicture}[scale=2.2,>=Stealth]
\coordinate (O) at (0,0);
\coordinate (M) at (1.3,0);
\coordinate (M1) at (50:1.3);
\coordinate (M2) at (110:1.3);
\draw[popline] (-0.35,0) -- (1.6,0);
\draw[popline] (0,-0.35) -- (0,1.6);
\draw[popmuted] (O) -- (M);
\draw[popblue,thick] (O) -- (M1);
\draw[poporange,thick] (O) -- (M2);
\draw[popblue,->,thick] (1.05,0) arc (0:50:1.05) node[midway,right] {$\alpha$};
\draw[poporange,->,thick] (50:0.85) arc (50:110:0.85) node[midway,above left] {$\beta$};
\draw[popgreen,->,thick] (0.55,0) arc (0:110:0.55) node[midway,above right,pos=0.75] {$\alpha+\beta$};
\fill[popdark] (O) circle (0.018) node[below left] {$O$};
\fill[popdark] (M) circle (0.018) node[below right] {$M$};
\fill[popblue] (M1) circle (0.018) node[right] {$M_1 = r_{(O,\alpha)}(M)$};
\fill[poporange] (M2) circle (0.018) node[above left] {$M_2 = r_{(O,\beta)}(M_1)$};
\draw[dashed,popmuted] (M) -- (M1) -- (M2);
\end{tikzpicture}
\end{center}
\legende{Enchaînement de deux rotations de même centre $O$ : $M_1 = r_{(O,\alpha)}(M)$, puis $M_2 = r_{(O,\beta)}(M_1)$. L'angle total $\left(\overrightarrow{OM},\overrightarrow{OM_2}\right)$ vaut $\alpha+\beta$ : la composée est la rotation $r_{(O,\alpha+\beta)}$.}
\end{popfigure}

\courssub{Lecture matricielle : le produit de deux matrices de rotation}

Le cœur de la démonstration peut se résumer en un tableau, à savoir refaire rapidement le jour du Probatoire.

\begin{center}
\begin{tabularx}{\linewidth}{|l|X|X|}\hline
\rowcolor{popblueL} \textbf{Objet} & \textbf{Rotation d'angle $\alpha$} & \textbf{Rotation d'angle $\beta$} \\ \hline
Matrice & $R_\alpha = \begin{pmatrix} \cos\alpha & -\sin\alpha \\ \sin\alpha & \cos\alpha \end{pmatrix}$ & $R_\beta = \begin{pmatrix} \cos\beta & -\sin\beta \\ \sin\beta & \cos\beta \end{pmatrix}$ \\ \hline
\rowcolor{popgreenL} \multicolumn{3}{|l|}{\textbf{Produit (composée $r_\beta \circ r_\alpha$)}} \\ \hline
Produit brut & \multicolumn{2}{l|}{$\begin{pmatrix} \cos\beta\cos\alpha - \sin\beta\sin\alpha & -(\cos\beta\sin\alpha + \sin\beta\cos\alpha) \\ \sin\beta\cos\alpha + \cos\beta\sin\alpha & \cos\beta\cos\alpha - \sin\beta\sin\alpha \end{pmatrix}$} \\ \hline
Après formules d'addition & \multicolumn{2}{l|}{$R_\beta R_\alpha = \begin{pmatrix} \cos(\alpha+\beta) & -\sin(\alpha+\beta) \\ \sin(\alpha+\beta) & \cos(\alpha+\beta) \end{pmatrix} = R_{\alpha+\beta}$} \\ \hline
Conclusion & \multicolumn{2}{l|}{$r_{(O,\beta)} \circ r_{(O,\alpha)} = r_{(O,\alpha+\beta)}$ : rotation de centre $O$, d'angle $\alpha+\beta$.} \\ \hline
\end{tabularx}
\end{center}

\courssub{Réduction modulo $2\pi$ et cas particuliers}

Un angle de rotation n'est défini qu'à un multiple de $2\pi$ près : tourner de $\dfrac{9\pi}{4}$ ou de $\dfrac{\pi}{4}$, c'est la même transformation. Il faut donc toujours \emph{réduire} la somme $\alpha+\beta$ modulo $2\pi$ avant de conclure.

\begin{propriete}[Cas particuliers fondamentaux]
Soit $r_{(O,\alpha)}$ et $r_{(O,\beta)}$ deux rotations de même centre $O$.
\begin{itemize}
\item Si $\alpha + \beta \equiv 0 \ [2\pi]$, alors $r_{(O,\beta)} \circ r_{(O,\alpha)} = \mathrm{Id}_P$, l'\cle{identité} du plan. En particulier, la réciproque de $r_{(O,\alpha)}$ est $r_{(O,-\alpha)}$.
\item Si $\alpha + \beta \equiv \pi \ [2\pi]$, alors $r_{(O,\beta)} \circ r_{(O,\alpha)}$ est la \cle{symétrie centrale} de centre $O$, c'est-à-dire la rotation d'angle $\pi$, encore égale à l'homothétie de centre $O$ et de rapport $-1$.
\end{itemize}
\end{propriete}

Vérifions le second point matriciellement : $R_\pi = \begin{pmatrix} \cos\pi & -\sin\pi \\ \sin\pi & \cos\pi \end{pmatrix} = \begin{pmatrix} -1 & 0 \\ 0 & -1 \end{pmatrix}$, qui envoie $(x\,;y)$ sur $(-x\,;-y)$ : c'est bien la symétrie centrale de centre $O$, et aussi l'homothétie $h(O,-1)$.

\begin{attention}[Erreur fréquente : oublier la réduction modulo $2\pi$]
Beaucoup d'élèves concluent trop vite. Par exemple, pour $r_{\left(O,\frac{3\pi}{2}\right)} \circ r_{\left(O,\frac{\pi}{2}\right)}$, la somme des angles vaut $\dfrac{3\pi}{2}+\dfrac{\pi}{2} = 2\pi$. Dire « c'est la rotation d'angle $2\pi$ » est une réponse maladroite : comme $2\pi \equiv 0 \ [2\pi]$, la composée est l'\textbf{identité}. De même, une somme égale à $3\pi$ doit être réduite en $\pi$ : c'est une symétrie centrale. Toujours ramener l'angle dans $]-\pi\,;\pi]$ ou $[0\,;2\pi[$ avant de nommer la transformation.
\end{attention}

\begin{methode}[Caractériser la composée de deux rotations de même centre]
\begin{enumerate}
\item \textbf{Identifier} le centre commun $O$ et les deux angles $\alpha$ et $\beta$ (attention à l'ordre d'écriture : dans $r_\beta \circ r_\alpha$, c'est $r_\alpha$ qui agit en premier, mais ici l'ordre n'a pas d'importance puisque $\alpha+\beta = \beta+\alpha$).
\item \textbf{Additionner} les angles : $\alpha + \beta$.
\item \textbf{Réduire modulo $2\pi$} : écrire $\alpha+\beta \equiv \theta \ [2\pi]$ avec $\theta \in \ ]-\pi\,;\pi]$.
\item \textbf{Conclure} sur la nature et les éléments caractéristiques :
\begin{itemize}
\item si $\theta = 0$ : identité ;
\item si $\theta = \pi$ : symétrie centrale de centre $O$ ;
\item sinon : rotation de centre $O$ et d'angle $\theta$.
\end{itemize}
\end{enumerate}
\end{methode}

\begin{exemplebox}[Exemples chiffrés]
\textbf{Exemple 1.} $r_{\left(O,\frac{\pi}{3}\right)} \circ r_{\left(O,\frac{\pi}{6}\right)}$ : la somme des angles est $\dfrac{\pi}{3}+\dfrac{\pi}{6} = \dfrac{2\pi+\pi}{6} = \dfrac{\pi}{2}$. La composée est la rotation de centre $O$ et d'angle $\dfrac{\pi}{2}$ (le « quart de tour direct »).

\textbf{Exemple 2.} $r_{\left(O,\frac{2\pi}{3}\right)} \circ r_{\left(O,\frac{4\pi}{3}\right)}$ : la somme vaut $\dfrac{2\pi}{3}+\dfrac{4\pi}{3} = \dfrac{6\pi}{3} = 2\pi \equiv 0 \ [2\pi]$. La composée est l'\textbf{identité} du plan : les deux rotations sont réciproques l'une de l'autre.

\textbf{Exemple 3.} $r_{\left(O,\frac{3\pi}{4}\right)} \circ r_{\left(O,\frac{\pi}{4}\right)}$ : la somme vaut $\pi$. La composée est la \textbf{symétrie centrale} de centre $O$.
\end{exemplebox}

\courssub{Structure de groupe : une propriété de synthèse}

\begin{propriete}[L'ensemble des rotations de centre $O$ est un groupe]
Notons $\mathcal{R}_O$ l'ensemble des rotations de centre $O$. Muni de la composition $\circ$, c'est un \cle{groupe commutatif} :
\begin{itemize}
\item \textbf{stabilité} : $r_{(O,\alpha)} \circ r_{(O,\beta)} = r_{(O,\alpha+\beta)} \in \mathcal{R}_O$ ;
\item \textbf{associativité} : la composition des applications est toujours associative ;
\item \textbf{élément neutre} : $r_{(O,0)} = \mathrm{Id}_P$ ;
\item \textbf{symétrique} : la réciproque de $r_{(O,\alpha)}$ est $r_{(O,-\alpha)}$ ;
\item \textbf{commutativité} : $r_{(O,\alpha)} \circ r_{(O,\beta)} = r_{(O,\beta)} \circ r_{(O,\alpha)}$ car $\alpha+\beta = \beta+\alpha$.
\end{itemize}
L'application $\theta \mapsto r_{(O,\theta)}$ réalise un \cle{isomorphisme} entre le groupe $(\mathbb{R}/2\pi\mathbb{Z},+)$ (les angles modulo $2\pi$) et $(\mathcal{R}_O,\circ)$ : composer les rotations revient exactement à additionner les angles.
\end{propriete}

\begin{attention}[Centres distincts : tout change]
La commutativité et la belle formule « les angles s'additionnent » ne sont vraies que parce que les deux rotations ont le \textbf{même centre}. Si les centres sont distincts, la composée de deux rotations d'angles $\alpha$ et $\beta$ est encore une rotation d'angle $\alpha+\beta$ lorsque $\alpha+\beta \not\equiv 0 \ [2\pi]$, mais son centre doit être déterminé par construction, et la composée n'est plus commutative. Pire : si $\alpha+\beta \equiv 0 \ [2\pi]$ avec des centres distincts, la composée est une \textbf{translation}, pas l'identité. Cette étude est hors programme en Première, mais retiens l'essentiel : \emph{le centre commun est l'hypothèse qui fait tout fonctionner}.
\end{attention}

\begin{exoresolu}[Composée de deux rotations et image d'un point]
Dans le plan muni d'un repère orthonormé direct $(O\,;\vec{i},\vec{j})$, on considère les rotations $r_1 = r_{\left(O,\frac{\pi}{3}\right)}$ et $r_2 = r_{\left(O,\frac{\pi}{6}\right)}$, ainsi que le point $A(2\,;0)$.
\begin{enumerate}
\item Déterminer la nature et les éléments caractéristiques de $r_2 \circ r_1$.
\item Déterminer les coordonnées de $A'$, image de $A$ par $r_2 \circ r_1$.
\item Montrer que $r_1^6 = r_1 \circ r_1 \circ r_1 \circ r_1 \circ r_1 \circ r_1$ est l'identité.
\end{enumerate}
\tcblower
\textbf{1. Nature de la composée.} Les deux rotations ont le même centre $O$. D'après le théorème, $r_2 \circ r_1$ est la rotation de centre $O$ et d'angle
\[ \frac{\pi}{6}+\frac{\pi}{3} = \frac{\pi}{6}+\frac{2\pi}{6} = \frac{3\pi}{6} = \frac{\pi}{2}. \]
Donc $r_2 \circ r_1 = r_{\left(O,\frac{\pi}{2}\right)}$ : le quart de tour direct de centre $O$.

\textbf{2. Coordonnées de $A'$.} La matrice de la rotation d'angle $\dfrac{\pi}{2}$ est
\[ R_{\pi/2} = \begin{pmatrix} \cos\frac{\pi}{2} & -\sin\frac{\pi}{2} \\ \sin\frac{\pi}{2} & \cos\frac{\pi}{2} \end{pmatrix} = \begin{pmatrix} 0 & -1 \\ 1 & 0 \end{pmatrix}. \]
Alors
\[ \begin{pmatrix} x' \\ y' \end{pmatrix} = \begin{pmatrix} 0 & -1 \\ 1 & 0 \end{pmatrix} \begin{pmatrix} 2 \\ 0 \end{pmatrix} = \begin{pmatrix} 0 \\ 2 \end{pmatrix}. \]
Donc $A'(0\,;2)$. Vérification géométrique : $OA' = OA = 2$ et l'angle $\left(\overrightarrow{OA},\overrightarrow{OA'}\right)$ vaut bien $\dfrac{\pi}{2}$.

\textbf{3. Puissance sixième de $r_1$.} En composant six fois la rotation d'angle $\dfrac{\pi}{3}$, les angles s'additionnent : $r_1^6$ est la rotation de centre $O$ et d'angle $6 \times \dfrac{\pi}{3} = 2\pi \equiv 0 \ [2\pi]$. Donc $r_1^6 = \mathrm{Id}_P$. On dit que $r_1$ est d'\emph{ordre} $6$ : six rotations de $60$ degrés ramènent chaque point à sa position initiale.
\end{exoresolu}

\begin{exercice}[À toi de jouer]
On considère les rotations $r = r_{\left(O,\frac{5\pi}{6}\right)}$ et $r' = r_{\left(O,\frac{\pi}{6}\right)}$.
\begin{enumerate}
\item Déterminer la nature et les éléments caractéristiques de $r' \circ r$.
\item Même question pour $r \circ r'$ puis pour $r \circ r$.
\item Déterminer le plus petit entier $n \geq 1$ tel que $r'^n = \mathrm{Id}_P$.
\end{enumerate}
\end{exercice}

\begin{corrige}[Exercice 1]
\textbf{1.} Somme des angles : $\dfrac{\pi}{6}+\dfrac{5\pi}{6} = \dfrac{6\pi}{6} = \pi$. La composée $r' \circ r$ est la \textbf{symétrie centrale} de centre $O$.

\textbf{2.} $r \circ r' = r' \circ r$ (même centre, commutativité) : c'est la même symétrie centrale. Pour $r \circ r$ : somme $2 \times \dfrac{5\pi}{6} = \dfrac{5\pi}{3}$, que l'on réduit : $\dfrac{5\pi}{3} \equiv -\dfrac{\pi}{3} \ [2\pi]$. Donc $r \circ r = r_{\left(O,-\frac{\pi}{3}\right)}$, rotation de centre $O$ et d'angle $-\dfrac{\pi}{3}$.

\textbf{3.} $r'^n$ est la rotation d'angle $\dfrac{n\pi}{6}$. On cherche $\dfrac{n\pi}{6} \equiv 0 \ [2\pi]$, c'est-à-dire $\dfrac{n\pi}{6} = 2k\pi$, soit $n = 12k$. Le plus petit entier convenable est $n = 12$ : la rotation d'angle $\dfrac{\pi}{6}$ est d'ordre $12$.
\end{corrige}

\begin{savaistu}[Des rotations partout autour de nous]
En robotique, un bras articulé est une chaîne de rotations : lorsque deux articulations tournent autour du \emph{même axe} (par exemple l'épaule d'un robot soudeur dans une usine d'assemblage), l'orientation finale du poignet s'obtient en additionnant les angles — exactement le théorème de cette section. Les ingénieurs utilisent d'ailleurs les matrices $R_\alpha$ pour programmer ces mouvements. Plus près de nous, l'aiguille des minutes d'une horloge effectue une rotation d'angle $-\dfrac{\pi}{30}$ par minute (sens indirect, d'où le signe moins) : au bout de $60$ minutes, l'angle total est $-2\pi \equiv 0 \ [2\pi]$, et l'aiguille est revenue à sa position — une belle illustration de la réduction modulo $2\pi$ et de la notion d'ordre d'une rotation.
\end{savaistu}

\begin{aretenir}[Bilan de la section]
\begin{itemize}
\item Rotation de centre $O$, angle $\alpha$ : matrice $R_\alpha = \begin{pmatrix} \cos\alpha & -\sin\alpha \\ \sin\alpha & \cos\alpha \end{pmatrix}$, ou écriture complexe $z' = e^{i\alpha}z$.
\item \textbf{Théorème} : $r_{(O,\beta)} \circ r_{(O,\alpha)} = r_{(O,\alpha+\beta)}$ — composer, c'est additionner les angles.
\item Toujours \textbf{réduire modulo $2\pi$} : somme $\equiv 0 \ [2\pi] \Rightarrow$ identité ; somme $\equiv \pi \ [2\pi] \Rightarrow$ symétrie centrale de centre $O$.
\item $(\mathcal{R}_O, \circ)$ est un groupe commutatif, isomorphe à $(\mathbb{R}/2\pi\mathbb{Z}, +)$ ; la réciproque de $r_{(O,\alpha)}$ est $r_{(O,-\alpha)}$.
\end{itemize}
Dans la section suivante, nous verrons que les homothéties de même centre jouissent d'une propriété tout aussi élégante : composer deux homothéties de même centre, c'est \emph{multiplier} leurs rapports.
\end{aretenir}

\coursec{Composée de deux homothéties de même centre}

Dans la section précédente, nous avons découvert une belle surprise : composer deux rotations de même centre revient à \emph{additionner} leurs angles. La question se pose alors naturellement : que se passe-t-il lorsqu'on compose deux \cle{homothéties} de même centre ? Faut-il additionner les rapports ? Les multiplier ? Autre chose ? C'est toute l'ambition de cette section : montrer que les homothéties de même centre se comportent, pour la composition, exactement comme les nombres réels non nuls pour la multiplication. Autrement dit, la loi des rapports est la \emph{multiplication}, tout comme la loi des angles était l'addition pour les rotations.

\courssub{Rappels sur l'homothétie et sa définition analytique}

\begin{definition}[Homothétie de centre $O$ et de rapport $k$]
Soit $O$ un point du plan et $k$ un nombre réel non nul. L'\cle{homothétie} de centre $O$ et de rapport $k$, notée $H_{O,k}$, est la transformation du plan qui à tout point $M$ associe le point $M'$ tel que :
\[
\overrightarrow{OM'} = k\,\overrightarrow{OM}.
\]
Le point $O$ est le seul point fixe de $H_{O,k}$ lorsque $k \neq 1$.
\end{definition}

Plaçons-nous dans un repère $(O\,;\vec{\imath},\vec{\jmath})$ dont l'origine est précisément le centre $O$ de l'homothétie. Si $M$ a pour coordonnées $(x\,;y)$ et $M'$ pour coordonnées $(x'\,;y')$, la relation vectorielle $\overrightarrow{OM'} = k\,\overrightarrow{OM}$ se traduit par la \cle{définition analytique} :
\[
\begin{cases} x' = kx \\ y' = ky \end{cases}
\]
ce qui s'écrit matriciellement :
\[
\begin{pmatrix} x' \\ y' \end{pmatrix} = \begin{pmatrix} k & 0 \\ 0 & k \end{pmatrix} \begin{pmatrix} x \\ y \end{pmatrix} = kI_2 \begin{pmatrix} x \\ y \end{pmatrix},
\]
où $I_2 = \begin{pmatrix} 1 & 0 \\ 0 & 1 \end{pmatrix}$ est la matrice identité. La matrice d'une homothétie de centre $O$ est donc $kI_2$ : elle dilate (ou contracte) toutes les coordonnées du même facteur $k$.

\begin{aretenir}[Cas particuliers à connaître par cœur]
\begin{itemize}
\item Si $k = 1$ : $H_{O,1}$ est l'\cle{identité} du plan (chaque point est sa propre image).
\item Si $k = -1$ : $H_{O,-1}$ est la \cle{symétrie centrale} de centre $O$.
\item Si $|k| > 1$ : c'est un agrandissement ; si $0 < |k| < 1$ : c'est une réduction.
\item Si $k < 0$ : l'image est « de l'autre côté » du centre $O$.
\end{itemize}
\end{aretenir}

\courssub{Intuition : deux zooms successifs}

Imagine que tu photographies le marché de Mokolo avec ton téléphone. Tu appliques d'abord un zoom de facteur $3$, puis sur la photo obtenue, un zoom de facteur $2$. Le résultat final est identique à un zoom unique de facteur $6$ : les facteurs se sont \emph{multipliés}, pas additionnés. De même, si tu zoomes de facteur $2$ puis que tu « dézoomes » de facteur $\tfrac{1}{2}$, tu reviens exactement à l'image de départ : $2 \times \tfrac{1}{2} = 1$, et le rapport $1$ correspond à l'identité. Cette intuition de la vie courante est exactement ce que les mathématiques vont confirmer.

\courssub{Théorème fondamental}

\begin{propriete}[Théorème : composée de deux homothéties de même centre]
Soit $O$ un point du plan, et $k_1$, $k_2$ deux réels non nuls. La composée
\[
H_{O,k_2} \circ H_{O,k_1}
\]
est l'homothétie de centre $O$ et de rapport $k_1 k_2$ :
\[
H_{O,k_2} \circ H_{O,k_1} = H_{O,\,k_1 k_2}.
\]
En particulier, cette composée est commutative : $H_{O,k_2} \circ H_{O,k_1} = H_{O,k_1} \circ H_{O,k_2}$. \emph{Cette démonstration est exigible au Probatoire.}
\end{propriete}

\begin{experience}[Démonstration guidée du théorème]
Soit $M$ un point quelconque du plan. Notons $M_1 = H_{O,k_1}(M)$ puis $M_2 = H_{O,k_2}(M_1)$. Nous devons déterminer la nature de l'application $M \mapsto M_2$.

\textbf{Étape 1 : traduction vectorielle des deux homothéties.}
Par définition de $H_{O,k_1}$ : $\overrightarrow{OM_1} = k_1\,\overrightarrow{OM}$.
Par définition de $H_{O,k_2}$ : $\overrightarrow{OM_2} = k_2\,\overrightarrow{OM_1}$.

\textbf{Étape 2 : substitution.}
En injectant la première relation dans la seconde :
\[
\overrightarrow{OM_2} = k_2\,\overrightarrow{OM_1} = k_2\left(k_1\,\overrightarrow{OM}\right) = (k_1 k_2)\,\overrightarrow{OM}.
\]

\textbf{Étape 3 : conclusion.}
Le rapport $k_1 k_2$ est non nul (produit de deux réels non nuls). La relation $\overrightarrow{OM_2} = (k_1 k_2)\,\overrightarrow{OM}$ signifie exactement que $M_2 = H_{O,\,k_1 k_2}(M)$. Ceci étant vrai pour tout point $M$, on a bien :
\[
H_{O,k_2} \circ H_{O,k_1} = H_{O,\,k_1 k_2}.
\]

\textbf{Vérification analytique (seconde preuve).} Dans un repère d'origine $O$, si $M(x\,;y)$, alors $M_1(k_1 x\,; k_1 y)$ puis $M_2\big(k_2(k_1 x)\,; k_2(k_1 y)\big) = \big(k_1 k_2\, x\,; k_1 k_2\, y\big)$. On reconnaît la définition analytique de $H_{O,\,k_1 k_2}$. $\blacksquare$
\end{experience}

\begin{methode}[Déterminer la composée de deux homothéties de même centre]
\begin{enumerate}
\item Vérifier que les deux homothéties ont \textbf{le même centre} $O$ (sinon le théorème ne s'applique pas tel quel).
\item \textbf{Multiplier} les rapports : $k = k_1 \times k_2$.
\item Vérifier que $k_1 k_2 \neq 0$ (c'est automatique si $k_1 \neq 0$ et $k_2 \neq 0$).
\item Conclure : $H_{O,k_2} \circ H_{O,k_1} = H_{O,\,k_1 k_2}$, puis identifier les cas particuliers ($k_1k_2 = 1$ : identité ; $k_1k_2 = -1$ : symétrie centrale de centre $O$).
\end{enumerate}
\end{methode}

\begin{attention}[Piège classique : additionner au lieu de multiplier]
L'erreur la plus fréquente au Probatoire est de \textbf{confondre les deux lois} :
\begin{itemize}
\item Pour les \textbf{rotations} de même centre : on \emph{additionne} les angles : $r_{O,\theta_2} \circ r_{O,\theta_1} = r_{O,\,\theta_1 + \theta_2}$.
\item Pour les \textbf{homothéties} de même centre : on \emph{multiplie} les rapports : $H_{O,k_2} \circ H_{O,k_1} = H_{O,\,k_1 k_2}$.
\end{itemize}
Ainsi $H_{O,3} \circ H_{O,2} = H_{O,6}$ et surtout pas $H_{O,5}$ ! Retiens l'image du zoom : deux zooms successifs multiplient les grossissements.
\end{attention}

\courssub{Exemples chiffrés et cas particuliers}

\begin{exemplebox}[Premiers calculs de composées]
\textbf{a)} $H_{O,3} \circ H_{O,-\tfrac{1}{2}}$ : le rapport de la composée est $3 \times \left(-\dfrac{1}{2}\right) = -\dfrac{3}{2}$. Donc :
\[
H_{O,3} \circ H_{O,-\tfrac{1}{2}} = H_{O,\,-\tfrac{3}{2}}.
\]
C'est une homothétie de rapport négatif : l'image est agrandie (facteur $\tfrac{3}{2}$) et renversée par rapport à $O$.

\textbf{b)} $H_{O,4} \circ H_{O,\tfrac{1}{4}}$ : le rapport est $4 \times \dfrac{1}{4} = 1$. Donc :
\[
H_{O,4} \circ H_{O,\tfrac{1}{4}} = H_{O,1} = \mathrm{Id}.
\]
Les deux homothéties sont \emph{réciproques} l'une de l'autre : la seconde « annule » la première.
\end{exemplebox}

\begin{propriete}[Cas particuliers remarquables]
Soit $k_1, k_2 \in \mathbb{R}^*$.
\begin{itemize}
\item Si $k_1 k_2 = 1$ (c'est-à-dire $k_2 = \dfrac{1}{k_1}$), alors $H_{O,k_2} \circ H_{O,k_1} = \mathrm{Id}$ : les deux homothéties sont réciproques l'une de l'autre. Autrement dit :
\[
\left(H_{O,k}\right)^{-1} = H_{O,\,\tfrac{1}{k}}.
\]
\item Si $k_1 k_2 = -1$, alors $H_{O,k_2} \circ H_{O,k_1} = H_{O,-1}$ : la composée est la \cle{symétrie centrale} de centre $O$.
\end{itemize}
\end{propriete}

\begin{exoresolu}[Composer et reconnaître]
\textbf{Énoncé.} Dans le plan muni d'un repère d'origine $O$, on considère le point $A(2\,;-1)$. On note $A_1 = H_{O,2}(A)$, puis $A_2 = H_{O,-1}(A_1)$.
\begin{enumerate}
\item Déterminer les coordonnées de $A_1$ puis de $A_2$.
\item Déterminer la nature et les éléments caractéristiques de $H_{O,-1} \circ H_{O,2}$.
\end{enumerate}
\tcblower
\textbf{Solution.}
\begin{enumerate}
\item $A_1 = H_{O,2}(A)$ : on multiplie les coordonnées par $2$ : $A_1(4\,;-2)$.
Puis $A_2 = H_{O,-1}(A_1)$ : on multiplie par $-1$ : $A_2(-4\,;2)$.
\item D'après le théorème, $H_{O,-1} \circ H_{O,2} = H_{O,\,2 \times (-1)} = H_{O,-2}$. C'est l'homothétie de centre $O$ et de rapport $-2$. Vérification directe : $H_{O,-2}(A)$ a pour coordonnées $(-2 \times 2\,; -2 \times (-1)) = (-4\,;2)$, qui sont bien celles de $A_2$. $\blacksquare$
\end{enumerate}
\end{exoresolu}

\begin{popfigure}
\begin{center}
\begin{tikzpicture}[scale=1.15,>=Stealth]
\draw[popline] (-4.6,0) -- (4.6,0);
\fill[popdark] (0,0) circle (1.6pt) node[below left=1pt] {$O$};
% Cas identité : k1=2 puis k2=0.5
\fill[popblue] (1,0.9) circle (1.6pt) node[above] {$M$};
\fill[poporange] (2,1.8) circle (1.6pt) node[above] {$M_1=H_{O,2}(M)$};
\draw[->,poporange,thick] (1,0.9) -- (2,1.8);
\draw[->,popgreen,thick] (2,1.8) -- (1,0.9);
\node[popgreen] at (1.9,1.15) {\small $H_{O,\,0.5}$};
\node[popdark,align=center] at (2.6,-0.9) {\small $k_1=2$ puis $k_2=0.5$ :\\ \small $M_2=M$ (identité)};
% Cas symétrie centrale : k1=2 puis k2=-1
\fill[popblue] (1,-1.7) circle (1.6pt) node[below left] {$N$};
\fill[poporange] (2,-3.4) circle (1.6pt) node[below] {$N_1=H_{O,2}(N)$};
\fill[poppurple] (-2,3.4) circle (1.6pt) node[above] {$N_2=H_{O,-1}(N_1)$};
\draw[->,poporange,thick] (1,-1.7) -- (2,-3.4);
\draw[->,poppurple,thick] (2,-3.4) -- (-2,3.4);
\draw[dashed,popmuted] (-2,3.4) -- (2,-3.4);
\node[popdark,align=center] at (-3.4,-2.6) {\small $k_1=2$ puis $k_2=-1$ :\\ \small $N_2$ symétrique de $N_1$\\ \small rapport $-2$};
\end{tikzpicture}
\end{center}
\legende{Composition de deux homothéties de centre $O$. En haut : $H_{O,0.5} \circ H_{O,2} = \mathrm{Id}$. En bas : $H_{O,-1} \circ H_{O,2} = H_{O,-2}$, l'image $N_2$ est de l'autre côté de $O$.}
\end{popfigure}

\courssub{Structure de groupe des homothéties de centre $O$}

Le théorème précédent a une conséquence profonde : les homothéties de centre $O$ forment une structure algébrique remarquable, calquée sur $(\mathbb{R}^*, \times)$.

\begin{propriete}[Groupe des homothéties de centre $O$]
L'ensemble $\mathcal{H}_O = \left\{ H_{O,k} \mid k \in \mathbb{R}^* \right\}$, muni de la composition $\circ$, est un \cle{groupe commutatif} :
\begin{itemize}
\item \textbf{Stabilité} : $H_{O,k_2} \circ H_{O,k_1} = H_{O,k_1 k_2} \in \mathcal{H}_O$ car $k_1 k_2 \neq 0$ ;
\item \textbf{Associativité} : la composition des applications est toujours associative ;
\item \textbf{Élément neutre} : $H_{O,1} = \mathrm{Id}$ ;
\item \textbf{Symétrique} : toute homothétie $H_{O,k}$ admet une réciproque, $\left(H_{O,k}\right)^{-1} = H_{O,\,1/k}$ ;
\item \textbf{Commutativité} : $H_{O,k_1} \circ H_{O,k_2} = H_{O,k_2} \circ H_{O,k_1}$ car $k_1 k_2 = k_2 k_1$.
\end{itemize}
L'application $k \mapsto H_{O,k}$ réalise un \emph{isomorphisme} entre $(\mathbb{R}^*, \times)$ et $(\mathcal{H}_O, \circ)$ : composer des homothéties, c'est multiplier leurs rapports.
\end{propriete}

\begin{popfigure}
\begin{center}
\begin{tikzpicture}[>=Stealth,node distance=1.1cm]
\node[draw=popblue,fill=popblueL,rounded corners,thick,minimum width=2.6cm] (h1) {$H_{O,k_1}$};
\node[draw=poporange,fill=poporangeL,rounded corners,thick,minimum width=2.6cm,right=1.6cm of h1] (h2) {$H_{O,k_2}$};
\node[draw=poppurple,fill=poppurpleL,rounded corners,thick,minimum width=3.2cm,right=1.6cm of h2] (comp) {$H_{O,\,k_1 k_2}$};
\draw[->,thick,popdark] (h1) -- node[above]{\small puis} (h2);
\draw[->,thick,popdark] (h2) -- node[above]{\small $=$} (comp);
\node[draw=popgreen,fill=popgreenL,rounded corners,thick,minimum width=2.6cm,below=1.5cm of h1] (hk) {$H_{O,k}$};
\node[draw=popgreen,fill=popgreenL,rounded corners,thick,minimum width=2.6cm,right=1.6cm of hk] (hki) {$H_{O,\,1/k}$};
\node[draw=popgold,fill=popgoldL,rounded corners,thick,minimum width=3.2cm,right=1.6cm of hki] (id) {$\mathrm{Id} = H_{O,1}$};
\draw[->,thick,popdark] (hk) -- node[above]{\small réciproque} (hki);
\draw[->,thick,popdark] (hki) -- node[above]{\small $=$} (id);
\end{tikzpicture}
\end{center}
\legende{Schéma récapitulatif : la composée de deux homothéties de même centre a pour rapport le produit des rapports ; la réciproque de $H_{O,k}$ est $H_{O,1/k}$.}
\end{popfigure}

\courssub{Extension : composée de $n$ homothéties de même centre}

Par récurrence immédiate, le résultat se généralise à une chaîne de $n$ homothéties de même centre.

\begin{propriete}[Composée de $n$ homothéties de même centre]
Soit $O$ un point du plan et $k_1, k_2, \ldots, k_n$ des réels non nuls. Alors :
\[
H_{O,k_n} \circ \cdots \circ H_{O,k_2} \circ H_{O,k_1} = H_{O,\,k_1 k_2 \cdots k_n}.
\]
En particulier, la puissance $n$-ième d'une homothétie est $\left(H_{O,k}\right)^n = H_{O,\,k^n}$.
\end{propriete}

\begin{exoresolu}[Chaîne d'homothéties]
\textbf{Énoncé.} Un photographe de Yaoundé retouche une image en appliquant successivement trois homothéties de même centre $O$ : $H_{O,2}$, puis $H_{O,-3}$, puis $H_{O,\tfrac{1}{4}}$. Déterminer la nature et les éléments caractéristiques de la transformation globale.
\tcblower
\textbf{Solution.} Le rapport global est le produit des rapports :
\[
k = 2 \times (-3) \times \frac{1}{4} = -\frac{6}{4} = -\frac{3}{2}.
\]
La transformation globale est donc $H_{O,\,-\tfrac{3}{2}}$ : une homothétie de centre $O$, de rapport $-\dfrac{3}{2}$. L'image finale est agrandie d'un facteur $\dfrac{3}{2}$ et renversée par rapport à $O$ (rapport négatif). $\blacksquare$
\end{exoresolu}

\begin{exercice}[À toi de jouer]
Soit $O$ un point du plan.
\begin{enumerate}
\item Déterminer la nature et les éléments caractéristiques de $H_{O,-5} \circ H_{O,\tfrac{1}{5}}$.
\item Déterminer le réel $k$ tel que $H_{O,k} \circ H_{O,6} = H_{O,-1}$. Quelle est alors la nature de la composée ?
\item Montrer que $\left(H_{O,\sqrt{2}}\right)^4$ est une homothétie dont on précisera le rapport.
\end{enumerate}
\end{exercice}

\begin{corrige}[Exercice : À toi de jouer]
\begin{enumerate}
\item Rapport : $-5 \times \dfrac{1}{5} = -1$. La composée est $H_{O,-1}$, c'est-à-dire la \textbf{symétrie centrale de centre $O$}.
\item On veut $6k = -1$, donc $k = -\dfrac{1}{6}$. La composée $H_{O,-1/6} \circ H_{O,6} = H_{O,-1}$ est la symétrie centrale de centre $O$.
\item $\left(H_{O,\sqrt{2}}\right)^4 = H_{O,\,(\sqrt{2})^4} = H_{O,4}$ car $(\sqrt{2})^4 = 4$. C'est l'homothétie de centre $O$ et de rapport $4$.
\end{enumerate}
\end{corrige}

\begin{savaistu}[Les lentilles composent leurs grossissements]
En optique géométrique, une lentille mince produit d'un objet une image dont la taille est multipliée par un facteur appelé \emph{grandissement}. Lorsqu'on place plusieurs lentilles de même centre optique à la suite — comme dans les objectifs d'appareils photo ou les microscopes utilisés dans les laboratoires de biologie des lycées —, le grandissement total est le \emph{produit} des grandissements de chaque lentille. C'est exactement la loi de composition des homothéties de même centre ! C'est ce principe qui permet aux microscopes d'atteindre des grossissements de plusieurs centaines en combinant un objectif ($\times 40$) et un oculaire ($\times 10$) : $40 \times 10 = 400$.
\end{savaistu}

\begin{aretenir}[L'essentiel de la section]
\begin{itemize}
\item $H_{O,k_2} \circ H_{O,k_1} = H_{O,\,k_1 k_2}$ : on \textbf{multiplie} les rapports (alors qu'on \emph{additionne} les angles des rotations).
\item Si $k_1 k_2 = 1$ : la composée est l'\textbf{identité} ; $\left(H_{O,k}\right)^{-1} = H_{O,1/k}$.
\item Si $k_1 k_2 = -1$ : la composée est la \textbf{symétrie centrale} de centre $O$.
\item L'ensemble des homothéties de centre $O$ et de rapports non nuls est un \textbf{groupe commutatif} isomorphe à $(\mathbb{R}^*, \times)$.
\item Pour $n$ homothéties de même centre : le rapport de la composée est le \textbf{produit} des $n$ rapports.
\end{itemize}
\end{aretenir}

\coursec{Composée d'une homothétie et d'une translation}

Dans la section précédente, nous avons vu que composer deux homothéties de \emph{même centre} donne une homothétie de même centre, de rapport produit. Mais que se passe-t-il lorsqu'on enchaîne une homothétie avec une \emph{translation} ? C'est exactement ce que fait votre téléphone lorsque, après avoir zoomé sur une photo, vous la faites glisser avec le doigt : le résultat est encore un zoom... mais centré ailleurs ! Cette intuition simple — \og le zoom glisse \fg{} — est le cœur de cette section. Nous allons la transformer en théorème rigoureux, démontré analytiquement, puis apprendre à déterminer systématiquement la nature et les éléments caractéristiques de la composée.

\courssub{Mise en place : calcul analytique de la composée}

Le plan est rapporté à un repère $(O\,;\vec{\imath},\vec{\jmath})$. On note $H_{C,k}$ l'homothétie de centre $C(c\,;d)$ et de rapport $k\neq 0$, et $t_{\vec{u}}$ la translation de vecteur $\vec{u}(a\,;b)$.

\textbf{Premier ordre : $H_{C,k}$ puis $t_{\vec{u}}$, c'est-à-dire $t_{\vec{u}}\circ H_{C,k}$.}

Soit $M(x\,;y)$ un point du plan. Notons $M_1 = H_{C,k}(M)$ puis $M' = t_{\vec{u}}(M_1)$.

Par définition analytique de l'homothétie (section 3) :
\[ M_1 = kM + (1-k)C \quad\text{c'est-à-dire}\quad
\begin{cases} x_1 = kx + (1-k)c \\ y_1 = ky + (1-k)d \end{cases} \]

Puis par définition analytique de la translation :
\[ M' = M_1 + \vec{u} \quad\text{c'est-à-dire}\quad
\begin{cases} x' = x_1 + a \\ y' = y_1 + b \end{cases} \]

En substituant :
\[ \boxed{\; M' = kM + (1-k)C + \vec{u} \;} \]

\textbf{Second ordre : $t_{\vec{u}}$ puis $H_{C,k}$, c'est-à-dire $H_{C,k}\circ t_{\vec{u}}$.}

Cette fois $M_1 = M + \vec{u}$, puis $M' = kM_1 + (1-k)C$, d'où :
\[ \boxed{\; M' = k\bigl(M + \vec{u}\bigr) + (1-k)C = kM + k\vec{u} + (1-k)C \;} \]

\begin{attention}[Deux ordres, deux formules différentes]
Les expressions de $t_{\vec{u}}\circ H_{C,k}$ et de $H_{C,k}\circ t_{\vec{u}}$ ne coïncident pas en général : dans le premier cas le vecteur ajouté est $\vec{u}$, dans le second c'est $k\vec{u}$. \textbf{L'ordre de composition compte.} En revanche, dans les deux cas, le coefficient devant $M$ est le même : $k$. C'est lui qui va dicter la nature de la composée.
\end{attention}

\courssub{Le théorème fondamental}

Dans les deux cas, la composée s'écrit sous la forme $M' = kM + \vec{V}$, où $\vec{V}$ est un vecteur constant (indépendant de $M$). Reconnaissons cette écriture : c'est exactement la forme analytique d'une homothétie de rapport $k$... à condition que $k\neq 1$ !

\begin{propriete}[Composée d'une homothétie et d'une translation (démonstration exigible)]
Soit $H_{C,k}$ une homothétie de rapport $k\neq 1$ et $t_{\vec{u}}$ une translation. Alors :
\begin{itemize}
\item $t_{\vec{u}}\circ H_{C,k}$ est l'\textbf{homothétie de rapport $k$} dont le centre $\Omega_1$ est l'unique point fixe, donné par $\Omega_1 = C + \dfrac{\vec{u}}{1-k}$ ;
\item $H_{C,k}\circ t_{\vec{u}}$ est l'\textbf{homothétie de rapport $k$} dont le centre $\Omega_2$ est l'unique point fixe, donné par $\Omega_2 = C + \dfrac{k\vec{u}}{1-k}$.
\end{itemize}
De façon unifiée et infaillible : \textbf{dans les deux cas, la composée est une homothétie de rapport $k$, dont le centre est l'unique point fixe de la composée.}
\end{propriete}

\begin{experience}[Démonstration guidée du théorème]
Traitons le cas $f = t_{\vec{u}}\circ H_{C,k}$. On a établi : $M' = kM + (1-k)C + \vec{u}$.

\textbf{Étape 1 : chercher un point fixe.} Un point $\Omega$ est fixe par $f$ si $f(\Omega) = \Omega$, c'est-à-dire :
\[ \Omega = k\Omega + (1-k)C + \vec{u} \iff (1-k)\Omega = (1-k)C + \vec{u}. \]
Comme $k\neq 1$, on peut diviser par $1-k$ :
\[ \Omega = C + \frac{\vec{u}}{1-k}. \]
Il existe donc un \emph{unique} point fixe $\Omega_1 = C + \frac{\vec{u}}{1-k}$.

\textbf{Étape 2 : montrer que $f$ est l'homothétie de centre $\Omega_1$ et de rapport $k$.} Calculons, pour tout point $M$ d'image $M' = f(M)$ :
\begin{align*}
\overrightarrow{\Omega_1 M'} &= M' - \Omega_1 = \bigl(kM + (1-k)C + \vec{u}\bigr) - \bigl(k\Omega_1 + (1-k)C + \vec{u}\bigr) \\
&= k(M - \Omega_1) = k\,\overrightarrow{\Omega_1 M}.
\end{align*}
Par définition, $f$ est donc l'homothétie de centre $\Omega_1$ et de rapport $k$. $\blacksquare$

\textit{Remarque importante :} on rencontre parfois l'énoncé \og le centre de $t_{\vec{u}}\circ H_{C,k}$ est l'image de $C$ par la translation \fg{}. Cette affirmation est \textbf{fausse} (sauf si $k=0$, exclu). La formulation rigoureuse est : \textbf{le centre est l'unique point fixe de la composée}, qu'il faut calculer via $\Omega = \frac{\vec{V}}{1-k}$.
\end{experience}

\begin{aretenir}[Ce qu'il faut retenir absolument]
\begin{itemize}
\item $k \neq 1$ : homothétie $\circ$ translation $=$ \cle{homothétie de rapport $k$}, de centre l'\cle{unique point fixe} (qui n'est \textbf{pas} $C$ en général : le centre a \og glissé \fg{}).
\item Le rapport de la composée ne dépend que de $k$ ; la translation ne fait que \emph{déplacer le centre}.
\item $t_{\vec{u}}\circ H_{C,k}$ et $H_{C,k}\circ t_{\vec{u}}$ ont le \emph{même rapport} $k$ mais des \emph{centres différents} en général : la composition n'est pas commutative.
\end{itemize}
\end{aretenir}

\courssub{Cas particulier $k = 1$ : composée de deux translations}

Que devient la formule $M' = kM + \vec{V}$ lorsque $k = 1$ ? Elle se réduit à $M' = M + \vec{V}$ : c'est une translation ! Or une homothétie de rapport $1$ n'est autre que l'identité, donc composer $H_{C,1}$ avec $t_{\vec{u}}$ revient simplement à appliquer $t_{\vec{u}}$. Le cas vraiment intéressant est la composée de \emph{deux} translations.

\begin{propriete}[Composée de deux translations]
La composée de deux translations $t_{\vec{u}}$ et $t_{\vec{v}}$ est la translation de vecteur \cle{somme} $\vec{u}+\vec{v}$ :
\[ t_{\vec{v}}\circ t_{\vec{u}} = t_{\vec{u}+\vec{v}} = t_{\vec{u}}\circ t_{\vec{v}}. \]
En particulier, les translations \textbf{commutent} entre elles (contrairement aux homothéties et translations entre elles).
\end{propriete}

\begin{experience}[Démonstration (très simple, exigible)]
Soit $M$ un point, $M_1 = t_{\vec{u}}(M)$ et $M' = t_{\vec{v}}(M_1)$. Alors :
\[ \overrightarrow{MM'} = \overrightarrow{MM_1} + \overrightarrow{M_1M'} = \vec{u} + \vec{v}. \]
Donc $M' = t_{\vec{u}+\vec{v}}(M)$ pour tout point $M$. $\blacksquare$
\end{experience}

\begin{exemplebox}[Un déplacement en deux temps à Yaoundé]
Un moto-taxi part du rond-point Poste Centrale, effectue un premier trajet modélisé par $\vec{u}(3\,;-1)$ (en centaines de mètres, repère orthonormé), puis un second trajet $\vec{v}(-1\,;4)$. Le déplacement global est la translation de vecteur $\vec{u}+\vec{v} = (2\,;3)$ : comme si le taxi avait fait le trajet directement. Et peu importe l'ordre des deux trajets !
\end{exemplebox}

\courssub{Méthode générale de détermination de la nature et des éléments}

\begin{methode}[Étudier $H\circ t$ ou $t\circ H$]
\begin{enumerate}
\item \textbf{Écrire la composée sous la forme analytique} $M' = kM + \vec{V}$, en composant les expressions de chaque transformation (attention à l'ordre !).
\item \textbf{Regarder le coefficient $k$} :
\begin{itemize}
\item si $k = 1$ : la composée est la \textbf{translation} de vecteur $\vec{V}$ ;
\item si $k \neq 1$ : la composée est l'\textbf{homothétie de rapport $k$}.
\end{itemize}
\item \textbf{Déterminer le centre} (cas $k\neq 1$) en résolvant l'équation du point fixe :
\[ \Omega = k\Omega + \vec{V} \iff \Omega = \frac{\vec{V}}{1-k}. \]
\item \textbf{Conclure proprement} : nommer la transformation, donner son rapport et les coordonnées de son centre.
\end{enumerate}
\end{methode}

\begin{exoresolu}[Homothétie puis translation : calcul complet]
Le plan est muni d'un repère orthonormé $(O\,;\vec{\imath},\vec{\jmath})$. Soit $H$ l'homothétie de centre $O$ et de rapport $2$, et $t$ la translation de vecteur $\vec{u}(3\,;1)$. Déterminer la nature et les éléments caractéristiques de $f = t\circ H$.
\tcblower
\textbf{Étape 1 : écriture analytique.} Pour $M(x\,;y)$, $H(M) = (2x\,;2y)$ puis $t$ ajoute $\vec{u}$ :
\[ f(M) = \bigl(2x+3\,;\,2y+1\bigr), \quad\text{soit}\quad M' = 2M + \vec{u}. \]
On reconnaît la forme $M' = kM + \vec{V}$ avec $k = 2$ et $\vec{V} = (3\,;1)$.

\textbf{Étape 2 : nature.} $k = 2 \neq 1$ : $f$ est une \textbf{homothétie de rapport $2$}.

\textbf{Étape 3 : centre.} On résout $\Omega = 2\Omega + \vec{u}$, c'est-à-dire $-\Omega = \vec{u}$, donc :
\[ \Omega = \frac{\vec{u}}{1-2} = -\vec{u} = (-3\,;-1). \]

\textbf{Conclusion.} $f = t\circ H$ est l'homothétie de centre $\Omega(-3\,;-1)$ et de rapport $2$.

\textbf{Vérification.} $f(\Omega) = \bigl(2\times(-3)+3\,;\,2\times(-1)+1\bigr) = (-3\,;-1) = \Omega$. Le point est bien fixe. $\checkmark$
\end{exoresolu}

\begin{exoresolu}[Translation puis homothétie : l'ordre change le centre]
Avec les mêmes $H$ et $t$ que ci-dessus, déterminer la nature et les éléments caractéristiques de $g = H\circ t$.
\tcblower
\textbf{Étape 1.} $t(M) = (x+3\,;\,y+1)$, puis $H$ multiplie par $2$ :
\[ g(M) = \bigl(2x+6\,;\,2y+2\bigr), \quad\text{soit}\quad M' = 2M + 2\vec{u}. \]
\textbf{Étape 2.} $k = 2 \neq 1$ : $g$ est une \textbf{homothétie de rapport $2$} (même rapport que $f$).
\textbf{Étape 3.} Point fixe : $\Omega' = 2\Omega' + 2\vec{u} \iff \Omega' = \dfrac{2\vec{u}}{1-2} = -2\vec{u} = (-6\,;-2)$.
\textbf{Conclusion.} $g$ est l'homothétie de centre $\Omega'(-6\,;-2)$ et de rapport $2$. On constate que $\Omega' \neq \Omega$ : $t\circ H \neq H\circ t$, bien que les deux aient le même rapport.
\end{exoresolu}

\courssub{Interprétation géométrique : le centre qui glisse}

Géométriquement, tout se passe comme si la translation \og faisait glisser \fg{} le centre de l'homothétie. Appliquer $H_{C,k}$ puis $t_{\vec{u}}$ revient à appliquer une seule homothétie de même rapport, mais centrée en un point $\Omega$ décalé par rapport à $C$. La figure suivante illustre le trajet d'un point $M$ : d'abord $M' = H_{C,2}(M)$, puis $M'' = t_{\vec{u}}(M')$, et l'on vérifie que $M''$ est aussi l'image de $M$ par l'homothétie de centre $\Omega$ et de rapport $2$.

\begin{popfigure}
\begin{center}
\begin{tikzpicture}[scale=0.85,>=Stealth]
% Points
\coordinate (C) at (0,0);
\coordinate (M) at (1.2,0.9);
\coordinate (Mp) at (2.4,1.8);   % H(C,2)(M)
\coordinate (Mpp) at (4.4,2.8);  % translation (2,1)
\coordinate (Om) at (-2,-1);     % centre final = -u
% Vecteur u
\draw[->,popgreen,thick] (Mp) -- (Mpp) node[midway,above,sloped]{$\vec{u}$};
\draw[->,popgreen,thick,dashed] (C) -- (2,1) node[midway,below]{$\vec{u}$};
% Homothétie depuis C
\draw[->,popblue,thick] (C) -- (M);
\draw[->,popblue,thick] (C) -- (Mp);
\draw[popblue,dashed] (C) -- (Mp);
% Homothétie depuis Omega
\draw[poporange,dashed,thick] (Om) -- (M);
\draw[poporange,dashed,thick] (Om) -- (Mpp);
\draw[->,poporange,thick] (Om) -- ($(Om)!0.5!(M)$);
% Points et labels
\fill[popdark] (C) circle (2pt) node[below left]{$C$};
\fill[popblue] (M) circle (2pt) node[above right]{$M$};
\fill[popblue] (Mp) circle (2pt) node[above]{$M'=H(M)$};
\fill[popgreen] (Mpp) circle (2pt) node[above right]{$M''=t(M')$};
\fill[poporange] (Om) circle (2.5pt) node[below left]{$\Omega$};
\end{tikzpicture}
\end{center}
\legende{$H_{C,2}$ suivie de $t_{\vec{u}}$ : $M \mapsto M' \mapsto M''$. Le point $M''$ est aussi l'image de $M$ par l'homothétie de centre $\Omega$ (point fixe) et de rapport $2$ : les points $\Omega$, $M$, $M''$ sont alignés avec $\overrightarrow{\Omega M''} = 2\,\overrightarrow{\Omega M}$.}
\end{popfigure}

Le schéma comparatif ci-dessous montre que les deux ordres de composition donnent des homothéties de même rapport mais de centres distincts.

\begin{popfigure}
\begin{center}
\begin{tikzpicture}[scale=0.8,>=Stealth]
\coordinate (C) at (0,0);
\coordinate (u) at (2,1);
% Centre de t o H : Omega1 = C + u/(1-k), k=2 -> C - u
\coordinate (O1) at (-2,-1);
% Centre de H o t : Omega2 = C + k u/(1-k) = C -2u
\coordinate (O2) at (-4,-2);
\draw[->,popmuted] (-5,0) -- (3,0);
\draw[->,popmuted] (0,-2.8) -- (0,2);
\draw[->,popgreen,thick] (C) -- (u) node[midway,above left]{$\vec{u}$};
\fill[popdark] (C) circle (2pt) node[above right]{$C$};
\fill[popblue] (O1) circle (2.5pt) node[below left]{$\Omega_1 = C-\vec{u}$};
\fill[poporange] (O2) circle (2.5pt) node[below left]{$\Omega_2 = C-2\vec{u}$};
\draw[popblue,->,thick,dashed] (C) -- (O1) node[midway,above]{$-\vec{u}$};
\draw[poporange,->,thick,dashed] (O1) -- (O2) node[midway,above left]{$-\vec{u}$};
\node[popblue,align=center] at (2.6,-1.6) {$t_{\vec{u}}\circ H_{C,2}$\\ centre $\Omega_1$};
\node[poporange,align=center] at (-4,-0.6) {$H_{C,2}\circ t_{\vec{u}}$\\ centre $\Omega_2$};
\end{tikzpicture}
\end{center}
\legende{Comparaison des deux ordres pour $k=2$ : $t_{\vec{u}}\circ H_{C,2}$ a pour centre $\Omega_1 = C + \frac{\vec{u}}{1-2} = C - \vec{u}$, tandis que $H_{C,2}\circ t_{\vec{u}}$ a pour centre $\Omega_2 = C - 2\vec{u}$. Même rapport $2$, centres distincts.}
\end{popfigure}

\begin{attention}[Erreur fréquente n°1 : le centre fantôme]
L'erreur la plus répandue au Probatoire est d'affirmer que la composée $t_{\vec{u}}\circ H_{C,k}$ est \og l'homothétie de centre $C$ et de rapport $k$ \fg{}. C'est \textbf{faux} dès que $\vec{u} \neq \vec{0}$ : le centre $C$ n'est pas fixe par la composée (il est envoyé sur $C+\vec{u}$) ! Une homothétie de rapport $k\neq 1$ possède un \emph{unique} point fixe : son centre. Il faut donc \textbf{impérativement résoudre l'équation du point fixe} $\Omega = k\Omega + \vec{V}$. Toute réponse donnant le centre sans ce calcul est suspecte.
\end{attention}

\begin{attention}[Autres pièges classiques]
\begin{itemize}
\item \textbf{Oublier l'ordre} : $H\circ t$ signifie \og $t$ d'abord, $H$ ensuite \fg{} (on applique d'abord la transformation écrite à droite). Intervertir l'ordre change le vecteur $\vec{V}$ ($\vec{u}$ devient $k\vec{u}$) et donc le centre.
\item \textbf{Se tromper de signe} dans $\Omega = \dfrac{\vec{V}}{1-k}$ : avec $k = 2$, le dénominateur vaut $-1$, d'où $\Omega = -\vec{V}$. Une erreur de signe ici fausse tout l'exercice.
\item \textbf{Conclure \og translation \fg{} quand $k\neq 1$} : la présence d'un terme additif $\vec{V}$ ne fait pas de la composée une translation ; c'est le coefficient de $M$ qui décide.
\end{itemize}
\end{attention}

\begin{savaistu}[Zoom et panoramique sur ton smartphone]
Quand tu zoomes sur une photo (homothétie de rapport $k>1$ centrée sur l'écran) puis que tu fais glisser l'image avec le doigt (translation), le résultat affiché est rigoureusement une \emph{nouvelle homothétie de même rapport}, centrée en un autre point. Les logiciels de traitement d'image, les cartes GPS (comme celles utilisées par les applications de transport à Douala ou Yaoundé) et les jeux vidéo calculent en permanence de telles composées : c'est exactement la formule $M' = kM + \vec{V}$ de cette leçon qui tourne dans ton téléphone !
\end{savaistu}

\begin{exercice}[À toi de jouer]
Dans un repère orthonormé, on donne $C(1\,;2)$, $H = H_{C,3}$ et $\vec{u}(-2\,;4)$.
\begin{enumerate}
\item Déterminer l'expression analytique de $f = t_{\vec{u}}\circ H$.
\item En déduire la nature de $f$ et les coordonnées de son centre $\Omega$.
\item Mêmes questions pour $g = H\circ t_{\vec{u}}$. Les centres sont-ils les mêmes ?
\end{enumerate}
\end{exercice}

\begin{corrige}[Exercice 1]
\textbf{1.} $H(M) = 3M + (1-3)C = 3M - 2C = \bigl(3x-2\,;\,3y-4\bigr)$. Puis $t_{\vec{u}}$ ajoute $(-2\,;4)$ :
\[ f(M) = \bigl(3x-4\,;\,3y\bigr), \quad\text{soit}\quad M' = 3M + \vec{V}\ \text{avec}\ \vec{V}=(-4\,;0). \]
\textbf{2.} $k = 3 \neq 1$ : $f$ est une homothétie de rapport $3$. Point fixe : $\Omega = \dfrac{\vec{V}}{1-3} = \dfrac{(-4\,;0)}{-2} = (2\,;0)$. Donc $f = H_{\Omega,\,3}$ avec $\Omega(2\,;0)$.
\textbf{3.} $t_{\vec{u}}(M) = (x-2\,;\,y+4)$, puis $H$ : $g(M) = 3(M+\vec{u}) - 2C = \bigl(3x-6-2\,;\,3y+12-4\bigr) = \bigl(3x-8\,;\,3y+8\bigr)$. C'est une homothétie de rapport $3$ ; point fixe : $\Omega' = \dfrac{(-8\,;8)}{-2} = (4\,;-4)$. Les centres $\Omega(2\,;0)$ et $\Omega'(4\,;-4)$ sont \textbf{différents} : la composition n'est pas commutative.
\end{corrige}

\coursec{Composée d'une homothétie et d'une rotation de même centre (Similitude directe)}

Dans la section précédente, nous avons composé une homothétie avec une translation : le centre disparaissait, et la composée était soit une homothétie (de centre déplacé), soit une translation. Nous allons maintenant composer une homothétie avec une rotation \emph{de même centre}. Cette fois, le centre reste fixe, et la transformation obtenue combine deux gestes familiers : \textbf{tourner} et \textbf{agrandir} (ou réduire). Le résultat porte un nom célèbre : la \cle{similitude directe}. C'est elle qui gouverne les spirales des coquillages, des ouragans vus par satellite et même des galaxies.

\courssub{Rappels : rotation et homothétie de même centre}

On se place dans le plan muni d'un repère orthonormé direct $(O\,;\vec{i},\vec{j})$. Rappelons les deux définitions analytiques rencontrées dans les sections 3 et 4.

\begin{itemize}
\item La \cle{rotation} $R_{O,\alpha}$ de centre $O$ et d'angle $\alpha$ associe à tout point $M(x\,;y)$ le point $M'(x'\,;y')$ avec :
\[
\begin{pmatrix} x' \\ y' \end{pmatrix} = R_\alpha \begin{pmatrix} x \\ y \end{pmatrix}, \qquad R_\alpha = \begin{pmatrix} \cos\alpha & -\sin\alpha \\ \sin\alpha & \cos\alpha \end{pmatrix}.
\]
\item L'\cle{homothétie} $H_{O,k}$ de centre $O$ et de rapport $k \neq 0$ associe à $M(x\,;y)$ le point $M'(x'\,;y')$ avec :
\[
\begin{pmatrix} x' \\ y' \end{pmatrix} = k\,I \begin{pmatrix} x \\ y \end{pmatrix} = \begin{pmatrix} kx \\ ky \end{pmatrix}, \qquad I = \begin{pmatrix} 1 & 0 \\ 0 & 1 \end{pmatrix}.
\]
\end{itemize}

Les deux transformations ont le \textbf{même centre} $O$, qui est leur unique point fixe (sauf cas triviaux $\alpha = 0$ ou $k = 1$). Question naturelle : que vaut $H_{O,k} \circ R_{O,\alpha}$ ? Et l'ordre a-t-il de l'importance ?

\courssub{Commutativité et définition de la similitude directe}

\begin{experience}[Activité de découverte : tourner puis agrandir, ou agrandir puis tourner ?]
Prenons $k = 2$, $\alpha = \dfrac{\pi}{2}$ et le point $M(1\,;0)$.
\begin{enumerate}
\item \textbf{Rotation d'abord :} $R_{O,\pi/2}(M) = M_1(0\,;1)$, puis $H_{O,2}(M_1) = M'(0\,;2)$.
\item \textbf{Homothétie d'abord :} $H_{O,2}(M) = M_2(2\,;0)$, puis $R_{O,\pi/2}(M_2) = M''(0\,;2)$.
\end{enumerate}
On obtient $M' = M''$ ! Vérifions que ce n'est pas un hasard, avec les matrices :
\[
(kI)\,R_\alpha = k\,(I R_\alpha) = k R_\alpha \quad\text{et}\quad R_\alpha\,(kI) = k\,(R_\alpha I) = k R_\alpha.
\]
La matrice $kI$ commute avec toutes les matrices : les deux composées sont \textbf{égales}.
\end{experience}

\begin{propriete}[Commutativité — admise et justifiée ci-dessus]
Pour tout réel $k \neq 0$ et tout angle $\alpha$ :
\[
H_{O,k} \circ R_{O,\alpha} = R_{O,\alpha} \circ H_{O,k}.
\]
Cette application composée se note $S_{O,k,\alpha}$ et s'appelle la \cle{similitude directe} de centre $O$, de rapport $|k|$ et d'angle $\alpha$ (si $k>0$). Sa matrice dans le repère $(O\,;\vec{i},\vec{j})$ est :
\[
k\,R_\alpha = \begin{pmatrix} k\cos\alpha & -k\sin\alpha \\ k\sin\alpha & k\cos\alpha \end{pmatrix}.
\]
\end{propriete}

\begin{definition}[Similitude directe de centre $O$]
Soit $O$ un point du plan, $k$ un réel non nul et $\alpha$ un réel. On appelle \cle{similitude directe} de centre $O$, de rapport $|k|$ et d'angle $\alpha$ (modulo $2\pi$, à $\pi$ près si $k<0$) l'application $S$ du plan dans lui-même qui, à tout point $M$, associe le point $M' = S(M)$ défini par :
\[
\overrightarrow{OM'} = k\,R_\alpha\big(\overrightarrow{OM}\big),
\]
c'est-à-dire :
\[
OM' = |k|\,OM \qquad\text{et}\qquad \big(\overrightarrow{OM},\overrightarrow{OM'}\big) = \alpha \ [2\pi] \ \text{si } k>0.
\]
Si $k > 0$, l'orientation est conservée ; si $k < 0$, on peut écrire $S = H_{O,|k|} \circ R_{O,\alpha+\pi}$ : l'angle devient $\alpha + \pi$.
\end{definition}

\begin{attention}[Rapport $k$ négatif : un piège classique]
Dans la définition analytique $X' = kR_\alpha X$, le coefficient $k$ peut être négatif. Dans ce cas, la similitude a pour \textbf{rapport} $|k|$ (un rapport est toujours positif !) et pour \textbf{angle} $\alpha + \pi$. Par exemple, $X' = -2 R_{\pi/3} X$ est la similitude de rapport $2$ et d'angle $\dfrac{\pi}{3} + \pi = \dfrac{4\pi}{3}$. Ne jamais écrire « rapport $-2$ » dans une copie de Probatoire.
\end{attention}

\begin{popfigure}
\begin{center}
\begin{tikzpicture}[scale=1.6,>=Stealth]
\coordinate (O) at (0,0);
\coordinate (M) at (2.6,0);
\coordinate (M1) at (0,2.6);
\coordinate (Mp) at (0,5.2);
% cercles
\draw[popmuted,dashed] (O) circle (2.6);
\draw[popmuted,dashed] (O) circle (5.2);
% points
\fill[popblue] (O) circle (1.5pt) node[below left]{$O$};
\fill[popink] (M) circle (1.5pt) node[below right]{$M$};
\fill[poporange] (M1) circle (1.5pt) node[right]{$M_1 = R_{O,\alpha}(M)$};
\fill[poppurple] (Mp) circle (1.5pt) node[right]{$M' = S(M)$};
% vecteurs
\draw[->,popblue,thick] (O) -- (M);
\draw[->,poporange,thick] (O) -- (M1);
\draw[->,poppurple,thick] (O) -- (Mp);
% arc angle
\draw[->,popgreen,thick] (1.1,0) arc (0:90:1.1) node[midway,above right=2pt]{$\alpha$};
% annotations
\node[popmuted] at (2.2,-0.9) {\small $OM$};
\node[popmuted] at (-1.6,4.6) {\small $OM' = |k|\,OM$};
\end{tikzpicture}
\end{center}
\legende{La similitude directe $S = H_{O,k}\circ R_{O,\alpha}$ : on tourne $M$ de l'angle $\alpha$ (point $M_1$), puis on multiplie la distance par $|k|$ (point $M'$). L'angle $(\overrightarrow{OM},\overrightarrow{OM'})$ vaut $\alpha$ et $OM' = |k|\,OM$.}
\end{popfigure}

\courssub{Théorème fondamental : la composée est une similitude directe}

\begin{propriete}[Théorème — démonstration exigible]
Soit $O$ un point du plan, $k \in \mathbb{R}^*$ et $\alpha \in \mathbb{R}$. La composée, dans n'importe quel ordre, de l'homothétie $H_{O,k}$ et de la rotation $R_{O,\alpha}$ est la similitude directe de centre $O$, de rapport $|k|$ et d'angle $\alpha$ (si $k>0$) :
\[
H_{O,k} \circ R_{O,\alpha} = R_{O,\alpha} \circ H_{O,k} = S_{O,k,\alpha}.
\]
Réciproquement, toute similitude directe de centre $O$ se décompose de manière \textbf{unique} en une homothétie et une rotation de centre $O$.
\end{propriete}

\begin{experience}[Démonstration guidée]
Soit $M(x\,;y)$ un point quelconque, de vecteur position $X = \begin{pmatrix} x \\ y \end{pmatrix}$.
\begin{enumerate}
\item \textbf{Image par la composée.} Par définition de la composée et par les définitions analytiques des sections 2 et 3 :
\[
X' = (kI)\big(R_\alpha X\big) = k\,R_\alpha X.
\]
\item \textbf{Effet sur les distances.} La rotation conserve les normes : $\|R_\alpha X\| = \|X\|$. Donc :
\[
OM' = \|X'\| = |k|\,\|R_\alpha X\| = |k|\,\|X\| = |k|\,OM.
\]
\item \textbf{Effet sur les angles.} Pour $k > 0$, $X'$ est l'image de $R_\alpha X$ par une homothétie de rapport positif, qui conserve les directions : $(\overrightarrow{OM},\overrightarrow{OM'}) = (\overrightarrow{OM}, R_\alpha\overrightarrow{OM}) = \alpha \ [2\pi]$.
\item \textbf{Conclusion.} $M'$ vérifie exactement les deux conditions de la définition : la composée est bien la similitude $S_{O,k,\alpha}$. L'unicité de la décomposition vient de ce que $k$ est imposé par le rapport des distances et $\alpha$ par l'angle. $\square$
\end{enumerate}
\end{experience}

\begin{exemplebox}[Exemple chiffré : $H_{O,2} \circ R_{O,\pi/2}$]
La matrice de la composée est :
\[
2\,R_{\pi/2} = 2\begin{pmatrix} \cos\frac{\pi}{2} & -\sin\frac{\pi}{2} \\ \sin\frac{\pi}{2} & \cos\frac{\pi}{2} \end{pmatrix} = 2\begin{pmatrix} 0 & -1 \\ 1 & 0 \end{pmatrix} = \begin{pmatrix} 0 & -2 \\ 2 & 0 \end{pmatrix}.
\]
C'est la similitude directe de centre $O$, de rapport $k = 2$ et d'angle $\alpha = \dfrac{\pi}{2}$. Vérification sur le point $M(1\,;0)$ :
\[
X' = \begin{pmatrix} 0 & -2 \\ 2 & 0 \end{pmatrix}\begin{pmatrix} 1 \\ 0 \end{pmatrix} = \begin{pmatrix} 0 \\ 2 \end{pmatrix}.
\]
On a bien $OM' = 2 = 2\,OM$ et $(\overrightarrow{OM},\overrightarrow{OM'}) = \dfrac{\pi}{2}$ : tout est cohérent.
\end{exemplebox}

\courssub{Définition analytique et écriture complexe}

\begin{aretenir}[Définition analytique de la similitude directe de centre $O$]
Dans le repère orthonormé direct $(O\,;\vec{i},\vec{j})$, la similitude directe $S_{O,k,\alpha}$ associe à $M(x\,;y)$ le point $M'(x'\,;y')$ avec :
\[
\boxed{\;X' = k\,R_\alpha\, X\;} \qquad\text{c'est-à-dire}\qquad \begin{cases} x' = k(x\cos\alpha - y\sin\alpha) \\ y' = k(x\sin\alpha + y\cos\alpha) \end{cases}
\]
En identifiant le plan au plan complexe d'origine $O$, si $M$ a pour affixe $z$, alors $M'$ a pour affixe :
\[
z' = k\,e^{i\alpha}\, z \qquad (k>0).
\]
Multiplier par $e^{i\alpha}$, c'est tourner de $\alpha$ ; multiplier par $k$, c'est dilater de rapport $k$. L'écriture complexe résume la similitude en \textbf{une seule multiplication}.
\end{aretenir}

\begin{propriete}[Propriétés métriques de la similitude directe]
Soit $S$ une similitude directe de rapport $k > 0$.
\begin{enumerate}
\item $S$ \textbf{conserve les angles orientés} : pour tous points $A,B,C$ d'images $A',B',C'$, $(\overrightarrow{A'B'},\overrightarrow{A'C'}) = (\overrightarrow{AB},\overrightarrow{AC}) \ [2\pi]$.
\item $S$ \textbf{multiplie les longueurs par $k$} : $A'B' = k\,AB$.
\item $S$ \textbf{multiplie les aires par $k^2$}.
\item $S$ transforme une droite en une droite, un cercle en un cercle, et conserve le parallélisme, l'orthogonalité et les contacts.
\end{enumerate}
\end{propriete}

\begin{propriete}[Réciproque d'une similitude directe]
La similitude directe $S_{O,k,\alpha}$ est bijective, et sa réciproque est la similitude directe de même centre, de rapport $\dfrac{1}{|k|}$ et d'angle $-\alpha$ :
\[
S_{O,k,\alpha}^{-1} = S_{O,\,1/k,\,-\alpha} = R_{O,-\alpha} \circ H_{O,\,1/k}.
\]
En effet, matriciellement : $(kR_\alpha)^{-1} = \dfrac{1}{k}R_\alpha^{-1} = \dfrac{1}{k}R_{-\alpha}$.
\end{propriete}

\courssub{Caractériser analytiquement une similitude directe}

\begin{methode}[Reconnaître une similitude directe à partir de sa matrice]
On donne une application $M(x\,;y) \mapsto M'(x'\,;y')$ définie par $X' = AX$ avec
\[
A = \begin{pmatrix} a & -b \\ b & a \end{pmatrix}, \qquad (a,b) \neq (0,0).
\]
\begin{enumerate}
\item \textbf{Forme.} Vérifier que la matrice a bien la forme $\begin{pmatrix} a & -b \\ b & a \end{pmatrix}$ : mêmes coefficients sur la diagonale, coefficients opposés sur l'antidiagonale. C'est la signature d'une similitude directe de centre $O$.
\item \textbf{Rapport.} Calculer $k = \sqrt{a^2 + b^2}$.
\item \textbf{Angle.} Déterminer $\alpha$ par $\cos\alpha = \dfrac{a}{k}$ et $\sin\alpha = \dfrac{b}{k}$ (placer le point $(a\,;b)$ sur le cercle trigonométrique).
\item \textbf{Centre.} Le centre est l'unique point fixe ; ici $X' = X \iff (A - I)X = 0$, qui donne $X = 0$ car $\det(A - I) = (a-1)^2 + b^2 \neq 0$ dès que $(a,b) \neq (1,0)$. Le centre est donc $O$.
\end{enumerate}
\end{methode}

\begin{exoresolu}[Caractérisation d'une similitude]
Soit $f$ l'application du plan qui à $M(x\,;y)$ associe $M'(x'\,;y')$ avec
\[
\begin{cases} x' = -x - y\sqrt{3} \\ y' = x\sqrt{3} - y \end{cases}
\]
Montrer que $f$ est une similitude directe et préciser ses éléments caractéristiques.
\tcblower
\textbf{Solution.} La matrice de $f$ est $A = \begin{pmatrix} -1 & -\sqrt{3} \\ \sqrt{3} & -1 \end{pmatrix}$, de la forme $\begin{pmatrix} a & -b \\ b & a \end{pmatrix}$ avec $a = -1$ et $b = \sqrt{3}$.

\textbf{Rapport :} $k = \sqrt{a^2 + b^2} = \sqrt{1 + 3} = 2$.

\textbf{Angle :} $\cos\alpha = \dfrac{a}{k} = -\dfrac{1}{2}$ et $\sin\alpha = \dfrac{b}{k} = \dfrac{\sqrt{3}}{2}$, donc $\alpha = \dfrac{2\pi}{3} \ [2\pi]$.

\textbf{Centre :} $(a,b) = (-1,\sqrt{3}) \neq (1,0)$, donc l'unique point fixe est $O$.

\textbf{Conclusion :} $f$ est la similitude directe de centre $O$, de rapport $2$ et d'angle $\dfrac{2\pi}{3}$ : $f = H_{O,2} \circ R_{O,\,2\pi/3}$.
\end{exoresolu}

\courssub{Itérations et spirales logarithmiques}

Que se passe-t-il si l'on applique la \emph{même} similitude $S_{O,k,\alpha}$ encore et encore à un point $M_0$ ? Posons $M_{n+1} = S(M_n)$. En complexe, $z_n = (ke^{i\alpha})^n z_0$, donc :
\[
OM_n = k^n\,OM_0 \qquad\text{et}\qquad \big(\overrightarrow{OM_0},\overrightarrow{OM_n}\big) = n\alpha \ [2\pi].
\]
Chaque étape tourne de $\alpha$ et multiplie la distance par $k$ : les points $M_0, M_1, M_2, \dots$ s'enroulent autour de $O$ en s'en éloignant (si $k>1$) ou en s'en rapprochant (si $0<k<1$). Le lieu de ces points est une \cle{spirale logarithmique}.

\begin{popfigure}
\begin{center}
\begin{tikzpicture}[scale=1.15,>=Stealth]
\draw[->,popmuted] (-3.4,0) -- (3.4,0);
\draw[->,popmuted] (0,-3.2) -- (0,3.4);
\fill[popblue] (0,0) circle (1.2pt) node[below left=1pt]{$O$};
% spirale : r = 1.25^(t/60°), angle 60° par étape
\draw[popink,thick,domain=0:540,samples=300,variable=\t]
  plot ({0.28*1.25^(\t/60)*cos(\t)},{0.28*1.25^(\t/60)*sin(\t)});
% points itérés M_n = S^n(M0), k=1.25, alpha=60°
\foreach \n in {0,1,2,3,4,5,6,7,8}{
  \fill[poppurple] ({0.28*1.25^\n*cos(60*\n)},{0.28*1.25^\n*sin(60*\n)}) circle (1.3pt);
}
\node[poppurple] at ({0.28*cos(0)},{0.28*sin(0)}) [below right=-1pt]{\small $M_0$};
\node[poppurple] at ({0.28*1.25*cos(60)},{0.28*1.25*sin(60)}) [right=-1pt]{\small $M_1$};
\node[poppurple] at ({0.28*1.25^2*cos(120)},{0.28*1.25^2*sin(120)}) [above left=-2pt]{\small $M_2$};
\node[poppurple] at ({0.28*1.25^3*cos(180)},{0.28*1.25^3*sin(180)}) [above=-1pt]{\small $M_3$};
\draw[->,poporange] (1.15,0) arc (0:60:1.15) node[midway,right=1pt]{\small $\alpha$};
\end{tikzpicture}
\end{center}
\legende{Itérées d'une similitude directe ($k = 1.25$, $\alpha = 60^\circ$) : les points $M_n = S^n(M_0)$ s'enroulent sur une spirale logarithmique. Chaque pas tourne de $\alpha$ et multiplie la distance à $O$ par $k$.}
\end{popfigure}

\begin{savaistu}[Les spirales de la nature]
La spirale logarithmique fut surnommée \emph{spira mirabilis} (« spirale merveilleuse ») par Jacques Bernoulli, qui la voulut gravée sur sa tombe. On la retrouve partout : la coquille du nautile qui grandit en conservant sa forme, les bras des galaxies spirales, l'œil d'un cyclone photographié par satellite, le vol d'un insecte vers une lampe. Le point commun de tous ces phénomènes : une croissance qui conserve les angles, c'est-à-dire une \textbf{similitude directe itérée}. En Terminale C, l'écriture complexe $z' = az$ avec $a = ke^{i\alpha}$ te donnera un outil d'une élégance redoutable pour étudier toutes les similitudes.
\end{savaistu}

\courssub{Applications à la géométrie : triangles semblables et configurations}

Les similitudes sont l'outil moderne pour traiter les \textbf{triangles semblables} : dire que deux triangles $ABC$ et $A'B'C'$ sont directement semblables, c'est dire qu'il existe une similitude directe envoyant $A$ sur $A'$, $B$ sur $B'$, $C$ sur $C'$. Toutes les propriétés de conservation (angles) et de proportionnalité (longueurs multipliées par $k$, aires par $k^2$) s'en déduisent immédiatement.

\begin{exoresolu}[Cercles et similitude]
Soit $\mathcal{C}$ le cercle de centre $\Omega(1\,;0)$ et de rayon $1$, et $S$ la similitude directe de centre $O$, de rapport $3$ et d'angle $\dfrac{\pi}{2}$. Déterminer la nature et les éléments caractéristiques de l'image $\mathcal{C}' = S(\mathcal{C})$, puis comparer les aires des disques.
\tcblower
\textbf{Solution.} Une similitude transforme un cercle en un cercle, de rayon multiplié par le rapport.

\textbf{Image du centre :} $S(\Omega) = \Omega'$ avec $X_{\Omega'} = 3R_{\pi/2}\begin{pmatrix} 1 \\ 0 \end{pmatrix} = 3\begin{pmatrix} 0 \\ 1 \end{pmatrix} = \begin{pmatrix} 0 \\ 3 \end{pmatrix}$, donc $\Omega'(0\,;3)$.

\textbf{Rayon :} $r' = 3 \times 1 = 3$.

\textbf{Conclusion :} $\mathcal{C}'$ est le cercle de centre $\Omega'(0\,;3)$ et de rayon $3$. Les aires vérifient $\mathcal{A}' = k^2\,\mathcal{A} = 9\,\mathcal{A} = 9\pi$ (unités d'aire), ce que confirme le calcul direct $\pi \times 3^2 = 9\pi$.
\end{exoresolu}

\begin{attention}[Directe ou indirecte : ne pas confondre]
Une similitude \textbf{directe} conserve l'orientation (les angles orientés) : sa matrice est $\begin{pmatrix} a & -b \\ b & a \end{pmatrix}$, de déterminant $a^2+b^2 > 0$. Une similitude \textbf{indirecte} (composée d'une similitude directe et d'une symétrie axiale) inverse l'orientation : sa matrice est $\begin{pmatrix} a & b \\ b & -a \end{pmatrix}$, de déterminant $-(a^2+b^2) < 0$. Avant de conclure, vérifie toujours la position des signes dans la matrice. Autre piège : oublier que le rapport d'une similitude est \emph{toujours positif} — un $k$ négatif dans l'écriture analytique se réabsorbe dans l'angle ($\alpha$ devient $\alpha + \pi$).
\end{attention}

\begin{exercice}[Pour s'entraîner]
\begin{enumerate}
\item Soit $S$ l'application définie par $\begin{cases} x' = x + y \\ y' = -x + y \end{cases}$. Montrer que $S$ est une similitude directe et déterminer son centre, son rapport et son angle.
\item On pose $T = S \circ S$ où $S$ est la similitude de la question 1. Déterminer la nature et les éléments caractéristiques de $T$.
\item Soit $A(1\,;0)$. Construire les points $A_1 = S(A)$, $A_2 = S(A_1)$, $A_3 = S(A_2)$ et vérifier que $OA_3 = 2\sqrt{2}$.
\end{enumerate}
\end{exercice}

\begin{corrige}[Exercice]
\textbf{1.} La matrice est $A = \begin{pmatrix} 1 & 1 \\ -1 & 1 \end{pmatrix} = \begin{pmatrix} a & -b \\ b & a \end{pmatrix}$ avec $a = 1$, $b = -1$. Rapport : $k = \sqrt{1+1} = \sqrt{2}$. Angle : $\cos\alpha = \dfrac{1}{\sqrt{2}}$, $\sin\alpha = -\dfrac{1}{\sqrt{2}}$, donc $\alpha = -\dfrac{\pi}{4}$. Centre : $O$ (seul point fixe car $(a-1)^2 + b^2 = 1 \neq 0$). Donc $S = S_{O,\sqrt{2},\,-\pi/4}$.

\textbf{2.} Composer deux similitudes de même centre revient à multiplier les rapports et additionner les angles : $T$ est la similitude de centre $O$, de rapport $\sqrt{2}\times\sqrt{2} = 2$ et d'angle $-\dfrac{\pi}{4}-\dfrac{\pi}{4} = -\dfrac{\pi}{2}$.

\textbf{3.} $A_1 = (1\,;-1)$, $A_2 = S(A_1) = (1+(-1)\,;\,-1+(-1)) = (0\,;-2)$, $A_3 = S(A_2) = (-2\,;-2)$. Alors $OA_3 = \sqrt{4+4} = 2\sqrt{2}$, ce qui confirme $OA_3 = k^3\,OA = (\sqrt{2})^3 = 2\sqrt{2}$.
\end{corrige}

\begin{aretenir}[L'essentiel de la section]
\begin{itemize}
\item La composée, dans n'importe quel ordre, d'une homothétie et d'une rotation \textbf{de même centre} $O$ est la \cle{similitude directe} $S_{O,k,\alpha} = H_{O,k} \circ R_{O,\alpha} = R_{O,\alpha} \circ H_{O,k}$, de matrice $kR_\alpha$.
\item Définition analytique : $X' = kR_\alpha X$ ; écriture complexe : $z' = ke^{i\alpha}z$.
\item Éléments caractéristiques : centre $O$ (unique point fixe), rapport $|k|$, angle $\alpha$ (ou $\alpha + \pi$ si $k<0$).
\item Elle conserve les angles orientés, multiplie les longueurs par $|k|$ et les aires par $k^2$.
\item Sa réciproque est $S_{O,\,1/k,\,-\alpha}$.
\item Les itérées d'une similitude directe ($k \neq 1$) décrivent une spirale logarithmique.
\end{itemize}
\end{aretenir}

\newpage

\coursec{Composées de symétries orthogonales, de rotations et d'homothéties de centres distincts}

Dans la leçon commune, tu as étudié les composées de transformations \emph{de même centre} : rotations, homothéties, similitudes. Mais le plan regorge de transformations dont les éléments caractéristiques sont \emph{distincts} : deux miroirs parallèles d'un couloir, deux homothéties centrées en deux villes différentes... Ce complément, réservé à certaines séries, te donne les outils pour traiter ces situations. L'idée directrice est magnifique : \textbf{toutes les transformations usuelles se fabriquent à partir des symétries orthogonales}.

\courssub{Composée de deux symétries orthogonales}

\begin{experience}[Découverte : deux miroirs]
Place un objet entre deux miroirs parallèles : tu observes une infinité d'images qui semblent se déplacer régulièrement. Avec deux miroirs inclinés formant un dièdre, les images semblent tourner autour de l'arête. Formalisons ces observations.
\end{experience}

\begin{propriete}[Axes parallèles]
Soient $(\Delta_1)$ et $(\Delta_2)$ deux droites strictement parallèles, et $S_1$, $S_2$ les symétries orthogonales d'axes respectifs $(\Delta_1)$ et $(\Delta_2)$. Alors
\[ S_2 \circ S_1 = t_{\vec{u}} \quad \text{où} \quad \vec{u} = 2\,\overrightarrow{H_1H_2}, \]
$H_1$ étant un point de $(\Delta_1)$ et $H_2$ son projeté orthogonal sur $(\Delta_2)$. Le vecteur $\vec{u}$ est \cle{normal} aux deux axes, de norme égale au double de la distance entre eux.
\end{propriete}

\begin{propriete}[Axes sécants]
Soient $(\Delta_1)$ et $(\Delta_2)$ deux droites sécantes en $\Omega$, d'angle orienté $(\Delta_1,\Delta_2) = \alpha$. Alors
\[ S_2 \circ S_1 = r_{\Omega,\,2\alpha}, \]
la rotation de centre $\Omega$ et d'angle $2\alpha$.
\end{propriete}

\begin{attention}[L'ordre compte !]
La composée n'est pas commutative : $S_2 \circ S_1 \neq S_1 \circ S_2$ en général. Pour des axes sécants, $S_2 \circ S_1$ est la rotation d'angle $2(\Delta_1,\Delta_2)$, tandis que $S_1 \circ S_2$ donne l'angle \emph{opposé}. Retiens : \textbf{on double l'angle du premier axe vers le second}.
\end{attention}

\begin{popfigure}
\begin{tikzpicture}[scale=1.1]
  \draw[popblue,thick] (0,-0.3) -- (0,2.3) node[above]{$(\Delta_1)$};
  \draw[poporange,thick] (1.2,-0.3) -- (1.2,2.3) node[above]{$(\Delta_2)$};
  \fill[popink] (-0.9,0.7) circle(1.5pt) node[left]{$M$};
  \fill[popgreen] (0.9,0.7) circle(1.5pt) node[right]{$M_1$};
  \fill[poppurple] (1.5,0.7) circle(1.5pt) node[right]{$M'$};
  \draw[->,popdark,thick] (-0.9,0.35) -- (1.5,0.35) node[midway,below]{$\vec{u}=2\overrightarrow{H_1H_2}$};
  \draw[<->,popmuted] (0,1.9) -- (1.2,1.9) node[midway,above]{$d$};
\end{tikzpicture}
\legende{Composée de deux symétries d'axes parallèles : $M \xrightarrow{S_1} M_1 \xrightarrow{S_2} M'$, et $\overrightarrow{MM'} = \vec{u}$.}
\end{popfigure}

\begin{methode}[Décomposer une translation ou une rotation]
\begin{enumerate}
  \item \textbf{Translation} $t_{\vec{u}}$ : choisir deux droites $(\Delta_1)$ et $(\Delta_2)$ parallèles, perpendiculaires à $\vec{u}$, distantes de $\dfrac{\|\vec{u}\|}{2}$, orientées dans le sens de $\vec{u}$. Alors $t_{\vec{u}} = S_{\Delta_2} \circ S_{\Delta_1}$. L'un des deux axes peut être choisi \emph{arbitrairement}.
  \item \textbf{Rotation} $r_{\Omega,\theta}$ : choisir deux droites sécantes en $\Omega$ telles que $(\Delta_1,\Delta_2) = \dfrac{\theta}{2}$. Alors $r_{\Omega,\theta} = S_{\Delta_2} \circ S_{\Delta_1}$. Là encore, un axe est libre.
\end{enumerate}
\end{methode}

\courssub{Composée d'une translation et d'une symétrie orthogonale}

\begin{propriete}[Symétrie glissée]
Soit $t_{\vec{u}}$ une translation et $S_{\Delta}$ une symétrie orthogonale dont l'axe $(\Delta)$ est \cle{normal} à $\vec{u}$. Alors $S_{\Delta} \circ t_{\vec{u}}$ (et $t_{\vec{u}} \circ S_{\Delta}$) est une \textbf{symétrie orthogonale} d'axe $(\Delta')$, image de $(\Delta)$ par la translation de vecteur $\dfrac{1}{2}\vec{u}$ (avec un signe selon l'ordre de composition).
\end{propriete}

En effet, en écrivant $t_{\vec{u}} = S_{\Delta_2} \circ S_{\Delta_1}$ avec $(\Delta_1) = (\Delta)$, on obtient $S_{\Delta} \circ t_{\vec{u}} = S_{\Delta} \circ S_{\Delta_2} \circ S_{\Delta} = S_{\Delta'}$ où $(\Delta')$ est le translaté de $(\Delta)$. Ce résultat est la clé de nombreuses démonstrations sur les frises et pavages.

\courssub{Composée de deux rotations de centres distincts}

\begin{propriete}[Théorème fondamental]
Soient $r_A$ et $r_B$ deux rotations de centres distincts $A$ et $B$, d'angles respectifs $\alpha$ et $\beta$.
\begin{itemize}
  \item Si $\alpha + \beta \not\equiv 0 \ [2\pi]$, alors $r_B \circ r_A$ est une \textbf{rotation d'angle $\alpha + \beta$}, dont le centre $\Omega$ est déterminé par les conditions angulaires :
  \[ (\overrightarrow{\Omega A}, \overrightarrow{\Omega B}) \ \text{avec} \ (\overrightarrow{A\Omega},\overrightarrow{AB}) = \frac{\alpha}{2} \ \text{et} \ (\overrightarrow{BA},\overrightarrow{B\Omega}) = \frac{\beta}{2}. \]
  \item Si $\alpha + \beta \equiv 0 \ [2\pi]$, alors $r_B \circ r_A$ est une \textbf{translation}.
\end{itemize}
\end{propriete}

\begin{experience}[Démonstration guidée]
Décomposons chaque rotation en symétries, en choisissant astucieusement l'axe commun $(AB)$ :
\begin{enumerate}
  \item Écris $r_A = S_{(AB)} \circ S_{d_1}$ où $d_1$ passe par $A$ avec $(d_1,(AB)) = \dfrac{\alpha}{2}$.
  \item Écris $r_B = S_{d_2} \circ S_{(AB)}$ où $d_2$ passe par $B$ avec $((AB),d_2) = \dfrac{\beta}{2}$.
  \item Compose : $r_B \circ r_A = S_{d_2} \circ S_{(AB)} \circ S_{(AB)} \circ S_{d_1} = S_{d_2} \circ S_{d_1}$, car $S_{(AB)} \circ S_{(AB)} = \mathrm{Id}$.
  \item Conclus : si $d_1$ et $d_2$ se coupent en $\Omega$, c'est la rotation de centre $\Omega$ et d'angle $2(d_1,d_2) = \alpha + \beta$ ; si elles sont parallèles (cas $\alpha + \beta \equiv 0\ [2\pi]$), c'est une translation.
\end{enumerate}
\end{experience}

\begin{exoresolu}[Un triangle équiléteral caché]
Soient $A$ et $B$ deux points distincts, $r_A$ la rotation de centre $A$ et d'angle $\dfrac{\pi}{3}$, $r_B$ la rotation de centre $B$ et d'angle $\dfrac{\pi}{3}$. Déterminer la nature et les éléments caractéristiques de $f = r_B \circ r_A$.
\tcblower
\textbf{Solution.} La somme des angles est $\dfrac{\pi}{3} + \dfrac{\pi}{3} = \dfrac{2\pi}{3} \not\equiv 0 \ [2\pi]$ : $f$ est une \textbf{rotation d'angle $\dfrac{2\pi}{3}$}. Son centre $\Omega$ vérifie :
\[ (\overrightarrow{A\Omega},\overrightarrow{AB}) = \frac{\alpha}{2} = \frac{\pi}{6} \quad \text{et} \quad (\overrightarrow{BA},\overrightarrow{B\Omega}) = \frac{\beta}{2} = \frac{\pi}{6}. \]
Le triangle $A\Omega B$ est donc isocèle en $\Omega$ avec des angles à la base de $\dfrac{\pi}{6}$ ; l'angle en $\Omega$ vaut $\pi - \dfrac{\pi}{3} = \dfrac{2\pi}{3}$. On construit $\Omega$ comme intersection des deux demi-droites ainsi définies (du côté imposé par le sens direct).
\end{exoresolu}

\courssub{Composée de deux homothéties de centres distincts}

\begin{propriete}[Nature de la composée]
Soient $h_A$ et $h_B$ deux homothéties de centres distincts $A$ et $B$, de rapports respectifs $k$ et $k'$.
\begin{itemize}
  \item Si $kk' \neq 1$, alors $h_B \circ h_A$ est une \textbf{homothétie de rapport $kk'$}, dont le centre $\Omega$ est \cle{aligné} avec $A$ et $B$.
  \item Si $kk' = 1$, alors $h_B \circ h_A$ est une \textbf{translation} de vecteur $(1-k')\,\overrightarrow{AB}$ (à adapter selon l'ordre de composition).
\end{itemize}
\end{propriete}

\begin{methode}[Trouver le centre par l'analytique]
Dans un repère, avec $A(a\,;a')$, $B(b\,;b')$ et $M(x\,;y)$ d'image $M''(x''\,;y'')$ :
\begin{enumerate}
  \item Écris les expressions analytiques : $h_A : \begin{cases} x' - a = k(x-a) \\ y' - a' = k(y-a') \end{cases}$ puis compose avec celles de $h_B$.
  \item Tu obtiens $x'' = kk'\,x + c$, $y'' = kk'\,y + c'$ : si $kk' \neq 1$, c'est l'écriture analytique d'une homothétie de rapport $kk'$ ; le centre est le point fixe, solution de $x = kk'x + c$.
  \item Si $kk' = 1$, tu reconnais l'écriture analytique d'une translation de vecteur $(c\,;c')$.
\end{enumerate}
\end{methode}

\begin{exemplebox}[Exemple chiffré]
Dans un repère, $A(0\,;0)$, $B(3\,;0)$, $h_A$ de rapport $2$, $h_B$ de rapport $\dfrac{1}{2}$. Pour $M(x\,;y)$ : $h_A(M) = M'(2x\,;2y)$, puis $h_B(M') = M''$ avec $x'' - 3 = \dfrac{1}{2}(2x-3)$, soit $x'' = x + \dfrac{3}{2}$, et $y'' = y$. Comme $2 \times \dfrac{1}{2} = 1$, $h_B \circ h_A$ est la \textbf{translation} de vecteur $\dfrac{3}{2}\,\vec{\imath}$, c'est-à-dire $(1-k')\overrightarrow{AB}$ avec $k' = \tfrac12$ : $\tfrac12 \overrightarrow{AB}$. Cohérent !
\end{exemplebox}

\begin{attention}[Pièges fréquents]
\begin{itemize}
  \item Ne jamais conclure « homothétie » avant d'avoir vérifié que $kk' \neq 1$ : le cas $kk' = 1$ donne une translation.
  \item Le centre de la composée n'est \textbf{pas} le milieu de $[AB]$ : il se calcule (point fixe) ou se construit en cherchant l'image d'un point bien choisi.
  \item Pour les rotations, la condition porte sur $\alpha + \beta \equiv 0 \ [2\pi]$, pas sur $\alpha + \beta = 0$ : deux angles de $\pi$ donnent aussi une translation !
\end{itemize}
\end{attention}

\begin{savaistu}[Des frises de l'art camerounais aux cristaux]
Les motifs répétitifs des pagnes, des sculptures bamiléké ou des clôtures décoratives sont gouvernés par les composées de symétries et de translations : les mathématiciens appellent ces structures des « groupes de frises » — il en existe exactement \textbf{sept types}. Les mêmes idées permettent aux cristallographes de classifier les 230 structures possibles des cristaux. Composer des transformations, c'est décrire les symétries du monde !
\end{savaistu}

\begin{exercice}[Pour s'entraîner]
Soient $A$ et $B$ distincts, $r_A$ de centre $A$ et d'angle $\dfrac{\pi}{2}$, $r_B$ de centre $B$ et d'angle $-\dfrac{\pi}{2}$.
\begin{enumerate}
  \item Déterminer la nature de $r_B \circ r_A$.
  \item En décomposant via l'axe $(AB)$, exprimer le vecteur de cette translation en fonction de $\overrightarrow{AB}$.
\end{enumerate}
\end{exercice}

\begin{corrige}[Exercice 1]
\textbf{1.} La somme des angles est $\dfrac{\pi}{2} - \dfrac{\pi}{2} = 0 \equiv 0 \ [2\pi]$ : $r_B \circ r_A$ est une \textbf{translation}.

\textbf{2.} Écrivons $r_A = S_{(AB)} \circ S_{d_1}$ avec $(d_1,(AB)) = \dfrac{\pi}{4}$, et $r_B = S_{d_2} \circ S_{(AB)}$ avec $((AB),d_2) = -\dfrac{\pi}{4}$. Alors $r_B \circ r_A = S_{d_2} \circ S_{d_1}$. Les droites $d_1$ et $d_2$ sont parallèles (angles $\dfrac{\pi}{4}$ et $-\dfrac{\pi}{4} + \dfrac{\pi}{2} = \dfrac{\pi}{4}$ par rapport à une même direction : elles font entre elles un angle nul modulo $\pi$). La distance entre $d_1$ et $d_2$, projetées depuis $A$ et $B$, donne un vecteur de translation de norme $2 \times \dfrac{AB}{\sqrt{2}} = AB\sqrt{2}$, dirigé perpendiculairement à $d_1$ : on trouve que le vecteur est l'image de $\overrightarrow{AB}$ par la rotation d'angle $\dfrac{\pi}{2}$, c'est-à-dire $\vec{u}$ tel que $\|\vec{u}\| = AB$ et $(\overrightarrow{AB},\vec{u}) = \dfrac{\pi}{2}$.
\end{corrige}

\newpage

\coursec{Activité d'intégration}

\courssub{Objectifs de l'activité}

À l'issue de cette activité d'intégration, tu seras capable de :

\begin{itemize}
\item écrire la définition analytique d'une rotation, d'une homothétie et d'une translation dans un repère orthonormé ;
\item composer plusieurs transformations et déterminer l'expression analytique globale de la composée sous la forme $X' = MX + V$ ;
\item reconnaître la nature d'une transformation à partir de son écriture analytique (similitude directe, translation, homothétie) ;
\item déterminer les éléments caractéristiques d'une similitude directe : centre (point fixe), rapport et angle ;
\item résoudre un problème inverse : ajuster une transformation pour atteindre un objectif imposé ;
\item vérifier un calcul algébrique par une construction géométrique sur papier millimétré.
\end{itemize}

\courssub{Situation}

La \cle{CNPS} (Caisse Nationale de Prévoyance Sociale), dont le siège est à Yaoundé, souhaite moderniser son identité visuelle. Le cabinet de communication retenu propose un logo construit à partir d'un \cle{motif de base} très simple : un triangle rectangle isocèle, symbole de stabilité et d'équilibre. Ce motif est ensuite déplacé, agrandi et pivoté par une suite de transformations géométriques programmées dans un logiciel de dessin.

Le graphiste travaille dans un repère orthonormé $(O\,;\vec{\imath},\vec{\jmath})$ du plan, l'unité étant le centimètre. Le motif de base est le triangle $ABC$ avec :
\[ A(0\,;0), \qquad B(2\,;0), \qquad C(0\,;2). \]

Pour obtenir le motif final $A'B'C'$, le logiciel applique successivement au triangle $ABC$ les quatre transformations suivantes :

\begin{enumerate}
\item $R_1$ : rotation de centre $A$ et d'angle $+90^\circ$ ;
\item $H$ : homothétie de centre $A$ et de rapport $k = 2$ ;
\item $t$ : translation de vecteur $\vec{u} = \begin{pmatrix} 5 \\ 5 \end{pmatrix}$ ;
\item $R_4$ : rotation de centre $\Omega(5\,;5)$ et d'angle $-90^\circ$.
\end{enumerate}

On note $F = R_4 \circ t \circ H \circ R_1$ la transformation globale qui envoie le motif initial sur le motif final : pour tout point $M$, $F(M) = R_4\bigl(t(H(R_1(M)))\bigr)$.

Le directeur artistique pose deux questions au graphiste :
\begin{itemize}
\item \textbf{Question 1 :} quelle est la nature exacte de la transformation globale $F$ ? S'agit-il d'une similitude directe ? Si oui, quels sont son centre, son rapport et son angle ?
\item \textbf{Question 2 (problème inverse) :} pour des raisons de mise en page, on exige que l'image finale du point $C$ soit exactement le point $C'(10\,;10)$. Comment faut-il modifier le vecteur de la translation de l'étape 3 (les trois autres transformations restant inchangées) pour atteindre cet objectif ?
\end{itemize}

\courssub{Consignes}

Travail en groupes de trois ou quatre élèves. Matériel : papier millimétré, règle, compas, rapporteur, calculatrice.

\begin{enumerate}
\item \textbf{Écriture analytique de chaque étape.} Pour un point $M(x\,;y)$ d'image $M'(x'\,;y')$, écrire les définitions analytiques de $R_1$, de $H$, de $t$ et de $R_4$.
\item \textbf{Calcul des images étape par étape.} Déterminer les coordonnées des images successives des points $A$, $B$ et $C$ après chacune des quatre transformations. Présenter les résultats dans un tableau.
\item \textbf{Expression globale.} En composant les quatre expressions analytiques, montrer que $F$ s'écrit sous la forme
\[ \begin{pmatrix} x' \\ y' \end{pmatrix} = M\begin{pmatrix} x \\ y \end{pmatrix} + V, \]
où $M$ est une matrice $2\times 2$ et $V$ un vecteur colonne que l'on déterminera.
\item \textbf{Nature de $F$.} Montrer que la matrice $M$ est de la forme $\begin{pmatrix} a & -b \\ b & a \end{pmatrix}$. En déduire que $F$ est une similitude directe, et déterminer son rapport $k$ et son angle $\theta$.
\item \textbf{Centre de $F$.} Le centre d'une similitude directe de rapport $k \neq 1$ est son unique point fixe. Résoudre le système $F(\Omega) = \Omega$ pour déterminer les coordonnées du centre. Comparer avec le résultat de la question 2.
\item \textbf{Vérification graphique.} Sur papier millimétré, placer le triangle $ABC$, construire le triangle final $A'B'C'$ à la règle et au compas, et vérifier que les longueurs ont bien été multipliées par le rapport trouvé et que la figure a bien tourné de l'angle trouvé.
\item \textbf{Problème inverse.} On remplace le vecteur de translation $\vec{u} = \begin{pmatrix} 5 \\ 5 \end{pmatrix}$ par un vecteur inconnu $\vec{w} = \begin{pmatrix} p \\ q \end{pmatrix}$, les transformations $R_1$, $H$ et $R_4$ restant inchangées. Déterminer $p$ et $q$ pour que l'image finale de $C$ soit le point de coordonnées $(10\,;10)$.
\end{enumerate}

\courssub{Ressources à mobiliser}

\begin{aretenir}[Boîte à outils de la leçon]
\begin{itemize}
\item \textbf{Translation} de vecteur $\vec{u} = \begin{pmatrix} a \\ b \end{pmatrix}$ : $\begin{cases} x' = x + a \\ y' = y + b \end{cases}$.
\item \textbf{Homothétie} de centre $\Omega(x_0\,;y_0)$ et de rapport $k$ : $\begin{cases} x' = kx + (1-k)x_0 \\ y' = ky + (1-k)y_0 \end{cases}$.
\item \textbf{Rotation} de centre $\Omega(x_0\,;y_0)$ et d'angle $\theta$ :
\[ \begin{cases} x' = x_0 + (x - x_0)\cos\theta - (y - y_0)\sin\theta \\ y' = y_0 + (x - x_0)\sin\theta + (y - y_0)\cos\theta \end{cases}. \]
Pour $\theta = 90^\circ$ : $(x'\,;y') = \bigl(x_0 - (y-y_0)\,;\, y_0 + (x-x_0)\bigr)$. Pour $\theta = -90^\circ$ : $(x'\,;y') = \bigl(x_0 + (y-y_0)\,;\, y_0 - (x-x_0)\bigr)$.
\item \textbf{Composée} de deux homothéties de même centre, de rapports $k$ et $k'$ : homothétie de même centre et de rapport $kk'$ (identité si $kk' = 1$).
\item \textbf{Similitude directe} : toute application dont l'écriture analytique est $\begin{pmatrix} x' \\ y' \end{pmatrix} = \begin{pmatrix} a & -b \\ b & a \end{pmatrix}\begin{pmatrix} x \\ y \end{pmatrix} + \begin{pmatrix} c \\ d \end{pmatrix}$ avec $(a\,;b) \neq (0\,;0)$. Son rapport est $k = \sqrt{a^2+b^2}$, son angle $\theta$ vérifie $\cos\theta = \dfrac{a}{k}$ et $\sin\theta = \dfrac{b}{k}$, et si $k \neq 1$ elle admet un unique point fixe : son centre.
\item \textbf{Point fixe} de $F$ : point $\Omega$ tel que $F(\Omega) = \Omega$, obtenu en résolvant le système $\begin{cases} x = ax - by + c \\ y = bx + ay + d \end{cases}$.
\end{itemize}
\end{aretenir}

\courssub{Démarche et corrigé commenté}

\begin{exoresolu}[Corrigé complet de l'activité « Le logo de la CNPS »]
\textbf{Énoncé rappelé.} $A(0\,;0)$, $B(2\,;0)$, $C(0\,;2)$. On applique $R_1$ (rotation centre $A$, $+90^\circ$), puis $H$ (homothétie centre $A$, rapport $2$), puis $t$ (translation de vecteur $(5\,;5)$), puis $R_4$ (rotation centre $\Omega(5\,;5)$, $-90^\circ$). On étudie $F = R_4 \circ t \circ H \circ R_1$.
\tcblower
\textbf{Question 1 : écritures analytiques.}

$R_1$, rotation de centre $A(0\,;0)$ et d'angle $+90^\circ$ : $\begin{cases} x_1 = -y \\ y_1 = x \end{cases}$.

$H$, homothétie de centre $A$ et de rapport $2$ : $\begin{cases} x_2 = 2x_1 \\ y_2 = 2y_1 \end{cases}$.

$t$, translation de vecteur $\vec{u} = \begin{pmatrix} 5 \\ 5 \end{pmatrix}$ : $\begin{cases} x_3 = x_2 + 5 \\ y_3 = y_2 + 5 \end{cases}$.

$R_4$, rotation de centre $\Omega(5\,;5)$ et d'angle $-90^\circ$ : $\begin{cases} x' = 5 + (y_3 - 5) \\ y' = 5 - (x_3 - 5) \end{cases}$, c'est-à-dire $\begin{cases} x' = y_3 \\ y' = 10 - x_3 \end{cases}$.

\textbf{Question 2 : images successives de $A$, $B$, $C$.}

\begin{center}
\begin{tabularx}{\linewidth}{|l|X|X|X|X|}
\hline
\rowcolor{popblueL}
Point & Après $R_1$ & Après $H$ & Après $t$ & Après $R_4$ (final) \\
\hline
$A(0\,;0)$ & $(0\,;0)$ & $(0\,;0)$ & $(5\,;5)$ & $A'(5\,;5)$ \\
\hline
$B(2\,;0)$ & $(0\,;2)$ & $(0\,;4)$ & $(5\,;9)$ & $B'(9\,;5)$ \\
\hline
$C(0\,;2)$ & $(-2\,;0)$ & $(-4\,;0)$ & $(1\,;5)$ & $C'(5\,;9)$ \\
\hline
\end{tabularx}
\end{center}

Détail pour $C$ : $R_1(C) = (-2\,;0)$ ; $H(-2\,;0) = (-4\,;0)$ ; $t(-4\,;0) = (1\,;5)$ ; enfin $R_4(1\,;5) = \bigl(5\,;\,10-1\bigr) = (5\,;9)$.

\textbf{Question 3 : expression globale de $F$.}

On compose les quatre expressions pour un point quelconque $M(x\,;y)$ :
\begin{align*}
x' &= y_3 = y_2 + 5 = 2y_1 + 5 = 2x + 5, \\
y' &= 10 - x_3 = 10 - (x_2 + 5) = 5 - 2x_1 = 5 - 2(-y) = 2y + 5.
\end{align*}
Donc :
\[ \begin{pmatrix} x' \\ y' \end{pmatrix} = \begin{pmatrix} 2 & 0 \\ 0 & 2 \end{pmatrix}\begin{pmatrix} x \\ y \end{pmatrix} + \begin{pmatrix} 5 \\ 5 \end{pmatrix}. \]

\textbf{Question 4 : nature de $F$.}

La matrice $M = \begin{pmatrix} 2 & 0 \\ 0 & 2 \end{pmatrix}$ est bien de la forme $\begin{pmatrix} a & -b \\ b & a \end{pmatrix}$ avec $a = 2$ et $b = 0$. Donc $F$ est une \cle{similitude directe} de rapport
\[ k = \sqrt{a^2 + b^2} = \sqrt{4} = 2, \]
et d'angle $\theta$ tel que $\cos\theta = \dfrac{2}{2} = 1$ et $\sin\theta = 0$, c'est-à-dire $\theta = 0^\circ$. Une similitude directe de rapport $2$ et d'angle nul est une \cle{homothétie} de rapport $2$.

\textbf{Vérification de cohérence :} $R_1$ et $R_4$ sont des rotations d'angles $+90^\circ$ et $-90^\circ$ : la somme des angles est nulle, donc la composée des deux rotations « encadrant » les autres transformations ne fait pas tourner la figure globalement. L'angle nul trouvé est donc logique.

\textbf{Question 5 : centre de $F$.}

Le centre est l'unique point fixe : on résout $F(x\,;y) = (x\,;y)$ :
\[ \begin{cases} x = 2x + 5 \\ y = 2y + 5 \end{cases} \iff \begin{cases} x = -5 \\ y = -5 \end{cases}. \]
$F$ est donc l'\textbf{homothétie de centre $\Omega_F(-5\,;-5)$ et de rapport $2$}.

\textbf{Vérification avec la question 2 :} $A' = (5\,;5)$. Or $\overrightarrow{\Omega_F A} = (5\,;5)$ et $\overrightarrow{\Omega_F A'} = (10\,;10) = 2\overrightarrow{\Omega_F A}$. De même $\overrightarrow{\Omega_F C} = (5\,;7)$ et $2\overrightarrow{\Omega_F C} = (10\,;14) = \overrightarrow{\Omega_F C'}$ puisque $C' - \Omega_F = (5-(-5)\,;\,9-(-5)) = (10\,;14)$. Le résultat est confirmé.

\textbf{Question 6 : vérification graphique.}

Sur papier millimétré, le triangle $ABC$ a pour côtés $AB = AC = 2$ et $BC = 2\sqrt{2}$. Le triangle final $A'B'C'$ avec $A'(5\,;5)$, $B'(9\,;5)$, $C'(5\,;9)$ a pour côtés $A'B' = A'C' = 4$ et $B'C' = 4\sqrt{2}$ : les longueurs sont bien doublées (rapport $2$), et les côtés correspondants sont parallèles (angle nul), ce qui confirme que $F$ est une homothétie.

\textbf{Question 7 : problème inverse.}

On remplace la translation par celle de vecteur $\vec{w} = \begin{pmatrix} p \\ q \end{pmatrix}$. Suivons le point $C(0\,;2)$ à travers la nouvelle chaîne :
\begin{align*}
R_1(C) &= (-2\,;0), \\
H(-2\,;0) &= (-4\,;0), \\
t_{\vec{w}}(-4\,;0) &= (p - 4\,;\,q), \\
R_4(p-4\,;\,q) &= \bigl(5 + (q - 5)\,;\,5 - (p - 4 - 5)\bigr) = (q\,;\,14 - p).
\end{align*}
On exige que l'image finale soit $(10\,;10)$ :
\[ \begin{cases} q = 10 \\ 14 - p = 10 \end{cases} \iff \begin{cases} p = 4 \\ q = 10 \end{cases}. \]
\textbf{Conclusion :} il faut remplacer le vecteur de translation par $\vec{w} = \begin{pmatrix} 4 \\ 10 \end{pmatrix}$. Vérification : $C \mapsto (-2\,;0) \mapsto (-4\,;0) \mapsto (0\,;10) \mapsto \bigl(5 + 5\,;\,5 - (-5)\bigr) = (10\,;10)$. L'objectif du directeur artistique est atteint.
\end{exoresolu}

\begin{attention}[Pièges rencontrés pendant l'activité]
\begin{itemize}
\item \textbf{Ordre de composition :} $F = R_4 \circ t \circ H \circ R_1$ signifie que $R_1$ est appliquée \emph{en premier}. Inverser l'ordre change complètement le résultat : la composition des transformations n'est pas commutative en général.
\item \textbf{Rotation de centre différent de l'origine :} pour $R_4$, il faut d'abord « recentrer » en soustrayant les coordonnées de $\Omega$, tourner, puis « retranslater ». Oublier le recentrage conduit à $x' = y_3$ au lieu de $x' = 5 + (y_3 - 5)$ — ici les deux coïncident car $\Omega$ est sur la première bissectrice, mais ce n'est pas une raison pour sauter l'étape !
\item \textbf{Angle global :} l'angle de la composée est la \emph{somme} des angles des rotations ($+90^\circ + (-90^\circ) = 0^\circ$), pas la différence des valeurs absolues ni un « mélange » avec le rapport de l'homothétie.
\item \textbf{Point fixe :} une homothétie de rapport $k \neq 1$ possède un unique point fixe (son centre). Ne pas conclure « pas de point fixe » si le système donne une solution unique : c'est justement le centre cherché.
\end{itemize}
\end{attention}

\courssub{Bilan de l'activité}

Cette activité a permis de mettre en évidence plusieurs idées fondamentales de la leçon :

\begin{itemize}
\item \textbf{Une chaîne de transformations se résume à une seule transformation.} Quatre transformations successives (deux rotations, une homothétie, une translation) se sont réduites à une unique homothétie de rapport $2$ et de centre $(-5\,;-5)$, que l'on aurait pu appliquer directement. C'est toute la puissance de l'écriture analytique $X' = MX + V$ : elle « absorbe » la composition.
\item \textbf{La forme de la matrice révèle la nature.} Dès que la matrice s'écrit $\begin{pmatrix} a & -b \\ b & a \end{pmatrix}$, on sait qu'on a affaire à une similitude directe ; le rapport $\sqrt{a^2+b^2}$ et l'angle se lisent alors immédiatement. Ici $b = 0$ a trahi une homothétie.
\item \textbf{Les angles s'ajoutent, les rapports se multiplient.} Les deux rotations d'angles opposés se sont « annulées » ($+90^\circ + (-90^\circ) = 0^\circ$), tandis que l'homothétie a imposé son rapport $2$ à la composée.
\item \textbf{Le point fixe caractérise le centre.} La recherche du point invariant $F(\Omega) = \Omega$ est la méthode systématique pour trouver le centre d'une similitude de rapport différent de $1$.
\item \textbf{Le calcul algébrique et la construction graphique se confirment mutuellement.} Le tracé sur papier millimétré a validé les résultats : compétence précieuse au Probatoire pour détecter une erreur de calcul.
\item \textbf{Le problème inverse} a montré qu'on peut \emph{concevoir} une transformation sur mesure pour atteindre un objectif imposé : c'est exactement le travail d'un infographiste, d'un cartographe ou d'un ingénieur.
\end{itemize}

\begin{savaistu}[Les transformations, outils des infographistes]
Les logiciels de conception graphique utilisés par les agences de communication de Douala ou de Yaoundé (comme ceux de la suite Adobe ou les logiciels libres Inkscape et GIMP) représentent chaque déplacement, agrandissement ou pivotement d'un logo par une matrice $2\times 2$ et un vecteur, exactement comme dans cette activité. Composer les transformations revient à multiplier les matrices : c'est ce que fait l'ordinateur des milliers de fois par seconde pour afficher les images à l'écran. Les mathématiques de Première sont donc littéralement au cœur des métiers du numérique, secteur en pleine croissance au Cameroun.
\end{savaistu}

\newpage

\coursec{Méthodes et exercices résolus}

\coursec{Rappels et cadre général : Transformations du plan}

\begin{definition}[Transformation du plan]
Une \cle{transformation} du plan est une application $f$ du plan dans lui-même qui est \textbf{bijective} : à tout point $M$ du plan, on associe un unique point $M' = f(M)$ (image), et réciproquement, tout point $M'$ admet un unique antécédent $M = f^{-1}(M')$.
\end{definition}

\begin{propriete}[Réciproque et composée]
Soient $f$ et $g$ deux transformations du plan.
\begin{itemize}
\item La \cle{réciproque} $f^{-1}$ est la transformation telle que $f^{-1} \circ f = f \circ f^{-1} = \mathrm{id}$ (identité).
\item La \cle{composée} $g \circ f$ est la transformation qui, à $M$, associe $g(f(M))$. L'ordre compte : en général $g \circ f \neq f \circ g$.
\end{itemize}
\end{propriete}

\begin{aretenir}[Définitions vectorielles des transformations élémentaires]
Dans un plan muni d'un repère $(O\,;\,\vec{i},\vec{j})$ :
\begin{itemize}
\item \textbf{Translation} $t_{\vec{u}}$ de vecteur $\vec{u}(a\,;\,b)$ : $\overrightarrow{MM'} = \vec{u}$.
\item \textbf{Homothétie} $h_{\Omega,k}$ de centre $\Omega(a\,;\,b)$ et de rapport $k \neq 0$ : $\overrightarrow{\Omega M'} = k\,\overrightarrow{\Omega M}$.
\item \textbf{Rotation} $r_{\Omega,\alpha}$ de centre $\Omega$ et d'angle $\alpha$ (orienté) : $\Omega M' = \Omega M$ et $\left(\overrightarrow{\Omega M}\,;\,\overrightarrow{\Omega M'}\right) = \alpha\ [2\pi]$.
\end{itemize}
Ces définitions vectorielles sont le point de départ pour établir les définitions analytiques.
\end{aretenir}

\begin{methode}[Montrer qu'une application est une transformation et définir sa réciproque]
Soit $f$ une application du plan dans lui-même définie analytiquement par $\begin{cases} x' = \varphi(x,y) \\ y' = \psi(x,y) \end{cases}$.
\begin{enumerate}
\item \textbf{Montrer que $f$ est une transformation} : prendre un point $M'(x'\,;\,y')$ quelconque et résoudre le système en $x,y$ (inconnues). Montrer qu'il admet une \textbf{unique} solution $(x\,;\,y)$ exprimée en fonction de $x',y'$.
\item La solution donne la \textbf{définition analytique de $f^{-1}$} : $\begin{cases} x = \varphi^{-1}(x',y') \\ y = \psi^{-1}(x',y') \end{cases}$ (on renomme ensuite les variables).
\item \textbf{Caractériser $f^{-1}$} (nature, éléments caractéristiques) en comparant les coefficients (Méthode 3).
\end{enumerate}
\textbf{Rappels utiles :}
\begin{itemize}
\item Réciproque de la translation de vecteur $\vec{u}$ : translation de vecteur $-\vec{u}$.
\item Réciproque de l'homothétie de centre $\Omega$ et rapport $k$ : homothétie de centre $\Omega$ et rapport $\dfrac{1}{k}$.
\item Réciproque de la rotation de centre $\Omega$ et angle $\alpha$ : rotation de centre $\Omega$ et angle $-\alpha$.
\end{itemize}
\end{methode}

\begin{exoresolu}[Transformation et réciproque]
Le plan est muni d'un repère $(O\,;\vec{i},\vec{j})$. On considère l'application $f$ qui, à tout point $M(x\,;\,y)$, associe $M'(x'\,;\,y')$ tel que :
\[ \begin{cases} x' = -x + 4 \\ y' = -y + 2 \end{cases}. \]

1. Montrer que $f$ est une transformation du plan.

2. Définir analytiquement sa réciproque $f^{-1}$ et la caractériser.
\tcblower
\textbf{1. $f$ est une transformation.}

Soit $M'(x'\,;\,y')$ un point quelconque du plan. Cherchons $M(x\,;\,y)$ tel que $f(M)=M'$ :
\[ \begin{cases} -x + 4 = x' \\ -y + 2 = y' \end{cases} \iff \begin{cases} x = -x' + 4 \\ y = -y' + 2 \end{cases}. \]
Pour tout $M'$, ce système admet une \textbf{unique solution}. Donc $f$ est bijective : c'est une \textbf{transformation du plan}.

\textbf{2. Réciproque de $f$.}

D'après la résolution, $f^{-1}$ est définie par (après renommage des variables) :
\[ f^{-1} : \begin{cases} x' = -x + 4 \\ y' = -y + 2 \end{cases}. \]
On constate que $f^{-1} = f$ : $f$ est \textbf{involutive} ($f \circ f = \mathrm{id}$).

\textbf{Caractérisation :} le coefficient commun est $\alpha = -1 \neq 1$, c'est une homothétie de rapport $-1$, c'est-à-dire une \textbf{symétrie centrale}. Son centre $\Omega(a\,;\,b)$ est l'unique point fixe :
\[ \begin{cases} a = -a + 4 \\ b = -b + 2 \end{cases} \iff \begin{cases} 2a = 4 \\ 2b = 2 \end{cases} \iff \Omega(2\,;\,1). \]
\textbf{Conclusion :} $f$ est la symétrie centrale de centre $\Omega(2\,;\,1)$, et $f^{-1}=f$.
\end{exoresolu}

\coursec{Définition analytique de la translation}

\begin{methode}[Définir analytiquement une translation]
Pour déterminer la définition analytique de la translation $t$ de vecteur $\vec{u}(a\,;\,b)$ dans un repère $(O\,;\vec{i},\vec{j})$ :
\begin{enumerate}
\item Poser $M(x\,;\,y)$ un point quelconque et $M'(x'\,;\,y')$ son image par $t$.
\item Traduire la définition vectorielle : $\overrightarrow{MM'} = \vec{u}$.
\item Exprimer les coordonnées de $\overrightarrow{MM'}$ : $(x' - x\,;\, y' - y)$.
\item Identifier les coordonnées : on obtient le système
\[ \begin{cases} x' = x + a \\ y' = y + b \end{cases} \]
\item Conclure : « $t$ a pour définition analytique : $\begin{cases} x' = x + a \\ y' = y + b \end{cases}$ ».
\end{enumerate}
\textbf{Réflexe :} pour l'image, on \emph{ajoute} les coordonnées du vecteur ; pour l'antécédent, on les \emph{soustrait}.
\end{methode}

\begin{exoresolu}[Définition analytique d'une translation]
Le plan est muni d'un repère orthonormé $(O\,;\vec{i},\vec{j})$. On considère les points $A(2\,;\,-1)$ et $B(5\,;\,3)$.

1. Déterminer la définition analytique de la translation $t$ de vecteur $\overrightarrow{AB}$.

2. Déterminer les coordonnées de $C'$, image du point $C(-1\,;\,4)$ par $t$.

3. Déterminer les coordonnées du point $D$ dont l'image par $t$ est $D'(0\,;\,7)$.
\tcblower
\textbf{1. Définition analytique de $t$.}

Les coordonnées du vecteur $\overrightarrow{AB}$ sont :
\[ \overrightarrow{AB} = \begin{pmatrix} x_B - x_A \\ y_B - y_A \end{pmatrix} = \begin{pmatrix} 5 - 2 \\ 3 - (-1) \end{pmatrix} = \begin{pmatrix} 3 \\ 4 \end{pmatrix}. \]
Soit $M(x\,;\,y)$ et $M'(x'\,;\,y') = t(M)$. Par définition :
\[ \overrightarrow{MM'} = \overrightarrow{AB} \iff \begin{cases} x' - x = 3 \\ y' - y = 4 \end{cases} \iff \begin{cases} x' = x + 3 \\ y' = y + 4 \end{cases}. \]
\textbf{Conclusion :} $t$ a pour définition analytique $\begin{cases} x' = x + 3 \\ y' = y + 4 \end{cases}$.

\textbf{2. Image de $C(-1\,;\,4)$.}

On applique la définition avec $x = -1$, $y = 4$ :
\[ \begin{cases} x_{C'} = -1 + 3 = 2 \\ y_{C'} = 4 + 4 = 8 \end{cases} \quad \text{donc} \quad C'(2\,;\,8). \]

\textbf{3. Antécédent de $D'(0\,;\,7)$.}

Soit $D(x\,;\,y)$ tel que $t(D)=D'$ :
\[ \begin{cases} x + 3 = 0 \\ y + 4 = 7 \end{cases} \iff \begin{cases} x = -3 \\ y = 3 \end{cases} \quad \text{donc} \quad D(-3\,;\,3). \]
\textbf{Vérification :} $t(D)$ donne $(-3+3\,;\,3+4) = (0\,;\,7) = D'$. Cohérent.
\end{exoresolu}

\coursec{Définition analytique de l'homothétie}

\begin{methode}[Définir analytiquement une homothétie]
Pour déterminer la définition analytique de l'homothétie $h$ de centre $\Omega(a\,;\,b)$ et de rapport $k$ ($k \neq 0$) :
\begin{enumerate}
\item Poser $M(x\,;\,y)$ et $M'(x'\,;\,y')$ son image par $h$.
\item Traduire la définition vectorielle : $\overrightarrow{\Omega M'} = k\,\overrightarrow{\Omega M}$.
\item Écrire les coordonnées : $\overrightarrow{\Omega M}(x - a\,;\, y - b)$ et $\overrightarrow{\Omega M'}(x' - a\,;\, y' - b)$.
\item Identifier et isoler $x'$ et $y'$ :
\[ \begin{cases} x' = kx + (1 - k)a \\ y' = ky + (1 - k)b \end{cases} \]
\end{enumerate}
\textbf{Formule à retenir (souvent plus rapide) :}
\[ \begin{cases} x' - a = k(x - a) \\ y' - b = k(y - b) \end{cases} \]
\end{methode}

\begin{exoresolu}[Définition analytique d'une homothétie]
Le plan est muni d'un repère orthonormé $(O\,;\vec{i},\vec{j})$. On considère l'homothétie $h$ de centre $\Omega(1\,;\,2)$ et de rapport $k = -2$.

1. Déterminer la définition analytique de $h$.

2. Déterminer l'image du point $A(3\,;\,-1)$ par $h$.

3. Déterminer l'antécédent du point $B'(7\,;\,-2)$ par $h$.
\tcblower
\textbf{1. Définition analytique de $h$.}

Soit $M(x\,;\,y)$ et $M'(x'\,;\,y') = h(M)$. Par définition :
\[ \overrightarrow{\Omega M'} = -2\,\overrightarrow{\Omega M} \iff \begin{cases} x' - 1 = -2(x - 1) \\ y' - 2 = -2(y - 2) \end{cases}. \]
On développe et on isole $x',y'$ :
\[ \begin{cases} x' = -2x + 2 + 1 = -2x + 3 \\ y' = -2y + 4 + 2 = -2y + 6 \end{cases}. \]
\textbf{Conclusion :} $h$ a pour définition analytique $\begin{cases} x' = -2x + 3 \\ y' = -2y + 6 \end{cases}$.

\textbf{2. Image de $A(3\,;\,-1)$.}
\[ \begin{cases} x_{A'} = -2 \times 3 + 3 = -3 \\ y_{A'} = -2 \times (-1) + 6 = 8 \end{cases} \quad \text{donc} \quad A'(-3\,;\,8). \]

\textbf{3. Antécédent de $B'(7\,;\,-2)$.}

Soit $B(x\,;\,y)$ tel que $h(B)=B'$ :
\[ \begin{cases} -2x + 3 = 7 \\ -2y + 6 = -2 \end{cases} \iff \begin{cases} -2x = 4 \\ -2y = -8 \end{cases} \iff \begin{cases} x = -2 \\ y = 4 \end{cases}. \]
Donc $B(-2\,;\,4)$.

\textbf{Vérification vectorielle :} $\overrightarrow{\Omega B}(-3\,;\,2)$ et $\overrightarrow{\Omega B'}(6\,;\,-4) = -2\,\overrightarrow{\Omega B}$. Confirmé.
\end{exoresolu}

\coursec{Caractérisation d'une transformation à partir de son expression analytique}

\begin{methode}[Reconnaître une translation ou une homothétie]
On donne une application $f$ du plan définie par $\begin{cases} x' = \alpha x + p \\ y' = \alpha y + q \end{cases}$ (même coefficient $\alpha$ devant $x$ et $y$).
\begin{enumerate}
\item \textbf{Comparer les coefficients} de $x$ et $y$ : s'ils sont différents, ce n'est ni une translation ni une homothétie (c'est une affinité).
\item \textbf{Cas $\alpha = 1$ :} $f$ est la \textbf{translation} de vecteur $\vec{u}(p\,;\,q)$ (si $p=q=0$, c'est l'identité).
\item \textbf{Cas $\alpha \neq 1$ :} $f$ est une \textbf{homothétie} de rapport $k = \alpha$. Le centre $\Omega(a\,;\,b)$ est l'unique \cle{point fixe} : on résout
\[ \begin{cases} a = \alpha a + p \\ b = \alpha b + q \end{cases} \iff \begin{cases} a = \dfrac{p}{1 - \alpha} \\[6pt] b = \dfrac{q}{1 - \alpha} \end{cases}. \]
\item Conclure en donnant la \textbf{nature} et les \textbf{éléments caractéristiques}.
\end{enumerate}
\textbf{Réflexe :} le centre d'une homothétie est son unique point fixe : on le cherche en posant $M' = M$.
\end{methode}

\begin{exoresolu}[Nature d'une application définie analytiquement]
Le plan est muni d'un repère $(O\,;\vec{i},\vec{j})$. Dans chaque cas, déterminer la nature et les éléments caractéristiques de l'application $f$.

1. $f_1 : \begin{cases} x' = x - 5 \\ y' = y + 2 \end{cases}$ \qquad 2. $f_2 : \begin{cases} x' = 3x - 4 \\ y' = 3y + 8 \end{cases}$ \qquad 3. $f_3 : \begin{cases} x' = 2x + 1 \\ y' = 3y - 2 \end{cases}$
\tcblower
\textbf{1. Étude de $f_1$.}

Coefficient de $x$ et $y$ vaut $1$ : c'est une \textbf{translation} de vecteur $\vec{u}(-5\,;\,2)$.

\textbf{2. Étude de $f_2$.}

Coefficient commun $\alpha = 3 \neq 1$ : \textbf{homothétie de rapport $3$}. Centre $\Omega(a\,;\,b)$ (point fixe) :
\[ \begin{cases} a = 3a - 4 \\ b = 3b + 8 \end{cases} \iff \begin{cases} -2a = -4 \\ -2b = 8 \end{cases} \iff \Omega(2\,;\,-4). \]
\textbf{Vérification :} $M=\Omega(2\,;\,-4) \Rightarrow x'=3\times2-4=2,\ y'=3\times(-4)+8=-4$. Point fixe.

\textbf{Conclusion :} $f_2$ est l'homothétie de centre $\Omega(2\,;\,-4)$ et de rapport $k=3$.

\textbf{3. Étude de $f_3$.}

Coefficient de $x$ est $2$, celui de $y$ est $3$ : \textbf{différents}. $f_3$ n'est \textbf{ni une translation, ni une homothétie}.
\end{exoresolu}

\coursec{Composée de deux rotations de même centre}

\begin{methode}[Composée de deux rotations de même centre]
Soient $r_1$ et $r_2$ deux rotations de \textbf{même centre} $\Omega$ et d'angles respectifs $\alpha$ et $\beta$ (orientés).
\begin{enumerate}
\item Calculer la \textbf{somme des angles} : $\theta = \alpha + \beta$.
\item Si $\theta \neq 0\ [2\pi]$ : $r_2 \circ r_1$ est la \textbf{rotation de centre $\Omega$ et d'angle $\theta$}.
\item Si $\theta = 0\ [2\pi]$ : $r_2 \circ r_1$ est l'\textbf{identité} du plan ($r_2 = r_1^{-1}$).
\end{enumerate}
\textbf{Remarque :} l'ordre n'a pas d'importance : $r_2 \circ r_1 = r_1 \circ r_2$ (même centre).
\end{methode}

\begin{exoresolu}[Composée de deux rotations de même centre]
Le plan est muni d'un repère orthonormé direct $(O\,;\vec{i},\vec{j})$. Soit $\Omega$ un point du plan. On considère les rotations $r_1$ de centre $\Omega$ et d'angle $\dfrac{\pi}{3}$, et $r_2$ de centre $\Omega$ et d'angle $\dfrac{\pi}{4}$.

1. Déterminer la nature et les éléments caractéristiques de $r = r_2 \circ r_1$.

2. Que vaut $r_1 \circ r_3$ où $r_3$ est la rotation de centre $\Omega$ et d'angle $-\dfrac{\pi}{3}$ ?
\tcblower
\textbf{1. Nature de $r = r_2 \circ r_1$.}

Même centre $\Omega$. Somme des angles :
\[ \theta = \dfrac{\pi}{3} + \dfrac{\pi}{4} = \dfrac{4\pi}{12} + \dfrac{3\pi}{12} = \dfrac{7\pi}{12}. \]
Comme $\dfrac{7\pi}{12} \neq 0\ [2\pi]$, $r$ est la \textbf{rotation de centre $\Omega$ et d'angle $\dfrac{7\pi}{12}$}.

\textbf{2. Nature de $r_1 \circ r_3$.}

Somme des angles : $\dfrac{\pi}{3} + \left(-\dfrac{\pi}{3}\right) = 0$. Donc $r_1 \circ r_3$ est l'\textbf{identité du plan} : $r_3$ est la réciproque de $r_1$.
\end{exoresolu}

\coursec{Composée de deux homothéties de même centre}

\begin{methode}[Composée de deux homothéties de même centre]
Soient $h_1$ et $h_2$ deux homothéties de \textbf{même centre} $\Omega$ et de rapports respectifs $k_1$ et $k_2$ ($k_1,k_2 \neq 0$).
\begin{enumerate}
\item Calculer le \textbf{produit des rapports} : $k = k_1 \times k_2$.
\item \textbf{Cas $k \neq 1$ :} $h_2 \circ h_1$ est l'\textbf{homothétie de centre $\Omega$ et de rapport $k$}.
\item \textbf{Cas $k = 1$ :} $h_2 \circ h_1$ est l'\textbf{identité} du plan ($h_2 = h_1^{-1}$).
\end{enumerate}
\textbf{Réflexe :} l'ordre n'a pas d'importance ($h_2 \circ h_1 = h_1 \circ h_2$). \textbf{Attention :} on \textbf{multiplie} les rapports (on n'additionne pas !).
\end{methode}

\begin{exoresolu}[Composée de deux homothéties de même centre]
Le plan est muni d'un repère $(O\,;\vec{i},\vec{j})$. Soit $\Omega(1\,;\,-1)$. On considère $h_1$ de centre $\Omega$ et rapport $2$, et $h_2$ de centre $\Omega$ et rapport $-\dfrac{1}{4}$.

1. Déterminer la nature et les éléments caractéristiques de $f = h_2 \circ h_1$.

2. Donner la définition analytique de $f$, puis l'image de $A(3\,;\,5)$ par $f$.

3. Que vaut $g = h_1 \circ h_3$ où $h_3$ est l'homothétie de centre $\Omega$ et de rapport $\dfrac{1}{2}$ ?
\tcblower
\textbf{1. Nature de $f = h_2 \circ h_1$.}

Même centre $\Omega$. Produit des rapports :
\[ k = k_1 \times k_2 = 2 \times \left(-\dfrac{1}{4}\right) = -\dfrac{1}{2}. \]
Comme $k \neq 1$, $f$ est l'\textbf{homothétie de centre $\Omega(1\,;\,-1)$ et de rapport $-\dfrac{1}{2}$}.

\textbf{2. Définition analytique de $f$.}

Soit $M(x\,;\,y)$ et $M'(x'\,;\,y') = f(M)$. $\overrightarrow{\Omega M'} = -\dfrac{1}{2}\overrightarrow{\Omega M}$ :
\[ \begin{cases} x' - 1 = -\dfrac{1}{2}(x - 1) \\[4pt] y' + 1 = -\dfrac{1}{2}(y + 1) \end{cases} \iff \begin{cases} x' = -\dfrac{1}{2}x + \dfrac{3}{2} \\[4pt] y' = -\dfrac{1}{2}y - \dfrac{3}{2} \end{cases}. \]
Image de $A(3\,;\,5)$ :
\[ x_{A'} = -\dfrac{3}{2} + \dfrac{3}{2} = 0 \quad ; \quad y_{A'} = -\dfrac{5}{2} - \dfrac{3}{2} = -4. \]
Donc $A'(0\,;\,-4)$.

\textbf{3. Nature de $g = h_1 \circ h_3$.}

Produit des rapports : $2 \times \dfrac{1}{2} = 1$. Donc $g$ est l'\textbf{identité du plan} : $h_3 = h_1^{-1}$.
\end{exoresolu}

\coursec{Composée d'une homothétie et d'une translation}

\begin{methode}[Composée d'une homothétie et d'une translation]
Soient $h$ l'homothétie de centre $\Omega$ et de rapport $k$, et $t$ la translation de vecteur $\vec{u}$. On étudie $f = t \circ h$ (ou $h \circ t$).
\begin{enumerate}
\item \textbf{Regarder le rapport $k$ :} c'est lui qui décide de la nature.
\item \textbf{Cas $k \neq 1$ :} $f$ est une \textbf{homothétie de rapport $k$}, mais de centre \textbf{différent de $\Omega$} en général. Pour trouver le nouveau centre $\Omega'$, on cherche l'\textbf{unique point fixe} de $f$ en résolvant $f(M) = M$ (analytiquement).
\item \textbf{Cas $k = 1$ :} $h$ est l'identité, donc $f = t$ (translation).
\item \textbf{Attention :} en général $t \circ h \neq h \circ t$ (les centres des deux composées sont différents). Respecter l'ordre : $t \circ h$ signifie « $h$ d'abord, puis $t$ ».
\end{enumerate}
\textbf{Astuce vectorielle (optionnelle) :} pour $f = t \circ h$ avec $k \neq 1$, $\overrightarrow{\Omega \Omega'} = \dfrac{1}{1-k}\,\vec{u}$. Mais la résolution analytique du point fixe est plus sûre au Probatoire.
\end{methode}

\begin{exoresolu}[Composée d'une homothétie et d'une translation]
Le plan est muni d'un repère orthonormé $(O\,;\vec{i},\vec{j})$. On considère $h$ homothétie de centre $\Omega(2\,;\,1)$ et rapport $k=3$, et $t$ translation de vecteur $\vec{u}(4\,;\,-6)$.

1. Définition analytique de $f = t \circ h$.

2. Nature et éléments caractéristiques de $f$.

3. On pose $g = h \circ t$. A-t-on $g = f$ ? Justifier.
\tcblower
\textbf{1. Définition analytique de $f = t \circ h$.}

Soit $M(x\,;\,y)$, $M_1 = h(M)$, $M' = t(M_1)$.

\textbf{Étape 1 — $h$ :} $\overrightarrow{\Omega M_1} = 3\,\overrightarrow{\Omega M}$ :
\[ \begin{cases} x_1 - 2 = 3(x - 2) \\ y_1 - 1 = 3(y - 1) \end{cases} \iff \begin{cases} x_1 = 3x - 4 \\ y_1 = 3y - 2 \end{cases}. \]

\textbf{Étape 2 — $t$ :} $\overrightarrow{M_1 M'} = \vec{u}$ :
\[ \begin{cases} x' = x_1 + 4 = 3x - 4 + 4 = 3x \\ y' = y_1 - 6 = 3y - 2 - 6 = 3y - 8 \end{cases}. \]
Donc $f : \begin{cases} x' = 3x \\ y' = 3y - 8 \end{cases}$.

\textbf{2. Nature de $f$.}

Coefficient commun $3 \neq 1$ : $f$ est une \textbf{homothétie de rapport $3$}. Centre $\Omega'(a\,;\,b)$ (point fixe) :
\[ \begin{cases} a = 3a \\ b = 3b - 8 \end{cases} \iff \begin{cases} -2a = 0 \\ -2b = -8 \end{cases} \iff \Omega'(0\,;\,4). \]
\textbf{Conclusion :} $f = t \circ h$ est l'homothétie de centre $\Omega'(0\,;\,4)$ et de rapport $3$. $\Omega' \neq \Omega$.

\textbf{3. Comparaison $g = h \circ t$ et $f$.}

$g = h \circ t$ : $M_2 = t(M)$ puis $M'' = h(M_2)$.
\[ t : \begin{cases} x_2 = x + 4 \\ y_2 = y - 6 \end{cases} \quad ; \quad h : \begin{cases} x'' = 3x_2 - 4 = 3x + 8 \\ y'' = 3y_2 - 2 = 3y - 20 \end{cases}. \]
$g$ est homothétie de rapport $3$, centre $\Omega''$ : $a = 3a+8 \Rightarrow a=-4$, $b=3b-20 \Rightarrow b=10$. Soit $\Omega''(-4\,;\,10)$.
Les centres $\Omega'(0\,;\,4)$ et $\Omega''(-4\,;\,10)$ sont distincts $\Rightarrow$ $\mathbf{g \neq f}$ : non commutativité.
\end{exoresolu}

\coursec{Composée d'une homothétie et d'une rotation de même centre (Similitude directe)}

\begin{methode}[Composée homothétie et rotation de même centre — Similitude directe]
Soient $h$ l'homothétie de centre $\Omega$ et de rapport $k > 0$, et $r$ la rotation de centre $\Omega$ et d'angle $\theta$.
\begin{enumerate}
\item La composée $s = h \circ r = r \circ h$ (l'ordre n'a pas d'importance car même centre) est une \cle{similitude directe} de centre $\Omega$, de rapport $k$ et d'angle $\theta$.
\item Une similitude directe transforme toute figure en une figure \textbf{semblable} : les angles sont conservés, les longueurs sont multipliées par $k$, les aires par $k^2$.
\end{enumerate}
\textbf{Réflexe rédaction :} pour une similitude, toujours annoncer la \textbf{nature} puis les \textbf{trois éléments caractéristiques} : centre, rapport, angle.
\end{methode}

\begin{exoresolu}[Similitude directe]
Le plan est muni d'un repère orthonormé direct $(O\,;\vec{i},\vec{j})$. Soit $\Omega$ un point du plan. Soit $h$ l'homothétie de centre $\Omega$ et de rapport $2$, et $r_1$ la rotation de centre $\Omega$ et d'angle $\dfrac{\pi}{3}$.

Caractériser $s = h \circ r_1$. Quelle est l'image par $s$ d'un segment $[AB]$ de longueur $5$~cm ?
\tcblower
\textbf{Nature de $s = h \circ r_1$.}

$h$ et $r_1$ ont le même centre $\Omega$, et le rapport $k=2$ est strictement positif. Donc $s$ est la \textbf{similitude directe de centre $\Omega$, de rapport $2$ et d'angle $\dfrac{\pi}{3}$}.

\textbf{Image du segment $[AB]$ :} une similitude de rapport $2$ multiplie toutes les longueurs par $2$. L'image de $[AB]$ est un segment $[A'B']$ tel que :
\[ A'B' = 2 \times AB = 2 \times 5 = 10\ \text{cm}. \]
De plus, l'angle orienté $\left(\overrightarrow{AB}\,;\,\overrightarrow{A'B'}\right)$ vaut $\dfrac{\pi}{3}$ : le segment image a « tourné » de $\dfrac{\pi}{3}$ par rapport au segment initial.
\end{exoresolu}

\begin{attention}[Les 5 erreurs qui coûtent des points au Probatoire]
\begin{enumerate}
\item \textbf{Confondre image et antécédent.} Dans $\begin{cases} x' = x + a \\ y' = y + b \end{cases}$, $(x'\,;\,y')$ est l'\textbf{image}. Pour un antécédent, on résout en $x,y$. Lire deux fois : « image de $A$ » $\neq$ « point dont l'image est $A$ ».
\item \textbf{Oublier la condition $k \neq 1$ pour le centre.} L'unique point fixe n'existe que si $k \neq 1$. Si $k_1 k_2 = 1$ (homothéties) ou $k=1$ (homothétie-translation), ne pas chercher de centre : la composée est l'\textbf{identité} ou une \textbf{translation}. Annoncer un « centre » ici est une erreur grave.
\item \textbf{Additionner les rapports au lieu de les multiplier.} Composée de deux homothéties même centre $\to$ rapport = \textbf{produit} $k_1 k_2$. Composée de deux rotations même centre $\to$ angle = \textbf{somme} $\alpha+\beta$. Ne pas mélanger les règles.
\item \textbf{Croire que $t \circ h = h \circ t$.} La composée homothétie-translation n'est \textbf{pas commutative} : même rapport, centres différents. Respecter l'ordre : $t \circ h$ = $h$ puis $t$.
\item \textbf{Rédaction incomplète.} « C'est une homothétie » ne suffit pas : il faut la \textbf{nature} ET les \textbf{éléments caractéristiques} (centre+rapport ; vecteur ; centre+rapport+angle). Pour prouver « transformation », écrire explicitement « tout point admet un \textbf{unique} antécédent » — sans « unique », justification incomplète.
\end{enumerate}
\end{attention}

\newpage

\coursec{Exercices}

\courssub{Je vérifie mes connaissances}

\begin{exercice}[QCM : une seule bonne réponse]
Pour chaque question, entoure la lettre de la bonne réponse.

\textbf{1.} L'application $f$ du plan dans lui-même qui, à tout point $M(x\,;y)$, associe $M'(x'\,;y')$ tel que $\begin{cases} x' = x + 3 \\ y' = y - 2 \end{cases}$ est :
\begin{itemize}
\item[a)] une homothétie de rapport $3$ ;
\item[b)] la translation de vecteur $\vec{u}\begin{pmatrix} 3 \\ -2 \end{pmatrix}$ ;
\item[c)] une rotation d'angle $\dfrac{\pi}{2}$.
\end{itemize}

\textbf{2.} L'homothétie de centre $\Omega(1\,;-2)$ et de rapport $k = -3$ a pour écriture analytique :
\begin{itemize}
\item[a)] $\begin{cases} x' = -3x + 4 \\ y' = -3y - 8 \end{cases}$ ;
\item[b)] $\begin{cases} x' = -3x - 2 \\ y' = -3y + 4 \end{cases}$ ;
\item[c)] $\begin{cases} x' = 3x - 2 \\ y' = 3y + 4 \end{cases}$.
\end{itemize}

\textbf{3.} La composée de deux rotations de même centre $O$, d'angles $\dfrac{\pi}{3}$ et $\dfrac{\pi}{6}$, est :
\begin{itemize}
\item[a)] une rotation de centre $O$ et d'angle $\dfrac{\pi}{2}$ ;
\item[b)] une rotation de centre $O$ et d'angle $\dfrac{\pi}{18}$ ;
\item[c)] une translation.
\end{itemize}

\textbf{4.} La composée de deux homothéties de même centre $O$ et de rapports $k = 2$ et $k' = \dfrac{1}{2}$ est :
\begin{itemize}
\item[a)] l'homothétie de centre $O$ et de rapport $1$, c'est-à-dire l'identité ;
\item[b)] l'homothétie de centre $O$ et de rapport $\dfrac{5}{2}$ ;
\item[c)] une translation.
\end{itemize}
\end{exercice}

\begin{exercice}[Vrai ou faux ? Justifier]
Réponds par VRAI ou FAUX en justifiant chaque réponse par une propriété du cours ou un contre-exemple.
\begin{enumerate}
\item Toute transformation du plan admet une application réciproque qui est aussi une transformation.
\item La composée d'une homothétie de rapport $k$ et d'une translation est toujours une homothétie, quel que soit $k$.
\item La composée d'une homothétie de centre $O$ et de rapport $k > 0$ et d'une rotation de même centre $O$ et d'angle $\theta$ est une similitude directe de rapport $k$ et d'angle $\theta$.
\item Si $f$ a pour écriture analytique $\begin{cases} x' = 2x + 1 \\ y' = 2y - 3 \end{cases}$, alors $f$ est une homothétie de rapport $2$.
\end{enumerate}
\end{exercice}

\begin{exercice}[Reconnaître une transformation]
Dans chaque cas, l'application $f$ est définie par son écriture analytique dans un repère orthonormé $(O\,;\vec{i},\vec{j})$. Indique sa nature (translation, homothétie, identité) et précise ses éléments caractéristiques.
\begin{enumerate}
\item $\begin{cases} x' = x - 4 \\ y' = y + 1 \end{cases}$
\item $\begin{cases} x' = 5x \\ y' = 5y \end{cases}$
\item $\begin{cases} x' = -x + 6 \\ y' = -y - 2 \end{cases}$
\item $\begin{cases} x' = x \\ y' = y \end{cases}$
\end{enumerate}
\end{exercice}

\begin{exercice}[Calculs directs]
Le plan est muni d'un repère orthonormé $(O\,;\vec{i},\vec{j})$.
\begin{enumerate}
\item Donne l'écriture analytique de la translation $t$ de vecteur $\vec{u}\begin{pmatrix} -2 \\ 5 \end{pmatrix}$, puis calcule l'image du point $A(3\,;-1)$ par $t$.
\item Donne l'écriture analytique de l'homothétie $h$ de centre $\Omega(2\,;1)$ et de rapport $k = -2$, puis calcule l'image du point $B(0\,;3)$ par $h$.
\item Soit $r_1$ la rotation de centre $O$ et d'angle $\dfrac{\pi}{4}$, et $r_2$ la rotation de centre $O$ et d'angle $\dfrac{\pi}{4}$. Quelle est la nature de $r_2 \circ r_1$ ? Précise ses éléments caractéristiques.
\end{enumerate}
\end{exercice}

\courssub{Je m'entraîne}

\begin{exercice}[Écriture analytique d'une homothétie]
Le plan est muni d'un repère orthonormé $(O\,;\vec{i},\vec{j})$. On considère l'homothétie $h$ de centre $\Omega(-1\,;2)$ et de rapport $k = 3$.
\begin{enumerate}
\item En utilisant la relation vectorielle $\overrightarrow{\Omega M'} = 3\,\overrightarrow{\Omega M}$, établis l'écriture analytique de $h$.
\item Détermine les images des points $A(1\,;0)$, $B(-1\,;2)$ et $C(0\,;-3)$ par $h$. Que remarques-tu pour le point $B$ ? Justifie.
\item Détermine l'écriture analytique de l'application réciproque $h^{-1}$. Quelle est sa nature ?
\end{enumerate}
\end{exercice}

\begin{exercice}[Composée de deux homothéties de même centre]
Le plan est muni d'un repère orthonormé $(O\,;\vec{i},\vec{j})$. On considère les homothéties $h_1$ de centre $I(2\,;-1)$ et de rapport $k_1 = 4$, et $h_2$ de même centre $I$ et de rapport $k_2 = -\dfrac{1}{2}$.
\begin{enumerate}
\item Écris les expressions analytiques de $h_1$ et de $h_2$.
\item Détermine l'écriture analytique de $f = h_2 \circ h_1$.
\item Quelle est la nature de $f$ ? Précise son centre et son rapport. Le résultat était-il prévisible sans calcul ?
\item Mêmes questions pour $g = h_1 \circ h_2$. Que constates-tu ?
\end{enumerate}
\end{exercice}

\begin{exercice}[Composée d'une homothétie et d'une translation]
Le plan est muni d'un repère orthonormé $(O\,;\vec{i},\vec{j})$. Soit $t$ la translation de vecteur $\vec{u}\begin{pmatrix} 3 \\ 1 \end{pmatrix}$ et $h$ l'homothétie de centre $O$ et de rapport $k = 2$.
\begin{enumerate}
\item Détermine l'écriture analytique de $f = t \circ h$.
\item Justifie que $f$ est une homothétie dont on précisera le rapport, puis calcule les coordonnées de son centre $\Omega$.
\item Détermine l'écriture analytique de $g = h \circ t$. A-t-on $f = g$ ? Calcule le centre de $g$ et compare avec celui de $f$.
\end{enumerate}
\end{exercice}

\begin{exercice}[Composée homothétie--rotation de même centre]
Le plan est muni d'un repère orthonormé direct $(O\,;\vec{i},\vec{j})$. On considère la rotation $r$ de centre $\Omega(1\,;1)$ et d'angle $\dfrac{\pi}{2}$, et l'homothétie $h$ de même centre $\Omega$ et de rapport $k = 3$.
\begin{enumerate}
\item Rappelle l'écriture analytique de la rotation $r$.
\item Détermine l'écriture analytique de $s = h \circ r$.
\item Quelle est la nature de $s$ ? Précise son centre, son rapport et son angle.
\item Calcule l'image du point $A(2\,;1)$ par $s$, puis vérifie géométriquement : que valent la distance $\Omega A'$ et l'angle $\left(\overrightarrow{\Omega A}\,;\overrightarrow{\Omega A'}\right)$ ?
\end{enumerate}
\end{exercice}

\courssub{Je me prépare au Probatoire}

\begin{exercice}[Type Probatoire : étude complète d'une application du plan]
Le plan est muni d'un repère orthonormé direct $(O\,;\vec{i},\vec{j})$. On considère l'application $f$ du plan dans lui-même qui, à tout point $M(x\,;y)$, associe le point $M'(x'\,;y')$ tel que :
\[
\begin{cases} x' = 2x - y + 1 \\ y' = x + 2y - 3 \end{cases}
\]
\begin{enumerate}
\item \textbf{$f$ est une transformation.}
\begin{enumerate}
\item Montre que pour tout point $M'(x'\,;y')$ donné, il existe un unique point $M(x\,;y)$ tel que $f(M) = M'$. (On résoudra le système d'inconnues $x$ et $y$.)
\item En déduis que $f$ est une transformation du plan et donne l'écriture analytique de sa réciproque $f^{-1}$.
\end{enumerate}
\item \textbf{Nature de $f$.}
\begin{enumerate}
\item Vérifie que l'écriture de $f$ est de la forme $\begin{cases} x' = ax - by + c \\ y' = bx + ay + c' \end{cases}$ et calcule $a^2 + b^2$.
\item On pose $k = \sqrt{a^2+b^2}$. Montre qu'il existe un réel $\theta$ tel que $\cos\theta = \dfrac{a}{k}$ et $\sin\theta = \dfrac{b}{k}$, et donne une mesure de $\theta$.
\item Détermine les coordonnées de l'unique point fixe $\Omega$ de $f$.
\item Conclus : donne la nature exacte de $f$ et l'ensemble de ses éléments caractéristiques (centre, rapport, angle).
\end{enumerate}
\item \textbf{Application.} Soit $A(1\,;0)$ et $B(0\,;1)$.
\begin{enumerate}
\item Calcule les coordonnées de $A' = f(A)$ et $B' = f(B)$.
\item Calcule les distances $AB$ et $A'B'$. Vérifie que $A'B' = k \times AB$.
\end{enumerate}
\end{enumerate}
\end{exercice}

\begin{exercice}[Type Probatoire : chantier de construction à Yaoundé]
Une entreprise de BTP modélise sur un plan, muni d'un repère orthonormé $(O\,;\vec{i},\vec{j})$ d'unité \SI{1}{cm} pour \SI{10}{m}, le déplacement d'une grue. Le plan d'une dalle rectangulaire est le rectangle $ABCD$ avec $A(0\,;0)$, $B(4\,;0)$, $C(4\,;2)$ et $D(0\,;2)$. On considère :
\begin{itemize}
\item la translation $t$ de vecteur $\vec{u}\begin{pmatrix} 2 \\ -1 \end{pmatrix}$ ;
\item l'homothétie $h$ de centre $O$ et de rapport $k = -\dfrac{1}{2}$.
\end{itemize}
\begin{enumerate}
\item \textbf{Étude de $f = t \circ h$.}
\begin{enumerate}
\item Détermine l'écriture analytique de $f$.
\item Montre que $f$ est une homothétie dont on précisera le rapport, et calcule les coordonnées de son centre $\Omega$.
\item Le centre $\Omega$ est-il l'image de $O$ par $t$ ? Justifie ta réponse.
\end{enumerate}
\item \textbf{Image de la dalle.}
\begin{enumerate}
\item Construis soigneusement le rectangle $ABCD$ et son image $A'B'C'D'$ par $f$ (figure à rendre avec la copie).
\item Calcule l'aire du rectangle $ABCD$ en $\mathrm{cm}^2$, puis celle de $A'B'C'D'$. Par quel nombre l'aire a-t-elle été multipliée ? Retrouve ce résultat par une propriété du cours.
\item Quelle est l'aire réelle, en $\mathrm{m}^2$, de la dalle représentée par $A'B'C'D'$ ?
\end{enumerate}
\item \textbf{Étude de $g = h \circ t$.}
\begin{enumerate}
\item Détermine l'écriture analytique de $g$ et les coordonnées de son centre $\Omega'$.
\item Les homothéties $f$ et $g$ ont-elles le même centre ? Calcule la distance $\Omega\Omega'$.
\end{enumerate}
\end{enumerate}
\end{exercice}

\begin{popfigure}
\begin{tikzpicture}[scale=1, >=Stealth]
    % Axes
    \draw[->, thin, popmuted] (-1,0) -- (5,0) node[right] {$x$};
    \draw[->, thin, popmuted] (0,-3) -- (0,3) node[above] {$y$};
    \foreach \x in {1,2,3,4} \draw (\x,0.1) -- (\x,-0.1) node[below, font=\tiny] {\x};
    \foreach \y in {-2,-1,1,2} \draw (0.1,\y) -- (-0.1,\y) node[left, font=\tiny] {\y};
    % Grid
    \draw[popline, step=1] (-0.5,-2.5) grid (4.5,2.5);
    % Rectangle ABCD
    \draw[popblue, thick] (0,0) -- (4,0) -- (4,2) -- (0,2) -- cycle;
    \node[popblue, below left] at (0,0) {$A(0,0)$};
    \node[popblue, below right] at (4,0) {$B(4,0)$};
    \node[popblue, above right] at (4,2) {$C(4,2)$};
    \node[popblue, above left] at (0,2) {$D(0,2)$};
    % Rectangle A'B'C'D' (image by f = t o h, k=-1/2, center (4/3, -2/3))
    % A(0,0) -> A'(2,-1)
    % B(4,0) -> B'(0,-1)
    % C(4,2) -> C'(0,-2)
    % D(0,2) -> D'(2,-2)
    \draw[poporange, thick] (2,-1) -- (0,-1) -- (0,-2) -- (2,-2) -- cycle;
    \node[poporange, below right] at (2,-1) {$A'(2,-1)$};
    \node[poporange, below left] at (0,-1) {$B'(0,-1)$};
    \node[poporange, above left] at (0,-2) {$C'(0,-2)$};
    \node[poporange, above right] at (2,-2) {$D'(2,-2)$};
    % Centers
    \node[poppurple, circle, fill, inner sep=1.5pt, label={above right:$\Omega(\frac{4}{3},-\frac{2}{3})$}] at (1.333,-0.666) {};
    \node[popgreen, circle, fill, inner sep=1.5pt, label={below left:$\Omega'(-\frac{2}{3},\frac{1}{3})$}] at (-0.666,0.333) {};
    % Vectors for translation
    \draw[poppink, ->, thick] (0,0) -- (2,-1) node[midway, above right] {$\vec{u}$};
\end{tikzpicture}
\legende{Rectangle $ABCD$ (bleu) et son image $A'B'C'D'$ (orange) par $f = t \circ h$. Centres $\Omega$ (violet) de $f$ et $\Omega'$ (vert) de $g = h \circ t$.}
\end{popfigure}

\begin{exercice}[Type Probatoire : triangle équilatéral et similitude]
Dans le plan orienté, on considère un triangle équilatéral $ABC$ de centre $O$, avec $\left(\overrightarrow{OA}\,;\overrightarrow{OB}\right) = \dfrac{2\pi}{3}$. On note $r$ la rotation de centre $O$ et d'angle $\dfrac{2\pi}{3}$, et $h$ l'homothétie de centre $O$ et de rapport $\dfrac{1}{2}$. On pose $f = r \circ h$.
\begin{enumerate}
\item \textbf{Nature de $f$.}
\begin{enumerate}
\item Justifie que $f$ est une similitude directe dont on précisera le centre, le rapport et l'angle.
\item $f$ est-elle une rotation ? Une homothétie ? Justifie.
\end{enumerate}
\item \textbf{Image du triangle.}
\begin{enumerate}
\item Détermine $r(A)$, $r(B)$ et $r(C)$ en justifiant.
\item On note $A' = f(A)$, $B' = f(B)$ et $C' = f(C)$. Montre que $A'$, $B'$, $C'$ sont les milieux respectifs des segments $[OB]$, $[OC]$, $[OA]$.
\item En déduis la nature du triangle $A'B'C'$ et le rapport entre son aire et celle de $ABC$.
\end{enumerate}
\item \textbf{Image du cercle circonscrit.} On note $\mathcal{C}$ le cercle circonscrit au triangle $ABC$, de rayon $R$.
\begin{enumerate}
\item Quelle est l'image de $\mathcal{C}$ par $f$ ? Précise son centre et son rayon.
\item Le cercle image est-il le cercle circonscrit au triangle $A'B'C'$ ? Justifie.
\end{enumerate}
\item \textbf{Écriture analytique.} On munit le plan du repère orthonormé direct $(O\,;\vec{i},\vec{j})$ avec $A(1\,;0)$.
\begin{enumerate}
\item Donne les coordonnées de $B$ et $C$.
\item Établis l'écriture analytique de $f$, puis vérifie par le calcul les coordonnées de $A'$, $B'$ et $C'$ obtenues à la question \textbf{2.b}.
\end{enumerate}
\end{enumerate}
\end{exercice}

\newpage

\coursec{Corrigés des exercices}

\begin{corrige}[Exercice 1 — QCM : une seule bonne réponse]
\textbf{1. Réponse b).} L'écriture $\begin{cases} x' = x + 3 \\ y' = y - 2 \end{cases}$ est de la forme $\begin{cases} x' = x + a \\ y' = y + b \end{cases}$ avec $a = 3$ et $b = -2$ : c'est la définition analytique de la \cle{translation} de vecteur $\vec{u}\begin{pmatrix} 3 \\ -2 \end{pmatrix}$. Ce n'est pas une homothétie (le coefficient de $x$ et de $y$ vaut $1$, pas un rapport $k \neq 1$) ni une rotation.

\textbf{2. Réponse a).} L'homothétie de centre $\Omega(x_\Omega\,;y_\Omega)$ et de rapport $k$ a pour écriture $\begin{cases} x' = kx + (1-k)x_\Omega \\ y' = ky + (1-k)y_\Omega \end{cases}$. Ici $k = -3$, $x_\Omega = 1$, $y_\Omega = -2$, donc $(1-k)x_\Omega = 4 \times 1 = 4$ et $(1-k)y_\Omega = 4 \times (-2) = -8$. D'où $\begin{cases} x' = -3x + 4 \\ y' = -3y - 8 \end{cases}$.

\textbf{3. Réponse a).} La composée de deux rotations de même centre $O$ est la rotation de centre $O$ dont l'angle est la \emph{somme} des angles : $\dfrac{\pi}{3} + \dfrac{\pi}{6} = \dfrac{2\pi}{6} + \dfrac{\pi}{6} = \dfrac{3\pi}{6} = \dfrac{\pi}{2}$.

\textbf{4. Réponse a).} La composée de deux homothéties de même centre $O$ et de rapports $k$ et $k'$ est l'homothétie de centre $O$ et de rapport $k \times k' = 2 \times \dfrac{1}{2} = 1$. L'homothétie de rapport $1$ est l'\cle{identité} du plan.
\end{corrige}

\begin{corrige}[Exercice 2 — Vrai ou faux ? Justifier]
\begin{enumerate}
\item \textbf{VRAI.} Par définition, une transformation du plan est une \cle{bijection} du plan dans lui-même. Or toute bijection $f$ admet une application réciproque $f^{-1}$, qui est également bijective : c'est donc aussi une transformation du plan.
\item \textbf{FAUX.} La composée d'une homothétie de rapport $k$ et d'une translation est une homothétie \emph{uniquement si} $k \neq 1$ (le rapport de la composée est alors $k$). Si $k = 1$, l'homothétie est l'identité et la composée se réduit à la translation, qui n'est pas une homothétie. Contre-exemple : $h = \mathrm{Id}$ et $t$ translation de vecteur $\vec{u} \neq \vec{0}$ donnent $t \circ h = t$, une translation.
\item \textbf{VRAI.} C'est une propriété du cours : la composée, dans un ordre quelconque, d'une homothétie de centre $O$ et de rapport $k > 0$ et d'une rotation de même centre $O$ et d'angle $\theta$ est une \cle{similitude directe} de centre $O$, de rapport $k$ et d'angle $\theta$.
\item \textbf{VRAI.} L'écriture $\begin{cases} x' = 2x + 1 \\ y' = 2y - 3 \end{cases}$ est de la forme $\begin{cases} x' = kx + a \\ y' = ky + b \end{cases}$ avec $k = 2 \neq 1$ : c'est bien une homothétie de rapport $2$. (Son centre $\Omega$ vérifie $x_\Omega = 2x_\Omega + 1$ et $y_\Omega = 2y_\Omega - 3$, soit $\Omega(-1\,;3)$.)
\end{enumerate}
\end{corrige}

\begin{corrige}[Exercice 3 — Reconnaître une transformation]
\begin{enumerate}
\item $\begin{cases} x' = x - 4 \\ y' = y + 1 \end{cases}$ : forme $\begin{cases} x' = x + a \\ y' = y + b \end{cases}$ avec $a = -4$, $b = 1$. C'est la \textbf{translation} de vecteur $\vec{u}\begin{pmatrix} -4 \\ 1 \end{pmatrix}$.
\item $\begin{cases} x' = 5x \\ y' = 5y \end{cases}$ : forme $\begin{cases} x' = kx \\ y' = ky \end{cases}$ avec $k = 5$. C'est l'\textbf{homothétie} de centre $O$ (l'origine du repère) et de rapport $5$.
\item $\begin{cases} x' = -x + 6 \\ y' = -y - 2 \end{cases}$ : forme $\begin{cases} x' = kx + a \\ y' = ky + b \end{cases}$ avec $k = -1 \neq 1$. C'est l'\textbf{homothétie} de rapport $-1$, c'est-à-dire la \textbf{symétrie centrale}. Son centre $\Omega$ est l'unique point fixe : $x_\Omega = -x_\Omega + 6$ donne $x_\Omega = 3$ ; $y_\Omega = -y_\Omega - 2$ donne $y_\Omega = -1$. Donc $\Omega(3\,;-1)$.
\item $\begin{cases} x' = x \\ y' = y \end{cases}$ : tout point est invariant. C'est l'\textbf{identité} du plan (que l'on peut voir comme l'homothétie de rapport $1$ ou la translation de vecteur nul).
\end{enumerate}
\end{corrige}

\begin{corrige}[Exercice 4 — Calculs directs]
\begin{enumerate}
\item La translation $t$ de vecteur $\vec{u}\begin{pmatrix} -2 \\ 5 \end{pmatrix}$ a pour écriture analytique :
\[ \begin{cases} x' = x - 2 \\ y' = y + 5 \end{cases} \]
Image de $A(3\,;-1)$ : $x' = 3 - 2 = 1$ et $y' = -1 + 5 = 4$. Donc $t(A) = A'(1\,;4)$.
\item L'homothétie $h$ de centre $\Omega(2\,;1)$ et de rapport $k = -2$ a pour écriture :
\[ \begin{cases} x' = -2x + (1-(-2)) \times 2 \\ y' = -2y + (1-(-2)) \times 1 \end{cases} \quad \text{soit} \quad \begin{cases} x' = -2x + 6 \\ y' = -2y + 3 \end{cases} \]
Image de $B(0\,;3)$ : $x' = -2 \times 0 + 6 = 6$ et $y' = -2 \times 3 + 3 = -3$. Donc $h(B) = B'(6\,;-3)$.
\item La composée de deux rotations de même centre $O$ est la rotation de centre $O$ dont l'angle est la somme des angles : $\dfrac{\pi}{4} + \dfrac{\pi}{4} = \dfrac{\pi}{2}$. Donc $r_2 \circ r_1$ est la \textbf{rotation de centre $O$ et d'angle $\dfrac{\pi}{2}$} (le quart de tour direct de centre $O$).
\end{enumerate}
\end{corrige}

\begin{corrige}[Exercice 5 — Écriture analytique d'une homothétie]
\begin{enumerate}
\item Soit $M(x\,;y)$ et $M'(x'\,;y') = h(M)$. La relation $\overrightarrow{\Omega M'} = 3\,\overrightarrow{\Omega M}$ se traduit, avec $\Omega(-1\,;2)$, par :
\[ \begin{cases} x' - (-1) = 3\big(x - (-1)\big) \\ y' - 2 = 3(y - 2) \end{cases} \iff \begin{cases} x' + 1 = 3x + 3 \\ y' - 2 = 3y - 6 \end{cases} \iff \begin{cases} x' = 3x + 2 \\ y' = 3y - 4 \end{cases} \]
\item Images :
\begin{itemize}
\item $A(1\,;0)$ : $x' = 3 \times 1 + 2 = 5$, $y' = 3 \times 0 - 4 = -4$. Donc $A'(5\,;-4)$.
\item $B(-1\,;2)$ : $x' = 3 \times (-1) + 2 = -1$, $y' = 3 \times 2 - 4 = 2$. Donc $B'(-1\,;2) = B$.
\item $C(0\,;-3)$ : $x' = 3 \times 0 + 2 = 2$, $y' = 3 \times (-3) - 4 = -13$. Donc $C'(2\,;-13)$.
\end{itemize}
\textbf{Remarque :} le point $B$ est invariant par $h$. C'est normal : $B = \Omega$ est le \cle{centre} de l'homothétie, et le centre est l'unique point fixe d'une homothétie de rapport différent de $1$.
\item On exprime $x$ et $y$ en fonction de $x'$ et $y'$ :
\[ \begin{cases} x' = 3x + 2 \\ y' = 3y - 4 \end{cases} \iff \begin{cases} x = \dfrac{x' - 2}{3} \\ y = \dfrac{y' + 4}{3} \end{cases} \]
L'application réciproque $h^{-1}$ associe donc à $M'(x'\,;y')$ le point $M\left(\dfrac{x'-2}{3}\,;\dfrac{y'+4}{3}\right)$. On reconnaît l'écriture $\begin{cases} x = \dfrac{1}{3}x' + (1-\tfrac{1}{3})(-1) \\ y = \dfrac{1}{3}y' + (1-\tfrac{1}{3})(2) \end{cases}$ : $h^{-1}$ est l'\textbf{homothétie de même centre $\Omega(-1\,;2)$ et de rapport $\dfrac{1}{3}$}. C'est conforme au cours : la réciproque de l'homothétie de centre $\Omega$ et de rapport $k \neq 0$ est l'homothétie de centre $\Omega$ et de rapport $\dfrac{1}{k}$.
\end{enumerate}
\end{corrige}

\begin{corrige}[Exercice 6 — Composée de deux homothéties de même centre]
\begin{enumerate}
\item Avec la formule $\begin{cases} x' = kx + (1-k)x_I \\ y' = ky + (1-k)y_I \end{cases}$ et $I(2\,;-1)$ :
\[ h_1 \ (k_1 = 4) : \begin{cases} x' = 4x + (1-4)\times 2 \\ y' = 4y + (1-4)\times(-1) \end{cases} \text{ soit } \begin{cases} x' = 4x - 6 \\ y' = 4y + 3 \end{cases} \]
\[ h_2 \ (k_2 = -\tfrac{1}{2}) : \begin{cases} x' = -\dfrac{1}{2}x + \dfrac{3}{2}\times 2 \\ y' = -\dfrac{1}{2}y + \dfrac{3}{2}\times(-1) \end{cases} \text{ soit } \begin{cases} x' = -\dfrac{1}{2}x + 3 \\ y' = -\dfrac{1}{2}y - \dfrac{3}{2} \end{cases} \]
\item $f = h_2 \circ h_1$ : on applique d'abord $h_1$, puis $h_2$ au résultat :
\begin{align*}
x' &= -\dfrac{1}{2}(4x - 6) + 3 = -2x + 3 + 3 = -2x + 6 \\
y' &= -\dfrac{1}{2}(4y + 3) - \dfrac{3}{2} = -2y - \dfrac{3}{2} - \dfrac{3}{2} = -2y - 3
\end{align*}
Donc $f : \begin{cases} x' = -2x + 6 \\ y' = -2y - 3 \end{cases}$.
\item $f$ est de la forme $\begin{cases} x' = kx + a \\ y' = ky + b \end{cases}$ avec $k = -2 \neq 1$ : c'est une \textbf{homothétie de rapport $-2$}. Vérification du centre : $-2x + 6 = -2x + (1-(-2))\times 2$ et $-2y - 3 = -2y + 3\times(-1)$, donc le centre est $I(2\,;-1)$. \textbf{Oui, le résultat était prévisible} : le cours affirme que la composée de deux homothéties de même centre $I$ et de rapports $k_1$ et $k_2$ est l'homothétie de centre $I$ et de rapport $k_1 k_2 = 4 \times \left(-\dfrac{1}{2}\right) = -2$ (car $k_1k_2 \neq 1$).
\item $g = h_1 \circ h_2$ :
\begin{align*}
x' &= 4\left(-\dfrac{1}{2}x + 3\right) - 6 = -2x + 12 - 6 = -2x + 6 \\
y' &= 4\left(-\dfrac{1}{2}y - \dfrac{3}{2}\right) + 3 = -2y - 6 + 3 = -2y - 3
\end{align*}
On obtient exactement la même écriture : $g = f$, homothétie de centre $I$ et de rapport $-2$. \textbf{Constatation :} la composée de deux homothéties de \emph{même centre} est \cle{commutative} : $h_1 \circ h_2 = h_2 \circ h_1$.
\end{enumerate}
\end{corrige}

\begin{corrige}[Exercice 7 — Composée d'une homothétie et d'une translation]
\begin{enumerate}
\item $h$ de centre $O$ et de rapport $2$ : $M(x\,;y) \mapsto M_1(2x\,;2y)$. Puis $t$ de vecteur $\vec{u}\begin{pmatrix} 3 \\ 1 \end{pmatrix}$ : $M_1 \mapsto M'(2x+3\,;2y+1)$. Donc :
\[ f = t \circ h : \begin{cases} x' = 2x + 3 \\ y' = 2y + 1 \end{cases} \]
\item $f$ est de la forme $\begin{cases} x' = kx + a \\ y' = ky + b \end{cases}$ avec $k = 2 \neq 1$ : d'après le cours, la composée d'une homothétie de rapport $k \neq 1$ et d'une translation est une \textbf{homothétie de rapport $k = 2$}. Son centre $\Omega$ est l'unique point fixe :
\[ \begin{cases} x_\Omega = 2x_\Omega + 3 \\ y_\Omega = 2y_\Omega + 1 \end{cases} \iff \begin{cases} x_\Omega = -3 \\ y_\Omega = -1 \end{cases} \quad \text{donc } \Omega(-3\,;-1). \]
Vérification : $x' = 2x + (1-2)(-3) = 2x + 3$ \checkmark
\item $g = h \circ t$ : d'abord $t$ : $M(x\,;y) \mapsto M_1(x+3\,;y+1)$, puis $h$ : $M_1 \mapsto M'\big(2(x+3)\,;2(y+1)\big)$. Donc :
\[ g : \begin{cases} x' = 2x + 6 \\ y' = 2y + 2 \end{cases} \]
On n'a \textbf{pas} $f = g$ (par exemple $f(O) = (3\,;1)$ et $g(O) = (6\,;2)$). Centre de $g$ : $x_{\Omega'} = 2x_{\Omega'} + 6$ donne $x_{\Omega'} = -6$ ; $y_{\Omega'} = 2y_{\Omega'} + 2$ donne $y_{\Omega'} = -2$. Donc $\Omega'(-6\,;-2)$. \textbf{Comparaison :} $f$ et $g$ sont deux homothéties de \emph{même rapport} $2$ mais de \emph{centres différents} $\Omega(-3\,;-1)$ et $\Omega'(-6\,;-2)$. La composée d'une homothétie et d'une translation n'est donc \emph{pas commutative} en général.
\end{enumerate}
\end{corrige}

\begin{corrige}[Exercice 8 — Composée homothétie--rotation de même centre]
\begin{enumerate}
\item La rotation $r$ de centre $\Omega(1\,;1)$ et d'angle $\dfrac{\pi}{2}$ (avec $\cos\dfrac{\pi}{2} = 0$ et $\sin\dfrac{\pi}{2} = 1$) :
\[ \begin{cases} x' - 1 = \cos\dfrac{\pi}{2}\,(x-1) - \sin\dfrac{\pi}{2}\,(y-1) \\ y' - 1 = \sin\dfrac{\pi}{2}\,(x-1) + \cos\dfrac{\pi}{2}\,(y-1) \end{cases} \iff \begin{cases} x' = -(y-1) + 1 = -y + 2 \\ y' = (x-1) + 1 = x \end{cases} \]
\item $s = h \circ r$ : on applique d'abord $r$ : $M(x\,;y) \mapsto M_1(-y+2\,;x)$, puis l'homothétie $h$ de centre $\Omega(1\,;1)$ et de rapport $3$ : $M' = \Omega + 3\overrightarrow{\Omega M_1}$ :
\begin{align*}
x' &= 1 + 3\big((-y+2) - 1\big) = 1 + 3(1 - y) = -3y + 4 \\
y' &= 1 + 3(x - 1) = 3x - 2
\end{align*}
Donc $s : \begin{cases} x' = -3y + 4 \\ y' = 3x - 2 \end{cases}$.
\item $s$ est la composée d'une rotation et d'une homothétie de rapport $k = 3 > 0$, de \emph{même centre} $\Omega$ : c'est une \textbf{similitude directe de centre $\Omega(1\,;1)$, de rapport $3$ et d'angle $\dfrac{\pi}{2}$}. (On vérifie la forme $\begin{cases} x' = a x - b y + c \\ y' = b x + a y + c' \end{cases}$ avec $a = 0$, $b = 3$, et $a^2 + b^2 = 9 = 3^2$.)
\item Image de $A(2\,;1)$ : $x' = -3 \times 1 + 4 = 1$, $y' = 3 \times 2 - 2 = 4$. Donc $A'(1\,;4)$.

\textbf{Vérification géométrique :} $\overrightarrow{\Omega A}\begin{pmatrix} 2-1 \\ 1-1 \end{pmatrix} = \begin{pmatrix} 1 \\ 0 \end{pmatrix}$ et $\overrightarrow{\Omega A'}\begin{pmatrix} 1-1 \\ 4-1 \end{pmatrix} = \begin{pmatrix} 0 \\ 3 \end{pmatrix}$.
\begin{itemize}
\item Distance : $\Omega A = 1$ et $\Omega A' = \sqrt{0^2 + 3^2} = 3 = k \times \Omega A$ \checkmark
\item Angle : $\left(\overrightarrow{\Omega A}\,;\overrightarrow{\Omega A'}\right) = \left(\vec{i}\,;3\vec{j}\right) = \dfrac{\pi}{2}$ \checkmark
\end{itemize}
La similitude multiplie bien les distances par $3$ et fait tourner de $\dfrac{\pi}{2}$.
\end{enumerate}
\end{corrige}

\begin{corrige}[Exercice 9 — Type Probatoire : étude complète d'une application du plan]
\begin{enumerate}
\item \textbf{a.} Soit $M'(x'\,;y')$ un point donné. On cherche $M(x\,;y)$ tel que $f(M) = M'$, c'est-à-dire on résout en $(x\,;y)$ :
\[ \begin{cases} 2x - y = x' - 1 \\ x + 2y = y' + 3 \end{cases} \]
Le déterminant de ce système est $\begin{vmatrix} 2 & -1 \\ 1 & 2 \end{vmatrix} = 2 \times 2 - (-1) \times 1 = 5 \neq 0$ : le système admet une \emph{unique} solution. Par les formules de Cramer (ou par combinaison : $2L_1 + L_2$ donne $5x = 2x' + y' + 1$ ; $-L_1 + 2L_2$ donne $5y = -x' + 2y' + 7$) :
\[ x = \dfrac{2x' + y' + 1}{5} \qquad \text{et} \qquad y = \dfrac{-x' + 2y' + 7}{5} \]
Pour tout point $M'$, il existe donc un unique point $M$ tel que $f(M) = M'$.

\textbf{b.} $f$ est donc une bijection du plan dans lui-même : c'est une \textbf{transformation du plan}. Sa réciproque $f^{-1}$ associe à tout point $M'(x'\,;y')$ le point $M$ ci-dessus, autrement dit $f^{-1}$ a pour écriture analytique (en renommant les variables) :
\[ f^{-1} : \begin{cases} x' = \dfrac{2x + y + 1}{5} \\[2mm] y' = \dfrac{-x + 2y + 7}{5} \end{cases} \]

\item \textbf{a.} L'écriture $\begin{cases} x' = 2x - y + 1 \\ y' = x + 2y - 3 \end{cases}$ est bien de la forme $\begin{cases} x' = ax - by + c \\ y' = bx + ay + c' \end{cases}$ avec $a = 2$, $b = 1$, $c = 1$, $c' = -3$. On calcule : $a^2 + b^2 = 4 + 1 = 5$.

\textbf{b.} $k = \sqrt{a^2+b^2} = \sqrt{5}$. Alors $\left(\dfrac{a}{k}\right)^2 + \left(\dfrac{b}{k}\right)^2 = \dfrac{a^2+b^2}{k^2} = \dfrac{5}{5} = 1$ : le point de coordonnées $\left(\dfrac{2}{\sqrt{5}}\,;\dfrac{1}{\sqrt{5}}\right)$ est sur le cercle trigonométrique, donc il existe un réel $\theta$ (unique modulo $2\pi$) tel que $\cos\theta = \dfrac{2}{\sqrt{5}} = \dfrac{2\sqrt{5}}{5}$ et $\sin\theta = \dfrac{1}{\sqrt{5}} = \dfrac{\sqrt{5}}{5}$. Comme $\cos\theta > 0$ et $\sin\theta > 0$, $\theta$ est dans le premier quadrant : $\theta \approx \num{0,464}$ rad (soit environ $26,6\degres$).

\textbf{c.} Le point fixe $\Omega$ vérifie $f(\Omega) = \Omega$ :
\[ \begin{cases} x_\Omega = 2x_\Omega - y_\Omega + 1 \\ y_\Omega = x_\Omega + 2y_\Omega - 3 \end{cases} \iff \begin{cases} x_\Omega - y_\Omega = -1 \\ x_\Omega + y_\Omega = 3 \end{cases} \]
En additionnant les deux équations : $2x_\Omega = 2$, donc $x_\Omega = 1$ ; puis $y_\Omega = 3 - 1 = 2$.
Vérification : $f(1\,;2)$ donne $x' = 2 - 2 + 1 = 1$ et $y' = 1 + 4 - 3 = 2$. Donc $\Omega(1\,;2)$ est l'unique point fixe.

\textbf{d.} \textbf{Conclusion :} $f$ est la \textbf{similitude directe de centre $\Omega(1\,;2)$, de rapport $k = \sqrt{5}$ et d'angle $\theta$ tel que $\cos\theta = \dfrac{2\sqrt{5}}{5}$ et $\sin\theta = \dfrac{\sqrt{5}}{5}$} ($\theta \approx \num{0,464}$ rad). En effet, en retranchant les coordonnées de $\Omega$, l'écriture devient $\begin{cases} x' - 1 = 2(x-1) - (y-2) \\ y' - 2 = (x-1) + 2(y-2) \end{cases}$, forme canonique de la similitude directe de centre $\Omega$.

\item \textbf{a.} $A(1\,;0)$ : $x' = 2 - 0 + 1 = 3$, $y' = 1 + 0 - 3 = -2$, donc $A'(3\,;-2)$. $B(0\,;1)$ : $x' = 0 - 1 + 1 = 0$, $y' = 0 + 2 - 3 = -1$, donc $B'(0\,;-1)$.

\textbf{b.} $AB = \sqrt{(0-1)^2 + (1-0)^2} = \sqrt{2}$. $A'B' = \sqrt{(0-3)^2 + (-1-(-2))^2} = \sqrt{9 + 1} = \sqrt{10}$. Or $k \times AB = \sqrt{5} \times \sqrt{2} = \sqrt{10}$. On a bien $A'B' = k \times AB$ : la similitude multiplie les distances par son rapport $\sqrt{5}$.
\end{enumerate}

\textbf{Barème indicatif (sur 10) :} 1.a : 2 pts (déterminant non nul ou résolution menée) ; 1.b : 1 pt ; 2.a : 1 pt ; 2.b : 1,5 pt ; 2.c : 1,5 pt ; 2.d : 1 pt ; 3.a : 1 pt ; 3.b : 1 pt.

\textbf{Points de vigilance :}
\begin{itemize}
\item Pour prouver qu'une application est une transformation, il faut établir l'\emph{existence ET l'unicité} de l'antécédent : citer le déterminant non nul est la justification attendue.
\item Dans la forme $x' = ax - by + c$, attention aux signes : ici $b = 1$ car le coefficient de $y$ est $-1$.
\item Le point fixe se calcule en posant $x' = x$ et $y' = y$ ; toute erreur de signe se détecte par la vérification finale — exige-la de toi-même.
\item La conclusion doit donner les \emph{trois} éléments caractéristiques : centre, rapport, angle.
\end{itemize}
\end{corrige}

\begin{corrige}[Exercice 10 — Type Probatoire : chantier de construction à Yaoundé]
\begin{enumerate}
\item \textbf{a.} $h$ de centre $O$ et de rapport $-\dfrac{1}{2}$ : $M(x\,;y) \mapsto M_1\left(-\dfrac{x}{2}\,;-\dfrac{y}{2}\right)$. Puis $t$ de vecteur $\vec{u}\begin{pmatrix} 2 \\ -1 \end{pmatrix}$ : $M_1 \mapsto M'\left(-\dfrac{x}{2}+2\,;-\dfrac{y}{2}-1\right)$. Donc :
\[ f = t \circ h : \begin{cases} x' = -\dfrac{1}{2}x + 2 \\[1mm] y' = -\dfrac{1}{2}y - 1 \end{cases} \]

\textbf{b.} $f$ est de la forme $\begin{cases} x' = kx + a \\ y' = ky + b \end{cases}$ avec $k = -\dfrac{1}{2} \neq 1$ : d'après le cours, la composée d'une homothétie de rapport $k \neq 1$ et d'une translation est une \textbf{homothétie de rapport $-\dfrac{1}{2}$}. Son centre $\Omega$ est l'unique point fixe :
\[ x_\Omega = -\dfrac{1}{2}x_\Omega + 2 \iff \dfrac{3}{2}x_\Omega = 2 \iff x_\Omega = \dfrac{4}{3} \ ; \quad y_\Omega = -\dfrac{1}{2}y_\Omega - 1 \iff \dfrac{3}{2}y_\Omega = -1 \iff y_\Omega = -\dfrac{2}{3}. \]
Donc $\Omega\left(\dfrac{4}{3}\,;-\dfrac{2}{3}\right)$.

\textbf{c.} $t(O) = (2\,;-1)$. Or $\Omega\left(\dfrac{4}{3}\,;-\dfrac{2}{3}\right) \neq (2\,;-1)$ : le centre de $f$ n'est \textbf{pas} l'image de $O$ par $t$. En effet, le centre de $f = t \circ h$ est le point fixe de la composée, et non l'image du centre de $h$ ; il faut résoudre l'équation de point fixe et non se fier à l'intuition.

\item \textbf{a.} Images des sommets par $f$ :
\begin{itemize}
\item $A(0\,;0) \mapsto A'(2\,;-1)$ ;
\item $B(4\,;0) \mapsto B'\left(-2+2\,;-1\right) = B'(0\,;-1)$ ;
\item $C(4\,;2) \mapsto C'(0\,;-2)$ ;
\item $D(0\,;2) \mapsto D'(2\,;-2)$.
\end{itemize}
Le quadrilatère $A'B'C'D'$ est un rectangle de longueur \SI{2}{cm} et de largeur \SI{1}{cm}, image de $ABCD$ par l'homothétie de rapport $-\dfrac{1}{2}$ (le rapport négatif explique le « retournement » : $B$ et $D$ échangent leurs positions relatives).

\textbf{b.} Aire de $ABCD$ : $4 \times 2 = 8\ \mathrm{cm}^2$. Aire de $A'B'C'D'$ : $2 \times 1 = 2\ \mathrm{cm}^2$. L'aire a été multipliée par $\dfrac{2}{8} = \dfrac{1}{4}$. Propriété du cours : une homothétie de rapport $k$ multiplie les aires par $k^2$ ; ici $k^2 = \left(-\dfrac{1}{2}\right)^2 = \dfrac{1}{4}$ \checkmark

\textbf{c.} Échelle : \SI{1}{cm} sur le plan représente \SI{10}{m} en réalité, donc $1\ \mathrm{cm}^2$ représente $10^2 = 100\ \mathrm{m}^2$. L'aire réelle de la dalle $A'B'C'D'$ est donc $2 \times 100 = 200\ \mathrm{m}^2$.

\item \textbf{a.} $g = h \circ t$ : d'abord $t$ : $M(x\,;y) \mapsto M_1(x+2\,;y-1)$, puis $h$ : $M_1 \mapsto M'\left(-\dfrac{x+2}{2}\,;-\dfrac{y-1}{2}\right)$. Donc :
\[ g : \begin{cases} x' = -\dfrac{1}{2}x - 1 \\[1mm] y' = -\dfrac{1}{2}y + \dfrac{1}{2} \end{cases} \]
Centre $\Omega'$ : $x_{\Omega'} = -\dfrac{1}{2}x_{\Omega'} - 1 \iff \dfrac{3}{2}x_{\Omega'} = -1 \iff x_{\Omega'} = -\dfrac{2}{3}$ ; $y_{\Omega'} = -\dfrac{1}{2}y_{\Omega'} + \dfrac{1}{2} \iff \dfrac{3}{2}y_{\Omega'} = \dfrac{1}{2} \iff y_{\Omega'} = \dfrac{1}{3}$. Donc $\Omega'\left(-\dfrac{2}{3}\,;\dfrac{1}{3}\right)$.

\textbf{b.} $\Omega\left(\dfrac{4}{3}\,;-\dfrac{2}{3}\right) \neq \Omega'\left(-\dfrac{2}{3}\,;\dfrac{1}{3}\right)$ : $f$ et $g$ n'ont \textbf{pas} le même centre. Distance :
\[ \Omega\Omega' = \sqrt{\left(\dfrac{4}{3}+\dfrac{2}{3}\right)^2 + \left(-\dfrac{2}{3}-\dfrac{1}{3}\right)^2} = \sqrt{2^2 + (-1)^2} = \sqrt{5}\ \mathrm{cm} \approx \num{2,24}\ \mathrm{cm}. \]
\end{enumerate}

\textbf{Barème indicatif (sur 10) :} 1.a : 1 pt ; 1.b : 2 pts ; 1.c : 1 pt ; 2.a : 1,5 pt ; 2.b : 1,5 pt ; 2.c : 1 pt ; 3.a : 1 pt ; 3.b : 1 pt.

\textbf{Points de vigilance :}
\begin{itemize}
\item Dans $t \circ h$, c'est $h$ qui s'applique \emph{en premier} : ne pas inverser l'ordre de composition.
\item Citer la condition $k \neq 1$ pour conclure que la composée est une homothétie.
\item Pour les aires, le facteur est $k^2$ (et non $k$) ; pour l'échelle, $1\ \mathrm{cm}^2 \leftrightarrow 100\ \mathrm{m}^2$ (et non $10\ \mathrm{m}^2$).
\item Le rapport négatif entraîne un renversement de la figure : bien placer les images $B'$ et $D'$.
\end{itemize}
\end{corrige}

\begin{corrige}[Exercice 11 — Type Probatoire : triangle équilatéral et similitude]
\begin{enumerate}
\item \textbf{a.} $f = r \circ h$ est la composée d'une homothétie de centre $O$ et de rapport $k = \dfrac{1}{2} > 0$ et d'une rotation de \emph{même centre} $O$ et d'angle $\dfrac{2\pi}{3}$. D'après le cours, $f$ est une \textbf{similitude directe de centre $O$, de rapport $\dfrac{1}{2}$ et d'angle $\dfrac{2\pi}{3}$}.

\textbf{b.} $f$ n'est \textbf{pas une rotation} car son rapport $\dfrac{1}{2} \neq 1$ (une rotation conserve les distances, alors que $f$ les divise par $2$). $f$ n'est \textbf{pas une homothétie} car son angle $\dfrac{2\pi}{3}$ n'est ni $0$ ni $\pi$ modulo $2\pi$ (une homothétie envoie tout point $M$ sur la droite $(OM)$, ce qui imposerait un angle nul ou plat).

\item \textbf{a.} $ABC$ est équilatéral de centre $O$ (centre du cercle circonscrit), donc $OA = OB = OC = R$ et les angles entre les rayons valent $\dfrac{2\pi}{3}$ : $\left(\overrightarrow{OA}\,;\overrightarrow{OB}\right) = \left(\overrightarrow{OB}\,;\overrightarrow{OC}\right) = \left(\overrightarrow{OC}\,;\overrightarrow{OA}\right) = \dfrac{2\pi}{3}$. La rotation $r$ de centre $O$ et d'angle $\dfrac{2\pi}{3}$ envoie donc :
\[ r(A) = B, \qquad r(B) = C, \qquad r(C) = A. \]

\textbf{b.} $h$ de centre $O$ et de rapport $\dfrac{1}{2}$ envoie tout point $M$ sur le milieu de $[OM]$. Donc :
\[ A' = f(A) = r\big(h(A)\big) = r(\text{milieu de } [OA]) = \text{milieu de } [Or(A)] = \text{milieu de } [OB], \]
car $r$ conserve les milieux (c'est une isométrie). De même, $B' = f(B)$ est le milieu de $[OC]$ et $C' = f(C)$ est le milieu de $[OA]$.

\textbf{c.} L'image d'un triangle équilatéral par une similitude est un triangle équilatéral (les similitudes conservent les angles) : $A'B'C'$ est \textbf{équilatéral}. Le rapport des aires est le carré du rapport de la similitude : $\mathcal{A}(A'B'C') = k^2 \times \mathcal{A}(ABC) = \dfrac{1}{4}\mathcal{A}(ABC)$.

\item \textbf{a.} L'image du cercle $\mathcal{C}$ de centre $O$ et de rayon $R$ par la similitude $f$ de centre $O$ et de rapport $\dfrac{1}{2}$ est le \textbf{cercle $\mathcal{C}'$ de centre $O$ et de rayon $\dfrac{R}{2}$} (l'image d'un cercle par une similitude est un cercle dont le rayon est multiplié par le rapport).

\textbf{b.} $OA' = OB' = OC' = \dfrac{R}{2}$ (par exemple $A'$ milieu de $[OB]$ avec $OB = R$), donc $A'$, $B'$, $C'$ appartiennent à $\mathcal{C}'$ : $\mathcal{C}'$ est bien le \textbf{cercle circonscrit au triangle $A'B'C'$}.

\item \textbf{a.} Avec $A(1\,;0)$ (donc $R = 1$) et les angles de $\dfrac{2\pi}{3}$ :
\[ B\left(\cos\dfrac{2\pi}{3}\,;\sin\dfrac{2\pi}{3}\right) = \left(-\dfrac{1}{2}\,;\dfrac{\sqrt{3}}{2}\right), \qquad C\left(-\dfrac{1}{2}\,;-\dfrac{\sqrt{3}}{2}\right). \]

\textbf{b.} Écriture analytique de $f = r \circ h$. D'abord $h$ : $M(x\,;y) \mapsto M_1\left(\dfrac{x}{2}\,;\dfrac{y}{2}\right)$. Puis $r$ de centre $O$ et d'angle $\dfrac{2\pi}{3}$, avec $\cos\dfrac{2\pi}{3} = -\dfrac{1}{2}$ et $\sin\dfrac{2\pi}{3} = \dfrac{\sqrt{3}}{2}$ :
\begin{align*}
x' &= -\dfrac{1}{2}\cdot\dfrac{x}{2} - \dfrac{\sqrt{3}}{2}\cdot\dfrac{y}{2} = -\dfrac{x}{4} - \dfrac{\sqrt{3}\,y}{4} \\
y' &= \dfrac{\sqrt{3}}{2}\cdot\dfrac{x}{2} - \dfrac{1}{2}\cdot\dfrac{y}{2} = \dfrac{\sqrt{3}\,x}{4} - \dfrac{y}{4}
\end{align*}
Donc $f : \begin{cases} x' = -\dfrac{1}{4}x - \dfrac{\sqrt{3}}{4}y \\[1mm] y' = \dfrac{\sqrt{3}}{4}x - \dfrac{1}{4}y \end{cases}$.

\textbf{Vérifications :}
\begin{itemize}
\item $A(1\,;0)$ : $x' = -\dfrac{1}{4}$, $y' = \dfrac{\sqrt{3}}{4}$. Or le milieu de $[OB]$ est $\left(-\dfrac{1}{4}\,;\dfrac{\sqrt{3}}{4}\right)$ \checkmark
\item $B\left(-\dfrac{1}{2}\,;\dfrac{\sqrt{3}}{2}\right)$ : $x' = \dfrac{1}{8} - \dfrac{\sqrt{3}}{4}\times\dfrac{\sqrt{3}}{2} = \dfrac{1}{8} - \dfrac{3}{8} = -\dfrac{1}{4}$ ; $y' = \dfrac{\sqrt{3}}{4}\times\left(-\dfrac{1}{2}\right) - \dfrac{1}{4}\times\dfrac{\sqrt{3}}{2} = -\dfrac{\sqrt{3}}{8} - \dfrac{\sqrt{3}}{8} = -\dfrac{\sqrt{3}}{4}$. Or le milieu de $[OC]$ est $\left(-\dfrac{1}{4}\,;-\dfrac{\sqrt{3}}{4}\right)$ \checkmark
\item $C\left(-\dfrac{1}{2}\,;-\dfrac{\sqrt{3}}{2}\right)$ : $x' = \dfrac{1}{8} + \dfrac{3}{8} = \dfrac{1}{2}$ ; $y' = -\dfrac{\sqrt{3}}{8} + \dfrac{\sqrt{3}}{8} = 0$. Or le milieu de $[OA]$ est $\left(\dfrac{1}{2}\,;0\right)$ \checkmark
\end{itemize}
Les résultats analytiques confirment les résultats géométriques de la question \textbf{2.b}.
\end{enumerate}

\textbf{Barème indicatif (sur 10) :} 1.a : 1,5 pt ; 1.b : 1 pt ; 2.a : 1,5 pt ; 2.b : 1,5 pt ; 2.c : 1 pt ; 3.a : 1 pt ; 3.b : 0,5 pt ; 4.a : 1 pt ; 4.b : 1 pt.

\textbf{Points de vigilance :}
\begin{itemize}
\item Pour justifier la nature de $f$, citer explicitement l'hypothèse « même centre $O$ » et $k > 0$ : c'est la condition d'application du théorème.
\item Pour nier que $f$ est une rotation ou une homothétie, argumenter avec le rapport ($\neq 1$) et l'angle ($\neq 0, \pi$) : une simple réponse « non » sans justification ne rapporte aucun point.
\item Dans l'écriture analytique, ne pas oublier que $h$ s'applique d'abord : les coordonnées de $M_1$ sont $\left(\dfrac{x}{2}\,;\dfrac{y}{2}\right)$ avant d'appliquer les formules de rotation.
\item Les vérifications finales (recoupement géométrie/calcul) sont très appréciées au Probatoire : prends le temps de les rédiger.
\end{itemize}
\end{corrige}

\newpage

\coursec{Fiche bilan}

\begin{popfigure}
\begin{center}\tcbox[colback=popblue,colframe=popblue,coltext=white,arc=9pt,boxsep=4pt,left=6pt,right=6pt]{\bfseries\small Transformations du plan}\end{center}
\vspace{-2pt}
\begin{tcbraster}[raster columns=3,raster equal height=rows,raster column skip=6pt,raster row skip=6pt,enhanced,blankest]
\begin{tcolorbox}[enhanced,colback=white,colframe=poporange,boxrule=0.9pt,arc=3pt,title={Cadre g\'en\'eral},fonttitle=\bfseries\scriptsize,coltitle=white,colbacktitle=poporange,left=4pt,right=4pt,top=2pt,bottom=2pt,boxsep=1pt,toptitle=1.5pt,bottomtitle=1.5pt,halign=left]
\begin{itemize}[nosep,leftmargin=*,itemsep=1pt,font=\scriptsize]
\item Bijection du plan
\item R\'eciproque $f^{-1}$
\end{itemize}
\end{tcolorbox}
\begin{tcolorbox}[enhanced,colback=white,colframe=popgreen,boxrule=0.9pt,arc=3pt,title={Translation $t_{\vec{u}}$},fonttitle=\bfseries\scriptsize,coltitle=white,colbacktitle=popgreen,left=4pt,right=4pt,top=2pt,bottom=2pt,boxsep=1pt,toptitle=1.5pt,bottomtitle=1.5pt,halign=left]
\begin{itemize}[nosep,leftmargin=*,itemsep=1pt,font=\scriptsize]
\item $\begin{cases}x'=x+a y'=y+b\end{cases}$
\item $\vec{u}=(a\,;\,b)$
\end{itemize}
\end{tcolorbox}
\begin{tcolorbox}[enhanced,colback=white,colframe=poppurple,boxrule=0.9pt,arc=3pt,title={Homoth\'etie $h(\Omega,k)$},fonttitle=\bfseries\scriptsize,coltitle=white,colbacktitle=poppurple,left=4pt,right=4pt,top=2pt,bottom=2pt,boxsep=1pt,toptitle=1.5pt,bottomtitle=1.5pt,halign=left]
\begin{itemize}[nosep,leftmargin=*,itemsep=1pt,font=\scriptsize]
\item $\begin{cases}x'=kx+(1-k)a y'=ky+(1-k)b\end{cases}$
\item $\overrightarrow{\Omega M'}=k\overrightarrow{\Omega M}$
\end{itemize}
\end{tcolorbox}
\begin{tcolorbox}[enhanced,colback=white,colframe=poppink,boxrule=0.9pt,arc=3pt,title={Compos\'ees m\^eme centre},fonttitle=\bfseries\scriptsize,coltitle=white,colbacktitle=poppink,left=4pt,right=4pt,top=2pt,bottom=2pt,boxsep=1pt,toptitle=1.5pt,bottomtitle=1.5pt,halign=left]
\begin{itemize}[nosep,leftmargin=*,itemsep=1pt,font=\scriptsize]
\item $r_\alpha\circ r_\beta=r_{\alpha+\beta}$
\item $h_k\circ h_{k'}=h_{kk'}$
\end{itemize}
\end{tcolorbox}
\begin{tcolorbox}[enhanced,colback=white,colframe=popgold,boxrule=0.9pt,arc=3pt,title={$h\circ t$ et $t\circ h$},fonttitle=\bfseries\scriptsize,coltitle=white,colbacktitle=popgold,left=4pt,right=4pt,top=2pt,bottom=2pt,boxsep=1pt,toptitle=1.5pt,bottomtitle=1.5pt,halign=left]
\begin{itemize}[nosep,leftmargin=*,itemsep=1pt,font=\scriptsize]
\item Homoth\'etie si $k\neq 1$
\item Translation si $k=1$
\end{itemize}
\end{tcolorbox}
\begin{tcolorbox}[enhanced,colback=white,colframe=popblue!70,boxrule=0.9pt,arc=3pt,title={$h\circ r$ m\^eme centre},fonttitle=\bfseries\scriptsize,coltitle=white,colbacktitle=popblue!70,left=4pt,right=4pt,top=2pt,bottom=2pt,boxsep=1pt,toptitle=1.5pt,bottomtitle=1.5pt,halign=left]
\begin{itemize}[nosep,leftmargin=*,itemsep=1pt,font=\scriptsize]
\item Similitude directe
\item Rapport $|k|$, angle $\theta$
\end{itemize}
\end{tcolorbox}
\end{tcbraster}

\legende{Carte mentale de la le\c{c}on : transformations du plan, d\'efinitions analytiques et compos\'ees.}
\end{popfigure}

\begin{bilan}
\textbf{1. D\'efinitions cl\'es}
\begin{itemize}
  \item Une \cle{transformation} du plan $\mathcal{P}$ est une \cle{bijection} de $\mathcal{P}$ dans lui-m\^eme : tout point a un unique ant\'ec\'edent. Sa \cle{r\'eciproque} $f^{-1}$ v\'erifie $f^{-1}\circ f = f\circ f^{-1}=\mathrm{Id}_{\mathcal{P}}$.
  \item \cle{Translation} $t_{\vec{u}}$ : $\overrightarrow{MM'}=\vec{u}$. \cle{Homoth\'etie} $h(\Omega,k)$ : $\overrightarrow{\Omega M'}=k\,\overrightarrow{\Omega M}$.
  \item \cle{Rotation} $r(\Omega,\theta)$ : $\Omega M'=\Omega M$ et $(\overrightarrow{\Omega M}\,;\,\overrightarrow{\Omega M'})=\theta$.
  \item \cle{Similitude directe} : compos\'ee d'une homoth\'etie de rapport $k>0$ et d'une rotation de m\^eme centre ; elle multiplie les distances par $k$ et conserve les angles orient\'es.
\end{itemize}

\textbf{2. Formules \`a conna\^itre par c\oe ur}
\begin{center}
\begin{tabularx}{\linewidth}{|l|X|}
\hline
\rowcolor{popblueL} \textbf{Transformation} & \textbf{D\'efinition analytique (rep\`ere $(O\,;\vec{\imath},\vec{\jmath})$)} \\
\hline
Translation $t_{\vec{u}}$, $\vec{u}=(a\,;\,b)$ & $\begin{cases}x'=x+a\\ y'=y+b\end{cases}$ \\
\hline
Homoth\'etie $h(\Omega,k)$, $\Omega=(a\,;\,b)$ & $\begin{cases}x'=kx+(1-k)a\\ y'=ky+(1-k)b\end{cases}$ \\
\hline
R\'eciproque de $t_{\vec{u}}$ & $t_{-\vec{u}}$ : $\begin{cases}x'=x-a\\ y'=y-b\end{cases}$ \\
\hline
R\'eciproque de $h(\Omega,k)$, $k\neq 0$ & $h\!\left(\Omega,\dfrac{1}{k}\right)$ \\
\hline
\end{tabularx}
\end{center}

\textbf{3. Compos\'ees : r\'esultats fondamentaux}
\begin{itemize}
  \item Deux \cle{rotations de m\^eme centre} : $r(\Omega,\alpha)\circ r(\Omega,\beta)=r(\Omega,\alpha+\beta)$.
  \item Deux \cle{homoth\'eties de m\^eme centre} : $h(\Omega,k)\circ h(\Omega,k')=h(\Omega,kk')$ si $kk'\neq 1$ ; c'est l'\cle{identit\'e} si $kk'=1$.
  \item \cle{Homoth\'etie et translation} : $t_{\vec{u}}\circ h(\Omega,k)$ (et $h\circ t_{\vec{u}}$) est une translation si $k=1$ ; sinon une \cle{homoth\'etie de rapport $k$} dont le centre $\Omega'$ se d\'etermine par les points fixes ou par l'image de $\Omega$.
  \item \cle{Homoth\'etie et rotation de m\^eme centre} ($k>0$) : similitude directe de centre $\Omega$, de rapport $k$ et d'angle $\theta$ ; $h$ et $r$ commutent : $h\circ r=r\circ h$.
\end{itemize}

\textbf{4. M\'ethodes express}
\begin{itemize}
  \item \emph{Reconna\^itre analytiquement} : forme $\begin{cases}x'=x+a\\ y'=y+b\end{cases}$ $\Rightarrow$ translation ; forme $\begin{cases}x'=kx+p\\ y'=ky+q\end{cases}$ avec $k\neq 1$ $\Rightarrow$ homoth\'etie de rapport $k$, centre obtenu en r\'esolvant $x=kx+p$, $y=ky+q$.
  \item \emph{Prouver qu'une application est une transformation} : montrer que tout $M'(x'\,;\,y')$ poss\`ede un unique ant\'ec\'edent $M(x\,;\,y)$ en exprimant $x,y$ en fonction de $x',y'$ ; les formules obtenues d\'efinissent $f^{-1}$.
  \item \emph{Composer} : \'ecrire les expressions analytiques (ou vectorielles) de chaque transformation, substituer, puis identifier la nature de la compos\'ee et ses \'el\'ements caract\'eristiques.
  \item \emph{Points fixes} : le centre d'une homoth\'etie ou d'une rotation est l'unique point invariant (sauf cas d\'eg\'en\'er\'es) : r\'esoudre $f(M)=M$.
\end{itemize}
\end{bilan}

\begin{aretenir}[Checklist avant le Probatoire]
\begin{itemize}
  \item[$\square$] Je sais d\'efinir une transformation du plan et citer sa propri\'et\'e fondamentale (bijection).
  \item[$\square$] Je sais \'ecrire sans h\'esiter la d\'efinition analytique d'une translation de vecteur $\vec{u}=(a\,;\,b)$.
  \item[$\square$] Je sais \'ecrire la d\'efinition analytique d'une homoth\'etie de centre $\Omega=(a\,;\,b)$ et de rapport $k$.
  \item[$\square$] Je sais reconna\^itre une translation ou une homoth\'etie \`a partir d'expressions du type $x'=\ldots$, $y'=\ldots$.
  \item[$\square$] Je sais d\'eterminer le centre d'une homoth\'etie en r\'esolvant le syst\`eme des points fixes.
  \item[$\square$] Je sais montrer qu'une application est bijective et expliciter sa r\'eciproque analytiquement.
  \item[$\square$] Je conna\^is la r\'eciproque de $t_{\vec{u}}$ ($t_{-\vec{u}}$) et celle de $h(\Omega,k)$ ($h(\Omega,\tfrac{1}{k})$).
  \item[$\square$] Je sais que $r(\Omega,\alpha)\circ r(\Omega,\beta)=r(\Omega,\alpha+\beta)$ (angles modulo $2\pi$).
  \item[$\square$] Je sais que $h(\Omega,k)\circ h(\Omega,k')=h(\Omega,kk')$, avec le cas particulier $kk'=1$ (identit\'e).
  \item[$\square$] Je sais d\'eterminer la nature et le rapport de la compos\'ee d'une homoth\'etie et d'une translation, et trouver son centre.
  \item[$\square$] Je sais caract\'eriser la compos\'ee d'une homoth\'etie et d'une rotation de m\^eme centre (similitude directe : centre, rapport, angle).
  \item[$\square$] Je pense \`a utiliser les transformations pour justifier des propri\'et\'es de configurations (alignement, parall\'elisme, conservation des angles et des rapports de distances).
\end{itemize}
\end{aretenir}

\findecours

\end{document}
